% L'existence et la mort — Philosophie Tle Tech — Cours signé OnBuch+
% Programme officiel MINESEC. Compiler avec : tectonic N16-existence-mort.tex
\newcommand{\DOCMATIERE}{Philosophie}
\newcommand{\DOCNIVEAU}{Tle Tech}
\newcommand{\DOCMODULE}{Philosophie – classe de Terminale technique}
\newcommand{\DOCLECON}{N16}
\newcommand{\DOCTITRE}{L'existence et la mort}
% OnBuch+ — préambule « cours signés » ENRICHI (compilé avec tectonic / XeLaTeX)
% Système visuel « Pop » d'OnBuch+ : fond crème (jamais blanc), bordures encre
% épaisses, ombres « dures » décalées, titres Archivo Black en capitales.
% Version MATHÉMATIQUES : graphes (pgfplots/tikz), boîtes pédagogiques
% (définition, propriété, méthode, attention, le savais-tu, exercice résolu,
% exercice, TP, bilan), page de garde et signature OnBuch+ ancrée en bas.
%
% Macros attendues dans main.tex AVANT \input{preamble} :
%   \DOCMATIERE  \DOCNIVEAU  \DOCMODULE  \DOCLECON (numéro)  \DOCTITRE
\documentclass[11pt]{article}

\usepackage{fontspec}
\usepackage[french]{babel}
\usepackage[a4paper,margin=2cm,top=3.4cm,bottom=2.8cm,headheight=1.4cm,footskip=1.3cm]{geometry}
\usepackage{amsmath,amssymb,amsfonts}
\usepackage{mathtools} % \xmapsto, \coloneqq... souvent utilisés par les modèles
\usepackage{cancel} % simplifications barrées (bilans de forces)
% \abs{z} (module, valeur absolue) et \norm{u} (norme), utilisés par les modèles
\providecommand{\abs}{}\let\abs\relax\DeclarePairedDelimiter{\abs}{\lvert}{\rvert}
\providecommand{\norm}{}\let\norm\relax\DeclarePairedDelimiter{\norm}{\lVert}{\rVert}
\usepackage[table]{xcolor} % [table] : fournit aussi \rowcolors (alternance de lignes)
\usepackage{pagecolor}
\usepackage{tikz}
\usetikzlibrary{arrows.meta,positioning,calc,shapes.geometric,shapes.misc,shapes.arrows,decorations.pathmorphing,decorations.pathreplacing,decorations.markings,patterns,shadows,mindmap,trees,fit,backgrounds,matrix,intersections,angles,quotes}
\usepackage{pgfplots}
\usepackage[version=4]{mhchem} % équations nucléaires \ce{^{235}_{92}U}
\usepackage[european,siunitx]{circuitikz} % schémas électriques
\pgfplotsset{compat=1.18}
\usepgfplotslibrary{fillbetween,groupplots} % aires entre courbes (calcul intégral)
\usepackage{enumitem}
\usepackage{fancyhdr}
% [most] charge toutes les bibliothèques tcolorbox (listings, minted, raster,
% documentation...) et gonfle inutilement la mémoire TeX sur un document long
% (~300 boîtes) : on ne charge que ce qui est réellement utilisé.
\usepackage[skins,breakable]{tcolorbox}
\usepackage{graphicx}
\usepackage{colortbl}
\usepackage{booktabs}
\usepackage{tabularx}
\usepackage{multirow}
\usepackage{diagbox} % case d'en-tête diagonale des tableaux à double entrée
% Colonnes extensibles centrée / gauche / droite (C, L, R), souvent utilisées par les modèles
\newcolumntype{C}{>{\centering\arraybackslash}X}
\newcolumntype{L}{>{\raggedright\arraybackslash}X}
\newcolumntype{R}{>{\raggedleft\arraybackslash}X}
\usepackage{array}
\usepackage{multicol}
\usepackage{siunitx}
% parse-numbers=false : accepte aussi \SI{2,5 \times 10^{-3}}{...}
\sisetup{locale=FR,output-decimal-marker={,},group-minimum-digits=4,parse-numbers=false}
\DeclareSIUnit\ohm{\ensuremath{\symup{\Omega}}} % U+2126 absent de la police
\DeclareSIUnit\division{div} % réglages d'oscilloscope (ms/div, V/div)
\DeclareSIUnit\tr{tr} % tours (vitesse de rotation, ex. \SI{1500}{\tr\per\minute})
\DeclareSIUnit\molar{M} % concentration molaire (ex. \SI{3}{\molar}, \SI{1}{\micro\molar})
\DeclareSIUnit\an{an} % année (le modèle confond parfois avec \year, un compteur TeX) — ex. \SI{5}{\kilo\meter\per\an}
\DeclareSIUnit\FCFA{FCFA} % franc CFA (ex. \SI{500}{\FCFA\per\kilo\gram})
\DeclareSIUnit\degreeBrix{\textdegree Brix} % teneur en sucre d'un sirop/jus (réfractométrie)
\DeclareSIUnit\month{mois} % le modèle utilise parfois \month, un compteur TeX, comme unité de durée
\DeclareSIUnit\microliter{\micro\liter} % le modèle écrit parfois \microliter en un seul token
\DeclareSIUnit\million{millions} % dénombrement (ex. \SI{15}{\million\per\mL} spermatozoïdes)
\usepackage{qrcode}
% \degree : commande fréquemment utilisée par les modèles pour un angle ou une
% température en dehors de \SI{}{} (non fournie par les paquets chargés).
\providecommand{\degree}{^{\circ}}
% \qed (fin de démonstration), utilisé par les modèles sans amsthm
\providecommand{\qed}{\hfill$\square$}
% \barre : double barre des valeurs interdites dans un tableau de variations (tkz-tab)
\providecommand{\barre}{\ensuremath{\|}}
% environnement proof (amsthm non chargé) utilisé par les modèles
\ifdefined\proof\else\newenvironment{proof}[1][Démonstration]{\par\noindent\textit{#1.}\ }{\qed\par}\fi
\usepackage{lastpage}
\usepackage{hyperref}
\hypersetup{hidelinks}

% ============================================================
% Palette « Pop » OnBuch+ — mêmes hex que la classe P (plus_ui.dart)
% ============================================================
\definecolor{poporange}{HTML}{F59321}
\definecolor{popdark}{HTML}{E07A0C}
\definecolor{popcream}{HTML}{FFF6E8}
\definecolor{popink}{HTML}{1C1714}
\definecolor{poppurple}{HTML}{7A5AE0}
\definecolor{popgreen}{HTML}{2E9E6A}
\definecolor{poppink}{HTML}{E0457B}
\definecolor{popblue}{HTML}{2F7DE1}
\definecolor{popgold}{HTML}{C98A12}
\definecolor{popmuted}{HTML}{8A7F76}
\definecolor{popline}{HTML}{EADFD2}
\definecolor{popgray}{HTML}{8A7F76}
\colorlet{popgrey}{popgray}
% Alias tolérants (le modèle nomme parfois les couleurs différemment)
\colorlet{popwhite}{white}
\colorlet{popblack}{popink}
\colorlet{popred}{poppink}
\colorlet{popyellow}{popgold}
% Teintes claires (fonds de tableaux/encadrés) — le modèle nomme parfois
% les couleurs avec un suffixe L (ex. popblueL) sans les avoir définies.
\colorlet{poporangeL}{poporange!20}
\colorlet{popdarkL}{popdark!20}
\colorlet{poppurpleL}{poppurple!20}
\colorlet{popgreenL}{popgreen!20}
\colorlet{poppinkL}{poppink!20}
\colorlet{popblueL}{popblue!20}
\colorlet{popgoldL}{popgold!20}
\colorlet{popredL}{popred!20}
\colorlet{popgrayL}{popgray!20}
% Teintes claires pour fonds de tableaux / zones
\colorlet{popblueL}{popblue!10!white}
\colorlet{poporangeL}{poporange!14!white}
\colorlet{popgreenL}{popgreen!12!white}
\colorlet{poppurpleL}{poppurple!12!white}
\colorlet{poppinkL}{poppink!12!white}
\colorlet{popgoldL}{popgold!14!white}

\pagecolor{popcream}

% ============================================================
% Typographie
% ============================================================
\defaultfontfeatures{Ligatures=TeX}
\setmainfont{PlusJakartaSans-Medium.ttf}[Path=../../../fonts/,BoldFont=PlusJakartaSans-Bold.ttf]
% Maths en sans-serif (Fira Math) avec chiffres et lettres droites de Plus
% Jakarta, pour que les formules se fondent dans le texte.
\usepackage[math-style=ISO,bold-style=ISO]{unicode-math}
\setmathfont{FiraMath-Regular.otf}[Path=../../../fonts/]
\setmathfont{PlusJakartaSans-Medium.ttf}[Path=../../../fonts/,range={up/num,up/{Latin,latin},\mathup}]
\usepackage{icomma} % virgule décimale sans espace en mode math
% Caractères mathématiques Unicode absents des polices de texte (Plus Jakarta,
% Archivo Black) : rendus par leur équivalent mathématique, en texte comme en
% formule — sinon ils s'affichent en carré vide (ex. « Arithmétique dans ℤ »).
\usepackage{newunicodechar}
\newunicodechar{ℝ}{\ensuremath{\mathbb{R}}}
\newunicodechar{ℤ}{\ensuremath{\mathbb{Z}}}
\newunicodechar{ℂ}{\ensuremath{\mathbb{C}}}
\newunicodechar{ℕ}{\ensuremath{\mathbb{N}}}
\newunicodechar{ℚ}{\ensuremath{\mathbb{Q}}}
\newunicodechar{𝔻}{\ensuremath{\mathbb{D}}}
\newunicodechar{α}{\ensuremath{\alpha}}
\newunicodechar{β}{\ensuremath{\beta}}
\newunicodechar{γ}{\ensuremath{\gamma}}
\newunicodechar{δ}{\ensuremath{\delta}}
\newunicodechar{ε}{\ensuremath{\varepsilon}}
\newunicodechar{θ}{\ensuremath{\theta}}
\newunicodechar{λ}{\ensuremath{\lambda}}
\newunicodechar{μ}{\ensuremath{\mu}}
\newunicodechar{σ}{\ensuremath{\sigma}}
\newunicodechar{τ}{\ensuremath{\tau}}
\newunicodechar{φ}{\ensuremath{\varphi}}
\newunicodechar{ψ}{\ensuremath{\psi}}
\newunicodechar{ω}{\ensuremath{\omega}}
\newunicodechar{Σ}{\ensuremath{\Sigma}}
% majuscules produites par \MakeUppercase dans les titres (α-aminés -> Α)
\newunicodechar{Α}{\ensuremath{\alpha}}
\newunicodechar{Β}{\ensuremath{\beta}}
\newunicodechar{Γ}{\ensuremath{\Gamma}}
\newunicodechar{Δ}{\ensuremath{\Delta}}
\newunicodechar{Ε}{\ensuremath{\varepsilon}}
\newunicodechar{Θ}{\ensuremath{\Theta}}
\newunicodechar{Λ}{\ensuremath{\Lambda}}
\newunicodechar{Μ}{\ensuremath{\mu}}
\newunicodechar{Τ}{\ensuremath{\tau}}
\newunicodechar{Φ}{\ensuremath{\Phi}}
\newunicodechar{Ψ}{\ensuremath{\Psi}}
\newunicodechar{Ω}{\ensuremath{\Omega}}
\newunicodechar{↦}{\ensuremath{\mapsto}}
\newunicodechar{∈}{\ensuremath{\in}}
\newunicodechar{∉}{\ensuremath{\notin}}
\newunicodechar{∧}{\ensuremath{\wedge}}
\newunicodechar{⇔}{\ensuremath{\Leftrightarrow}}
\newunicodechar{⟺}{\ensuremath{\Longleftrightarrow}}
\newunicodechar{⇒}{\ensuremath{\Rightarrow}}
\newunicodechar{⟹}{\ensuremath{\Longrightarrow}}
\newunicodechar{∘}{\ensuremath{\circ}}
\newunicodechar{∪}{\ensuremath{\cup}}
\newunicodechar{∩}{\ensuremath{\cap}}
\newunicodechar{≡}{\ensuremath{\equiv}}
\newunicodechar{⊂}{\ensuremath{\subset}}
\newunicodechar{⊥}{\ensuremath{\perp}}
\newunicodechar{∃}{\ensuremath{\exists}}
\newunicodechar{∀}{\ensuremath{\forall}}
\newunicodechar{∛}{\ensuremath{\sqrt[3]{\,}}}
\newunicodechar{∜}{\ensuremath{\sqrt[4]{\,}}}
\newunicodechar{⃗}{\ensuremath{{}^{\to}}}
\newunicodechar{ⁿ}{\ifmmode^{n}\else\textsuperscript{\ensuremath{n}}\fi}
\newunicodechar{ˣ}{\ifmmode^{x}\else\textsuperscript{\ensuremath{x}}\fi}
\newunicodechar{ᵇ}{\ifmmode^{b}\else\textsuperscript{\ensuremath{b}}\fi}
\newunicodechar{ʳ}{\ifmmode^{r}\else\textsuperscript{\ensuremath{r}}\fi}
\newunicodechar{ᵘ}{\ifmmode^{u}\else\textsuperscript{\ensuremath{u}}\fi}
\newunicodechar{ʸ}{\ifmmode^{y}\else\textsuperscript{\ensuremath{y}}\fi}
\newunicodechar{ᵃ}{\ifmmode^{a}\else\textsuperscript{\ensuremath{a}}\fi}
\newunicodechar{ᵉ}{\ifmmode^{e}\else\textsuperscript{\ensuremath{e}}\fi}
\newunicodechar{ᵏ}{\ifmmode^{k}\else\textsuperscript{\ensuremath{k}}\fi}
\newunicodechar{ᵖ}{\ifmmode^{p}\else\textsuperscript{\ensuremath{p}}\fi}
\newunicodechar{ᴺ}{\ifmmode^{N}\else\textsuperscript{\ensuremath{N}}\fi}
\newunicodechar{ᶜ}{\ifmmode^{c}\else\textsuperscript{\ensuremath{c}}\fi}
\newunicodechar{ʰ}{\ifmmode^{h}\else\textsuperscript{\ensuremath{h}}\fi}
\newunicodechar{ᵠ}{\ifmmode^{q}\else\textsuperscript{\ensuremath{q}}\fi}
\newunicodechar{ₙ}{\ifmmode_{n}\else\textsubscript{\ensuremath{n}}\fi}
\newunicodechar{ₐ}{\ifmmode_{a}\else\textsubscript{\ensuremath{a}}\fi}
\newunicodechar{ᵢ}{\ifmmode_{i}\else\textsubscript{\ensuremath{i}}\fi}
\newunicodechar{ᵣ}{\ifmmode_{r}\else\textsubscript{\ensuremath{r}}\fi}
\newunicodechar{ₖ}{\ifmmode_{k}\else\textsubscript{\ensuremath{k}}\fi}
\newunicodechar{ᵦ}{\ifmmode_{\beta}\else\textsubscript{\ensuremath{\beta}}\fi}
\newunicodechar{ₓ}{\ifmmode_{x}\else\textsubscript{\ensuremath{x}}\fi}
\newfontfamily\popdisplay{ArchivoBlack-Regular.ttf}[Path=../../../fonts/]
\newfontfamily\poplogo{SpaceGrotesk-Bold.ttf}[Path=../../../fonts/]
\newfontfamily\popsemi{PlusJakartaSans-SemiBold.ttf}[Path=../../../fonts/]
\newfontfamily\popxbold{PlusJakartaSans-ExtraBold.ttf}[Path=../../../fonts/]
\newfontfamily\popxboldls{PlusJakartaSans-ExtraBold.ttf}[Path=../../../fonts/,LetterSpace=3.0]

\setlist[enumerate]{leftmargin=1.7em,itemsep=5pt,topsep=4pt}
\setlist[itemize]{leftmargin=1.7em,itemsep=4pt,topsep=4pt,label={\color{poporange}\textbullet}}
\setlength{\parindent}{0pt}
\setlength{\parskip}{5pt}
\renewcommand{\arraystretch}{1.25}

% pgfplots : style Pop par défaut
\pgfplotsset{
  every axis/.append style={axis line style={popink,line width=1pt},tick style={popink},
    grid style={popline,line width=0.5pt},label style={font=\small},tick label style={font=\footnotesize}},
  popcurve/.style={poporange,line width=1.8pt,no markers,smooth},
}
% Même style disponible dans un \draw TikZ simple (hors pgfplots)
\tikzset{popcurve/.style={poporange,line width=1.8pt}}
\tikzset{popgrid/.style={popline,very thin}}
\colorlet{popcurve}{poporange} % le nom du style est parfois utilisé comme couleur

% ============================================================
% Pill / squiggle
% ============================================================
\newcommand{\poppill}[3]{%
  \tikz[baseline=-0.55ex]\node[draw=popink,line width=1.2pt,rounded corners=2.6pt,%
    fill=#2,inner xsep=6.5pt,inner ysep=3pt]{%
    \popxboldls\fontsize{7}{7}\selectfont\color{#3}\MakeUppercase{#1}};%
}
\newcommand{\popsquiggle}[1]{%
  \tikz[baseline=-0.4ex]{
    \draw[#1,line width=2.6pt,line cap=round]
      plot [smooth] coordinates {(0,0.09) (0.34,0.01) (0.68,0.21) (1.02,0.01) (1.36,0.10)};
  }%
}
% Mot-clé surligné (marqueur orange sous le texte)
\newcommand{\cle}[1]{{\popxbold\color{popdark}#1}}

% ============================================================
% En-tête / pied
% ============================================================
\pagestyle{fancy}
\fancyhf{}
\renewcommand{\headrulewidth}{0pt}
\renewcommand{\footrulewidth}{0pt}
\lhead{\raisebox{-0.15ex}{\poplogo\fontsize{13}{13}\selectfont\color{popink}OnBuch\color{poporange}+}}
\rhead{\poppill{\DOCMATIERE\ \textbullet\ \DOCNIVEAU\ \textbullet\ Leçon \DOCLECON}{white}{popink}}
\lfoot{\popxbold\fontsize{7.6}{7.6}\selectfont\color{popmuted}\MakeUppercase{Cours signé OnBuch+}\ \textbullet\ {\popsemi\DOCTITRE}}
\rfoot{\tikz[baseline=-0.6ex]\node[rounded corners=8.5pt,fill=popink,minimum height=17pt,minimum width=17pt,inner xsep=5pt,inner sep=0]{\popxbold\fontsize{8}{8}\selectfont\color{popcream}\thepage\,/\,\pageref*{LastPage}};}

% ============================================================
% Titres
% ============================================================
\newcommand{\chaptitre}[2]{%
  \begin{tcolorbox}[enhanced,colback=poporange,colframe=popink,boxrule=2.6pt,arc=11pt,
      drop shadow=popink,left=15pt,right=15pt,top=12pt,bottom=11pt,before skip=3pt,after skip=18pt]
    \poppill{#2}{popink}{popcream}\par\vspace{9pt}
    {\raggedright\hyphenpenalty=10000\exhyphenpenalty=10000\popdisplay\fontsize{22}{24}\selectfont\color{popcream}\MakeUppercase{#1}\par}
  \end{tcolorbox}
}
% Section numérotée : pastille orange avec le numéro + titre Archivo
\newcounter{popsec}
\newcommand{\coursec}[1]{%
  \stepcounter{popsec}\setcounter{popsub}{0}%
  \par\vspace{16pt}\noindent
  \begin{minipage}{\linewidth}
  \tikz[baseline=-0.7ex]\node[draw=popink,line width=1.6pt,fill=poporange,rounded corners=4pt,minimum size=22pt,inner sep=0]{\popdisplay\fontsize{12}{12}\selectfont\color{popcream}\thepopsec};%
  \hspace{8pt}{\popdisplay\fontsize{15}{17}\selectfont\color{popink}\MakeUppercase{#1}}\par
  \vspace{4pt}\noindent{\color{poporange}\rule{36pt}{3.6pt}}
  \end{minipage}\par\nopagebreak\vspace{8pt}
}
\newcounter{popsub}
\newcommand{\courssub}[1]{%
  \stepcounter{popsub}%
  \par\vspace{9pt}\noindent{\popxbold\fontsize{11.5}{13}\selectfont\color{popink}{\color{poporange}\thepopsec.\thepopsub}\hspace{6pt}#1}\par\nopagebreak\vspace{4pt}
}

% ============================================================
% Boîtes pédagogiques — même recette : fond blanc, bordure épaisse colorée,
% ombre dure encre, badge plein. Titre optionnel [#1].
% ============================================================
\tcbset{popbase/.style={breakable,enhanced,colback=white,boxrule=2.3pt,arc=9pt,
  shadow={3pt}{-3pt}{0pt}{popink},left=14pt,right=14pt,top=11pt,bottom=11pt,before skip=11pt,after skip=15pt,
  fontupper=\popsemi\fontsize{10.6}{14.5}\selectfont\color{popink},
  fontlower=\popsemi\fontsize{10.6}{14.5}\selectfont\color{popink}}}
% Coche dessinée (le glyphe ✓ n'existe pas dans les polices du document)
\newcommand{\popcheck}[1]{\tikz[baseline=-0.5ex]\draw[#1,line width=1.6pt,line cap=round,line join=round](-0.13,0.0)--(-0.04,-0.09)--(0.13,0.1);}
\providecommand{\coche}{\popcheck{popgreen}}
\providecommand{\resultat}[1]{\par\smallskip\noindent\fcolorbox{popgreen}{popgreenL}{\parbox{\dimexpr\linewidth-2\fboxsep-2\fboxrule\relax}{\textbf{Résultat.} #1}}\par\smallskip}
\newcommand{\popboxhead}[3]{\poppill{#1}{#2}{white}\ifx\relax#3\relax\else\hspace{6pt}{\popxbold\fontsize{10.6}{12}\selectfont\color{popink}#3}\fi\par\vspace{7pt}}

\newtcolorbox{aretenir}[1][]{popbase,colframe=popgreen,before upper={\popboxhead{À retenir}{popgreen}{#1}}}
\newtcolorbox{exemplebox}[1][]{popbase,colframe=poporange,before upper={\popboxhead{Exemple}{poporange}{#1}}}
\newtcolorbox{definition}[1][]{popbase,colframe=popblue,before upper={\popboxhead{Définition}{popblue}{#1}}}
\newtcolorbox{propriete}[1][]{popbase,colframe=poppurple,before upper={\popboxhead{Propriété}{poppurple}{#1}}}
\newtcolorbox{methode}[1][]{popbase,colframe=popgold,before upper={\popboxhead{Méthode}{popgold}{#1}}}
\newtcolorbox{attention}[1][]{popbase,colframe=poppink,before upper={\popboxhead{Attention}{poppink}{#1}}}
\newtcolorbox{experience}[1][]{popbase,colframe=popblue,colback=popblueL,before upper={\popboxhead{Expérience}{popblue}{#1}}}
% « Le savais-tu ? » : carte violette pleine (contexte, culture, Cameroun)
\newtcolorbox{savaistu}[1][]{popbase,colback=poppurple,colframe=popink,
  fontupper=\popsemi\fontsize{10.4}{14.2}\selectfont\color{white},
  before upper={\poppill{Le savais-tu\,?}{white}{poppurple}\ifx\relax#1\relax\else\hspace{6pt}{\popxbold\color{white}#1}\fi\par\vspace{7pt}}}
% Exercice résolu : énoncé en haut, \tcblower puis la solution sur fond crème
\newcounter{popexo}\newcounter{popexr}
\newtcolorbox{exoresolu}[1][]{popbase,colframe=poporange,colbacklower=popcream,
  segmentation style={popink,line width=1.2pt,dashed},
  before upper={\stepcounter{popexr}\popboxhead{Exercice résolu \thepopexr}{poporange}{#1}},
  before lower={\poppill{Solution}{popink}{popcream}\par\vspace{6pt}}}
% Exercice (non résolu en place) — la correction est dans la partie Corrigés
\newtcolorbox{exercice}[1][]{popbase,colframe=popink,
  before upper={\stepcounter{popexo}\popboxhead{Exercice \thepopexo}{popink}{#1}}}
% Corrigé
\newtcolorbox{corrige}[1][]{popbase,colframe=popgreen,colback=popgreenL,
  before upper={\popboxhead{Corrigé}{popgreen}{#1}}}
% Objectifs / prérequis / bilan
\newtcolorbox{objectifs}{popbase,colframe=popink,colback=poporangeL,
  before upper={\popboxhead{Objectifs de la leçon}{popink}{}}}
\newtcolorbox{prerequis}{popbase,colframe=popink,colback=popblueL,
  before upper={\popboxhead{Prérequis}{popblue}{}}}
\newtcolorbox{bilan}{popbase,colframe=popink,colback=white,boxrule=2.8pt,
  before upper={\popboxhead{Fiche bilan}{poporange}{L'essentiel à connaître pour l'examen}}}
% Figure encadrée avec légende
\newtcolorbox{popfigure}[1][]{enhanced,breakable=false,colback=white,colframe=popline,boxrule=1.4pt,arc=8pt,
  left=10pt,right=10pt,top=10pt,bottom=8pt,before skip=10pt,after skip=12pt,halign=center,
  lower separated=false,halign lower=center,
  fontlower=\popsemi\fontsize{9}{11.5}\selectfont\color{popmuted},
  #1}
\newcommand{\legende}[1]{\par\vspace{6pt}{\popsemi\fontsize{9}{11.5}\selectfont\color{popmuted}#1}}

% ============================================================
% Signature OnBuch+
% ============================================================
\newcommand{\poplogoinline}[1]{{\poplogo\fontsize{#1}{#1}\selectfont\color{popink}OnBuch\color{poporange}+}}
\newcommand{\signatureonbuch}{%
  \begin{tcolorbox}[enhanced,colback=popink,colframe=popink,boxrule=2.2pt,arc=10pt,
      drop shadow=poporange,left=14pt,right=14pt,top=10pt,bottom=10pt,before skip=0pt,after skip=0pt]
    \tikz[baseline=-0.7ex]\node[circle,fill=popgreen,draw=popcream,line width=1.3pt,minimum size=22pt,inner sep=0]{\popcheck{white}};%
    \hspace{9pt}\begin{minipage}[c]{0.62\linewidth}
      {\popxbold\fontsize{12}{13}\selectfont\color{popcream}Signé\ \textbullet\ {\poplogo OnBuch\color{poporange}+}}\par\vspace{2pt}
      {\raggedright\popsemi\fontsize{8}{10.5}\selectfont\color{popline}\MakeUppercase{Équipe pédagogique OnBuch+}\\\MakeUppercase{Conforme au programme officiel MINESEC}\par}
    \end{minipage}\hfill
    {\poplogo\fontsize{22}{22}\selectfont\color{popcream}OnBuch\color{poporange}+}
  \end{tcolorbox}%
}

% ============================================================
% Page de garde
% #1 titre · #2 matière · #3 niveau · #4 module · #5 n° leçon · #6 durée ·
% #7 description / objectifs (texte ou itemize)
% ============================================================
\newcommand{\pagegarde}[7]{%
  \thispagestyle{empty}
  \newgeometry{a4paper,margin=1.7cm,top=1.6cm,bottom=1.4cm}%
  \begin{tikzpicture}[remember picture,overlay]
    \fill[poporange!18] ([xshift=-3.2cm,yshift=-3.6cm]current page.north east) circle (4.6cm);
    \fill[poppurple!14] ([xshift=2.4cm,yshift=7.6cm]current page.south west) circle (3.8cm);
  \end{tikzpicture}%
  {\poplogo\fontsize{30}{30}\selectfont\color{popink}OnBuch\color{poporange}+}\hfill
  \raisebox{0.6ex}{\poppill{Cours signé \textbullet\ Premium}{popink}{popcream}}\par
  \vspace{1pt}\popsquiggle{poppurple}\par
  \vspace{8pt}
  \begin{tcolorbox}[enhanced,colback=poporange,colframe=popink,boxrule=2.8pt,arc=13pt,
      drop shadow=popink,left=18pt,right=18pt,top=12pt,bottom=13pt,after skip=10pt]
    \poppill{#2 \textbullet\ #3}{popink}{popcream}\hspace{5pt}\poppill{#4}{white}{popink}\par\vspace{8pt}
    {\popdisplay\fontsize{14}{14}\selectfont\color{popink}LEÇON #5}\par\vspace{4pt}
    {\raggedright\hyphenpenalty=10000\exhyphenpenalty=10000\popdisplay\fontsize{25}{27}\selectfont\color{popcream}\MakeUppercase{#1}\par}
    \vspace{8pt}
    {\popxbold\fontsize{9.5}{9.5}\selectfont\color{popink}\MakeUppercase{Durée conseillée : #6}}
  \end{tcolorbox}
  \begin{tcolorbox}[enhanced,colback=white,colframe=popink,boxrule=2.2pt,arc=10pt,
      drop shadow=popink,left=16pt,right=16pt,top=10pt,bottom=10pt,after skip=8pt]
    \poppill{Dans cette leçon}{poppurple}{white}\par\vspace{5pt}
    \popsemi\fontsize{9.3}{12.6}\selectfont\color{popink}#7
  \end{tcolorbox}
  \noindent\begin{minipage}[t]{0.487\linewidth}
    \begin{tcolorbox}[enhanced,colback=popgreen,colframe=popink,boxrule=2.2pt,arc=10pt,
        drop shadow=popink,left=11pt,right=11pt,top=10pt,bottom=10pt,height=3.1cm,valign=center]
      \tcbox[colback=white,colframe=popink,boxrule=1.3pt,arc=5pt,boxsep=0pt,nobeforeafter,
        left=3pt,right=3pt,top=3pt,bottom=3pt,valign=center]{\qrcode[height=1.9cm]{https://play.google.com/store/apps/details?id=com.luvvix.onbuch&hl=fr}}%
      \hspace{8pt}\begin{minipage}[c]{0.5\linewidth}
        {\popdisplay\fontsize{11}{12.5}\selectfont\color{white}\MakeUppercase{Télécharge OnBuch}}\par\vspace{4pt}
        {\raggedright\popsemi\fontsize{8.4}{11}\selectfont\color{white}Gratuit \textbullet\ Google Play\\Tuteur IA, annales, concours, résultats\par}
      \end{minipage}
    \end{tcolorbox}
  \end{minipage}\hfill
  \begin{minipage}[t]{0.487\linewidth}
    \begin{tcolorbox}[enhanced,colback=poppurple,colframe=popink,boxrule=2.2pt,arc=10pt,
        drop shadow=popink,left=11pt,right=11pt,top=10pt,bottom=10pt,height=3.1cm,valign=center]
      \tikz[baseline=-0.7ex]\node[circle,fill=white,draw=popink,line width=1.3pt,minimum size=1.6cm,inner sep=0]{\popdisplay\fontsize{24}{24}\selectfont\color{poppurple}+};%
      \hspace{8pt}\begin{minipage}[c]{0.6\linewidth}
        {\popdisplay\fontsize{11}{12.5}\selectfont\color{white}\MakeUppercase{Passe à OnBuch+}}\par\vspace{4pt}
        {\raggedright\popsemi\fontsize{8.4}{11}\selectfont\color{white}Tous les cours signés, Tuteur IA illimité, fiches PDF — onglet OnBuch+ de l'app\par}
      \end{minipage}
    \end{tcolorbox}
  \end{minipage}
  \vspace{8pt}
  \popcontenu
  \vfill
  \signatureonbuch
  \restoregeometry
  \newpage
}

% Rangée « Ce cours contient » (page de garde)
\newcommand{\poptile}[3]{% #1 couleur #2 gros texte #3 légende
  \begin{tcolorbox}[enhanced,colback=#1,colframe=popink,boxrule=2pt,arc=9pt,shadow={2.5pt}{-2.5pt}{0pt}{popink},
      left=6pt,right=6pt,top=9pt,bottom=9pt,halign=center,valign=center,height=1.75cm,width=\linewidth,nobeforeafter]
    {\popdisplay\fontsize{12.5}{13}\selectfont\color{white}\MakeUppercase{#2}}\par\vspace{3pt}
    {\centering\popsemi\fontsize{7.8}{9.5}\selectfont\color{white}#3\par}
  \end{tcolorbox}}
\newcommand{\popcontenu}{%
  \par\noindent{\popxboldls\fontsize{7.5}{7.5}\selectfont\color{popmuted}CE COURS SIGNÉ CONTIENT}\par\vspace{7pt}
  \noindent\begin{minipage}[t]{0.235\linewidth}\poptile{popblue}{Cours}{complet, illustré, pas à pas}\end{minipage}\hfill
  \begin{minipage}[t]{0.235\linewidth}\poptile{poporange}{Méthodes}{et exercices résolus}\end{minipage}\hfill
  \begin{minipage}[t]{0.235\linewidth}\poptile{poppink}{Exercices}{type examen + corrigés}\end{minipage}\hfill
  \begin{minipage}[t]{0.235\linewidth}\poptile{popgreen}{Fiche bilan}{l'essentiel à retenir}\end{minipage}\par}

% Sommaire
\newtcolorbox{sommaire}{enhanced,colback=white,colframe=popink,boxrule=2.2pt,arc=10pt,
  drop shadow=popink,left=18pt,right=18pt,top=13pt,bottom=13pt,before skip=6pt,after skip=20pt,
  before upper={\poppill{Sommaire}{poppurple}{white}\par\vspace{9pt}\popsemi\fontsize{10.6}{14.5}\selectfont\color{popink}}}

% Signature de fin de cours — poussée tout en bas de la dernière page
\newcommand{\findecours}{%
  \par\vfill\nopagebreak
  \begin{center}\popsquiggle{poporange}\end{center}\vspace{4pt}
  \signatureonbuch
}
\AtBeginDocument{\def\llbracket{\mathopen{[\mkern-3.5mu[}}\def\rrbracket{\mathclose{]\mkern-3.5mu]}}} % intervalles d'entiers (glyphe ⟦ absent de la police)
% Glyphes absents de Fira Math : redéfinis à partir de glyphes disponibles
\AtBeginDocument{%
  \def\setminus{\mathbin{\backslash}}\let\smallsetminus\setminus
  \def\vdots{\mathord{\vbox{\baselineskip3.2pt\lineskiplimit0pt\kern4pt\hbox{.}\hbox{.}\hbox{.}}}}%
  \def\ddots{\mathinner{\mkern1mu\raise6.5pt\hbox{.}\mkern2mu\raise3.6pt\hbox{.}\mkern2mu\raise0.7pt\hbox{.}\mkern1mu}}%
  \def\triangle{\mathord{\scalebox{0.9}{$\symup{\Delta}$}}}%
  \def\nabla{\mathord{\raisebox{\depth}{\scalebox{1}[-1]{$\symup{\Delta}$}}}}%
  \def\checkmark{\popcheck{popdark}}%
  \def\star{\mathbin{\ast}}\def\bigstar{\mathord{\ast}}%
  \def\diamond{\mathbin{\scalebox{0.6}{$\Diamond$}}}%
  \def\bigcup{\mathop{\mathchoice{\raisebox{-0.3em}{\scalebox{1.7}{$\cup$}}}{\scalebox{1.25}{$\cup$}}{\cup}{\cup}}\displaylimits}%
  \def\bigcap{\mathop{\mathchoice{\raisebox{-0.3em}{\scalebox{1.7}{$\cap$}}}{\scalebox{1.25}{$\cap$}}{\cap}{\cap}}\displaylimits}%
  \def\lesseqqgtr{\mathrel{\lessgtr}}\def\bigoplus{\mathop{\oplus}\displaylimits}%
}
\providecommand{\ssi}{si et seulement si} % abréviation fréquente
\AtBeginDocument{\providecommand{\wideparen}{\overparen}} % arc de cercle

%% FIGFIX-BEGIN v4 — correctifs globaux des figures (docs/onbuch-generation/tools/figfix)
\usepackage{adjustbox}\usepackage{etoolbox}
% toute figure TikZ/pgfplots plus large que la ligne est réduite à la largeur disponible
\BeforeBeginEnvironment{tikzpicture}{\begin{adjustbox}{max width=\linewidth}}
\AfterEndEnvironment{tikzpicture}{\end{adjustbox}}
% légendes pgfplots lisibles par-dessus les courbes
\pgfplotsset{every axis/.append style={legend style={fill=white,fill opacity=0.92,text opacity=1,draw=black!15,inner sep=3pt}}}
% étiquettes d'arêtes (syntaxe edge["..."]) avec fond blanc
\tikzset{every edge quotes/.append style={fill=white,inner sep=1.2pt}}
%% FIGFIX-END
\begin{document}

\pagegarde{L'existence et la mort}{Philosophie}{Tle Tech}{Philosophie – classe de Terminale technique}{N16}{2 séances (fiche de progression harmonisée)}{Cette leçon de philosophie interroge deux notions fondamentales de la condition humaine : l'existence et la mort. À partir de situations concrètes de la vie camerounaise (funérailles, deuils, rites), elle définit ces concepts, examine les grandes conceptions philosophiques du sens de l'existence et outille l'élève pour rédiger une argumentation rigoureuse au Baccalauréat technique.
\vspace{4pt}
\begin{multicols}{2}\small
\begin{itemize}
\item À la fin de cette leçon, je sais définir les concepts d'existence et de mort et les distinguer des notions voisines (vie, être, néant)
\item À la fin de cette leçon, je sais analyser un document (texte philosophique, photographie, donnée chiffrée) portant sur la mort ou le sens de la vie
\item À la fin de cette leçon, je sais expliquer la spécificité de l'existence humaine par rapport à la simple vie biologique
\item À la fin de cette leçon, je sais présenter et comparer au moins trois conceptions philosophiques du sens de l'existence (religieuse, existentialiste, absurde)
\item À la fin de cette leçon, je sais appliquer les notions d'existence et de mort à des situations concrètes de la vie courante au Cameroun
\item À la fin de cette leçon, je sais construire un tableau comparatif et un schéma conceptuel des thèses étudiées
\end{itemize}\end{multicols}}

\begin{sommaire}
\begin{multicols}{2}
\begin{enumerate}
\item Introduction : problématiser l'existence et la mort
\item Définition du concept d'existence
\item Définition du concept de mort
\item Le sens de l'existence humaine : les grandes conceptions
\item Application aux situations concrètes de la vie courante
\item Méthode : rédiger et justifier une démarche argumentative
\item Activité d'intégration
\item Méthodes et exercices résolus
\item Exercices
\item Corrigés des exercices
\item Fiche bilan
\end{enumerate}
\end{multicols}
\end{sommaire}

\begin{objectifs}
\begin{itemize}
\item À la fin de cette leçon, je sais définir les concepts d'existence et de mort et les distinguer des notions voisines (vie, être, néant)
\item À la fin de cette leçon, je sais analyser un document (texte philosophique, photographie, donnée chiffrée) portant sur la mort ou le sens de la vie
\item À la fin de cette leçon, je sais expliquer la spécificité de l'existence humaine par rapport à la simple vie biologique
\item À la fin de cette leçon, je sais présenter et comparer au moins trois conceptions philosophiques du sens de l'existence (religieuse, existentialiste, absurde)
\item À la fin de cette leçon, je sais appliquer les notions d'existence et de mort à des situations concrètes de la vie courante au Cameroun
\item À la fin de cette leçon, je sais construire un tableau comparatif et un schéma conceptuel des thèses étudiées
\item À la fin de cette leçon, je sais rédiger et justifier une démarche argumentative complète (introduction problématisée, thèses argumentées, conclusion personnelle)
\item À la fin de cette leçon, je sais mobiliser des références philosophiques précises (Heidegger, Sartre, Camus, Platon) dans une copie d'examen
\end{itemize}
\end{objectifs}

\begin{prerequis}
\begin{itemize}
\item Notion de sujet et de conscience (leçons antérieures de Terminale)
\item Distinction entre opinion et savoir (introduction à la philosophie)
\item Notions de liberté et de responsabilité (classe de Première)
\item Méthode de l'analyse de texte philosophique
\item Vocabulaire de base : être, néant, vie, conscience
\end{itemize}
\end{prerequis}

\newpage

\coursec{Introduction : problématiser l'existence et la mort}

\courssub{Une situation d'accroche : les funérailles d'un jeune moto-taxi à Yaoundé}

En février 2024, à Yaoundé, les obsèques d'un jeune moto-taxi tué dans un accident à l'entrée de Nsam réunissent des centaines de personnes : pleurs, chants, danses traditionnelles, discours du chef de quartier. Pourquoi les Camerounais consacrent-ils autant de temps et de ressources aux funérailles ? Pourquoi la mort d'un inconnu nous émeut-elle alors que nous savons tous que nous mourrons un jour ? Et surtout : si la mort est certaine, qu'est-ce qui donne du sens à l'existence humaine ? Cette leçon invite à penser philosophiquement ces questions que chacun se pose mais que peu osent examiner rationnellement.

\begin{exemplebox}[Situation d'accroche analysée]
\textbf{Observation :} Un fait social massif (funérailles) mobilise une communauté entière autour d'un événement individuel (la mort d'un jeune homme).

\textbf{Questions qui surgissent :}
\begin{itemize}
    \item Pourquoi ce déploiement d'énergie, de temps, d'argent pour \emph{accompagner} le mort ?
    \item La mort est-elle seulement un fait biologique ou un événement \emph{humain} qui appelle un sens ?
    \item L'existence trouve-t-elle son sens \emph{malgré} la mort, \emph{à cause} de la mort, ou \emph{au-delà} de la mort ?
\end{itemize}
\tcblower
\textbf{Analyse philosophique :} Ce qui frappe, c'est le contraste entre la \textbf{certitude} de la mort (tous mourront) et l'\textbf{incertitude} du sens de la vie. Les funérailles ne sont pas une simple gestion du cadavre : elles sont une \emph{réponse symbolique} à l'angoisse du néant. Elles disent : « cette vie a compté », « ce départ n'est pas une disparition pure ». La philosophie commence là : quand l'émotion laisse place à l'interrogation rationnelle.
\end{exemplebox}

\courssub{Observation guidée de documents : la mort, une réalité socialement située}

\begin{popfigure}
\begin{tikzpicture}
\begin{axis}[
    ybar,
    bar width=18pt,
    width=\linewidth,
    height=9cm,
    ymin=0, ymax=90,
    ylabel={Espérance de vie à la naissance (années)},
    symbolic x coords={Cameroun,Nigeria,France,Japon},
    xtick=data,
    nodes near coords,
    nodes near coords align={vertical},
    every node near coord/.append style={font=\small,popdark},
    ymajorgrids=true,
    grid style={popline},
    axis line style={popdark},
    tick label style={font=\small,popdark},
    label style={font=\small,popdark},
    bar shift=0pt,
]
\addplot[popblue,fill=popblue!70] coordinates {(Cameroun,60) (Nigeria,55) (France,82) (Japon,84)};
\addplot[poporange,fill=poporange!70] coordinates {(Cameroun,62) (Nigeria,57) (France,85) (Japon,87)};
\legend{\textbf{Hommes},\textbf{Femmes}}
\end{axis}
\end{tikzpicture}
\legende{Figure 1 — Espérance de vie à la naissance par pays et par sexe (données INS Cameroun 2023, OMS 2023). L'écart entre le Cameroun (≈ 60 ans) et le Japon (≈ 84 ans) dépasse 20 ans : la mort n'est pas vécue au même moment ni dans les mêmes conditions selon le lieu de naissance.}
\end{popfigure}

\begin{exemplebox}[Analyse du graphique]
\textbf{Données brutes (approximatives, INS/OMS 2023) :}
\begin{center}
\begin{tabularx}{\linewidth}{|l|X|X|}\hline
\textbf{Pays} & \textbf{Hommes} & \textbf{Femmes} \\ \hline
Cameroun & 60 ans & 62 ans \\
Nigeria & 55 ans & 57 ans \\
France & 82 ans & 85 ans \\
Japon & 84 ans & 87 ans \\ \hline
\end{tabularx}
\end{center}

\textbf{Interprétation philosophique :} L'espérance de vie n'est pas une donnée purement biologique ; elle est \textbf{socialement déterminée} (système de santé, nutrition, sécurité, éducation). Mourir à 60 ans au Cameroun et à 84 ans au Japon ne renvoie pas à la même \emph{expérience existentielle} : le temps pour « faire sa vie », transmettre, accomplir des projets diffère radicalement. La mort est \emph{inégalitaire} avant d'être universelle.
\tcblower
\textbf{Question pour la dissertation :} « L'inégalité face à la mort (âge, conditions) contredit-elle l'universalité de la condition mortelle ? » — Cette tension structure le problème du sens de l'existence.
\end{exemplebox}

\begin{exoresolu}[Exercice : lire un indicateur démographique]
\textbf{Énoncé :} À partir du graphique ci-dessus, calculez l'écart d'espérance de vie entre le Cameroun et le Japon pour les hommes. Exprimez cet écart en pourcentage de l'espérance de vie camerounaise. Que signifie cet écart pour une philosophie de l'existence ?

\textbf{Solution détaillée :}
\begin{align*}
\text{Espérance de vie hommes (Cameroun)} &= 60 \text{ ans} \\
\text{Espérance de vie hommes (Japon)} &= 84 \text{ ans} \\
\text{Écart absolu} &= 84 - 60 = 24 \text{ ans} \\
\text{Écart relatif} &= \frac{24}{60} \times 100 = 40\%
\end{align*}
Un Camerounais vit en moyenne \textbf{40\% de temps en moins} qu'un Japonais. Cela signifie que la \emph{durée} disponible pour donner du sens à son existence (étudier, travailler, aimer, transmettre, créer) est structurellement réduite. La philosophie ne peut ignorer cette \textbf{contingence matérielle} : le sens de l'existence ne se pense pas dans l'abstrait, mais \emph{depuis} une situation concrète, historiquement et géographiquement située.
\end{exoresolu}

\begin{savaistu}[Funérailles au Cameroun : célébrer la vie autant que la mort]
Dans plusieurs cultures camerounaises (Bamiléké, Bassa, Beti, Douala, etc.), les funérailles durent \textbf{plusieurs jours}, parfois plusieurs semaines. Elles ne sont pas seulement un « enterrement » mais une \textbf{fête sociale} : danses (tamtam, balafon), chants de deuil et de louange, festins, discours du chef, partage des biens du défunt, rites de succession. On \emph{célèbre} le parcours du mort : ses enfants, ses œuvres, sa place dans la lignée. La mort est pensée comme un \emph{passage} vers le monde des ancêtres, non comme une annihilation. Cette pratique montre que la \textbf{conscience collective de la mort} structure le lien social et donne du sens à l'existence \emph{de son vivant} : « bien vivre, c'est préparer de belles funérailles ».
\end{savaistu}

\courssub{Dégagement des questions spontanées : de l'émotion à la philosophie}

Face à la mort, trois niveaux de questionnement se superposent. Le philosophe les distingue sans les séparer :

\begin{center}
\begin{tabularx}{\linewidth}{|l|X|X|}\hline
\textbf{Niveau} & \textbf{Question type} & \textbf{Exemple issu de la situation} \\ \hline
\textbf{Biologique / Médical} & Comment meurt-on ? Quelles causes ? & « Il a eu un traumatisme crânien suite à l'accident. » \\ \hline
\textbf{Psychologique / Existentiel} & Comment \emph{vais-je} mourir ? Quelle angoisse ? & « J'ai peur de souffrir, de laisser mes enfants seuls. » \\ \hline
\textbf{Philosophique / Ontologique} & \textbf{Qu'est-ce que mourir ? La mort a-t-elle un sens ?} & « La mort est-elle une fin absolue ou un passage ? L'existence prend-elle son sens \emph{parce qu'elle finit} ? » \\ \hline
\end{tabularx}
\end{center}

\begin{attention}[Piège fréquent : confondre la question biologique et la question philosophique]
\emph{Erreur :} Répondre « La mort, c'est l'arrêt des fonctions vitales » à la question « Qu'est-ce que la mort ? ».
\emph{Correction :} Cette définition est \textbf{biologique} (critère médical : mort encéphalique, arrêt cardiaque). La question philosophique porte sur le \textbf{statut ontologique} de la mort : est-elle la \emph{négation} de la vie (néant) ou son \emph{accomplissement} (achèvement) ? Est-elle un \emph{événement} qui m'arrive (extérieur) ou ma \emph{propre possibilité} la plus propre (Heidegger) ? \textbf{Au Bac, ne confondez jamais description scientifique et analyse philosophique.}
\end{attention}

\begin{methode}[Problématiser : de l'observation à la question philosophique]
\textbf{Étape 1 — Partir d'un fait observable (document, expérience, actualité).} \\
Exemple : « Les funérailles mobilisent la communauté ; l'espérance de vie varie de 20 ans selon les pays. »

\textbf{Étape 2 — Dégager une tension, une contradiction, un paradoxe.} \\
Exemple : « La mort est certaine pour tous, mais inégalement vécue ; on la fuit et on la ritualise ; elle détruit l'existence et lui donne sa valeur. »

\textbf{Étape 3 — Formuler une interrogation générale, ouverte, qui appelle une argumentation.} \\
Exemple : « \textbf{La conscience de la mort donne-t-elle ou retire-t-elle du sens à l'existence humaine ?} » (Problématique de la leçon)

\textbf{Étape 4 — Vérifier que la problématique :}
\begin{itemize}
    \item n'est pas une question fermée (oui/non) mais une \emph{invitation à nuancer} ;
    \item met en jeu des \textbf{thèses opposables} (ex. : la mort donne du sens / la mort retire tout sens / le sens est indépendant de la mort) ;
    \item concerne le \textbf{sujet humain} (liberté, responsabilité, finitude) et non seulement l'objet « mort ».
\end{itemize}
\end{methode}

\courssub{Formulation de la problématique de la leçon}

\begin{aretenir}[Problématique centrale]
\cle{La conscience de la mort donne-t-elle ou retire-t-elle du sens à l'existence humaine ?}

Cette question contient trois termes à préciser :
\begin{itemize}
    \item \textbf{Conscience de la mort :} Non pas la peur instinctive (que l'animal peut avoir), mais la \emph{représentation lucide} de sa propre finitude. L'homme est, selon Heidegger, l'être « \emph{être-pour-la-mort} » : il anticipe sa fin comme sa possibilité la plus propre, \emph{imminente} et \emph{inéluctable}.
    \item \textbf{Donner / retirer du sens :} Le sens n'est pas un objet qu'on trouve, mais une \emph{construction} (projet, valeurs, transmission). La mort peut être vue comme \textbf{horizon qui valorise le temps} (chaque instant compte parce qu'il est fini) ou comme \textbf{néant qui annule tout} (tout finit dans l'oubli).
    \item \textbf{Existence humaine :} Non la simple vie biologique (\emph{bios}), mais la \emph{vie vécue, choisie, projetée} (\emph{zoé} vs \emph{bios} chez les Grecs ; \emph{Dasein} chez Heidegger). L'existence est \textbf{liberté} et \textbf{responsabilité}.
\end{itemize}
\end{aretenir}

\begin{popfigure}
\begin{tikzpicture}[
    node distance=1.2cm and 1.5cm,
    every node/.style={font=\small},
    fleche/.style={->,>=Stealth,thick,popdark},
    interrog/.style={draw,rounded corners,fill=popcream,text width=4.5cm,align=center,minimum height=1.8cm,drop shadow={popmuted,opacity=0.3}},
    moment/.style={draw,circle,fill=popblue,text=white,minimum width=1.2cm,font=\bfseries\small},
    moment2/.style={draw,circle,fill=poporange,text=white,minimum width=1.2cm,font=\bfseries\small},
]
% Ligne de vie
\draw[thick,popdark] (0,0) -- (14,0);
% Points de vie
\node[moment] (naissance) at (0,0) {Naissance};
\node[moment] (enfance) at (3,0) {Enfance};
\node[moment] (adulte) at (6,0) {Vie adulte};
\node[moment] (vieux) at (9,0) {Vieillesse};
\node[moment2] (mort) at (12,0) {Mort};
% Flèche temps
\draw[fleche] (12.8,0) -- (14,0) node[right,popdark] {Temps};
% Interrogations au-dessus
\node[interrog, align=center] (q1) at (1.5,2.5){\textbf{Pourquoi suis-je né ?}\\(Sens de l'origine)};
\node[interrog, align=center] (q2) at (4.5,3.2){\textbf{Que faire de ma vie ?}\\(Projet, liberté)};
\node[interrog, align=center] (q3) at (7.5,2.5){\textbf{Ai-je réussi ma vie ?}\\(Bilan, transmission)};
\node[interrog, align=center] (q4) at (10.5,3.2){\textbf{Qu'est-ce qui reste ?}\\(Héritage, mémoire)};
% Interrogation centrale sous la mort
\node[interrog,fill=poppink!20,text width=5cm] (qm) at (12,-2.5) {\textbf{La mort est-elle une fin}\\\textbf{ou un accomplissement ?}\\(Néant / Passage / Achèvement)};
% Liens
\draw[fleche,dashed] (naissance) -- (q1);
\draw[fleche,dashed] (enfance) -- (q2);
\draw[fleche,dashed] (adulte) -- (q3);
\draw[fleche,dashed] (vieux) -- (q4);
\draw[fleche] (mort) -- (qm);
% Annotation
\node[align=center,font=\small\itshape,popmuted] at (6,-4) {Frise existentielle : chaque étape de la vie porte sa question philosophique.\\[2pt]La mort concentre l'interrogation sur le sens global.};
\end{tikzpicture}
\legende{Figure 2 — Frise existentielle : la vie comme parcours questionnant. La mort n'est pas seulement le point final ; elle rétro-éclaire tout le parcours.}
\end{popfigure}

\courssub{Annonce du plan de la leçon}

Cette introduction a posé le décor : un fait social (funérailles), une donnée chiffrée (espérance de vie), des questions spontanées, une problématique philosophique. La suite de la leçon déploiera l'argumentation en trois temps :

\begin{enumerate}
    \item \textbf{Définition rigoureuse des concepts :} Qu'entend-on philosophiquement par \cle{existence} (différence avec la vie biologique, l'essence, le Dasein) et par \cle{mort} (fin biologique, finitude ontologique, mort propre vs mort d'autrui) ?
    \item \textbf{Le sens de l'existence humaine : les grandes conceptions.} Nous confronterons trois familles de réponses :
    \begin{itemize}
        \item \emph{La mort donne du sens :} existentialisme (Heidegger, Sartre), stoïcisme (Sénèque, Marc Aurèle), certaines spiritualités — la finitude fonde l'urgence de vivre authentiquement.
        \item \emph{La mort retire le sens :} nihilisme (Schopenhauer, Cioran), matérialisme radical — le néant final rend tout projet vain.
        \item \emph{Le sens se construit \emph{malgré} la mort :} humanisme (Camus, Ricœur), éthique de la transmission — le sens n'est pas \emph{donné} par la mort mais \emph{créé} par le sujet libre dans le temps qui lui est imparti.
    \end{itemize}
    \item \textbf{Application aux situations concrètes et méthode :} Comment mobiliser ces concepts pour analyser une situation vécue (deuil, choix de fin de vie, engagement, transmission) et rédiger une argumentation structurée (dissertation, commentaire) ?
\end{enumerate}

\begin{exercice}[Pour vous tester avant la suite]
\begin{enumerate}
    \item Relevez dans l'actualité camerounaise récente (derniers 12 mois) un fait funéraire marquant. Identifiez : (a) le fait brut, (b) la question existentielle qu'il soulève, (c) une première ébauche de problématique philosophique.
    \item Expliquez en 3 lignes pourquoi « L'espérance de vie au Cameroun est de 60 ans » est un énoncé \textbf{démographique} et non \textbf{philosophique}. Que doit ajouter le philosophe pour en faire un objet de réflexion ?
    \item Complétez : « La problématique \emph{La conscience de la mort donne-t-elle ou retire-t-elle du sens à l'existence humaine ?} suppose que le sens de l'existence est \_\_\_\_\_\_\_\_\_\_\_\_ ».
\end{enumerate}
\end{exercice}


\coursec{Introduction : problématiser l'existence et la mort}

L'homme est le seul animal qui sait qu'il doit mourir. Cette certitude, loin de n'être qu'une angoisse, est le point de départ de toute philosophie : c'est \emph{parce que} notre temps est compté que nous nous demandons comment l'habiter. Le couple « existence / mort » structure ainsi l'interrogation fondamentale : \textbf{comment donner sens à une existence finie ?}

Dans cette leçon, nous suivrons trois étapes. D'abord, définir rigoureusement l'existence en la distinguant de la simple vie biologique. Ensuite, analyser la mort non comme un simple événement biologique mais comme une structure de l'existence humaine. Enfin, confronter les grandes réponses apportées au sens de la vie (religieuses, humanistes, existentielles) pour outiller l'élève face aux situations concrètes et aux épreuves du Baccalauréat technique.

\coursec{Définition du concept d'existence}

\courssub{L'étymologie : \emph{ex-sistere}, « se tenir hors de »}

Le mot « existence » vient du latin \emph{ex-sistere}, composé du préfixe \emph{ex-} (« hors de ») et du verbe \emph{sistere} (« se tenir », « se placer »). Exister, c'est littéralement \cle{se tenir hors de}, sortir de soi, se manifester au-dehors. L'être qui existe n'est pas enfermé en lui-même : il \emph{sort} dans le monde, il se rend présent aux autres et à lui-même.

\begin{definition}[Existence (sens étymologique et commun)]
Au sens commun, \cle{exister} signifie être là, être présent dans le monde, être réel et non pas seulement imaginé ou possible. Étymologiquement (latin \emph{ex-sistere}), exister c'est « se tenir hors de soi », sortir de soi-même pour se manifester dans le monde et devant les autres.
\end{definition}

Cette étymologie suggère que l'existence n'est pas un état passif (comme une pierre), mais un mouvement, une sortie de soi, une ouverture au monde. L'homme n'est pas un objet parmi les objets ; il est l'être qui \emph{sort de lui-même} — par la conscience, par le projet, par la liberté.

\courssub{Exister n'est pas simplement vivre}

La distinction fondamentale, attendue à l'examen, sépare la vie biologique de l'existence humaine.

Une manguière à Douala \emph{vit} : elle naît, grandit, fructifie, meurt. Une chèvre dans l'Adamaoua \emph{vit} : elle mange, boit, se reproduit. Mais la manguière sait-elle qu'elle vit ? La chèvre se demande-t-elle pourquoi elle est au monde ? Non. La plante et l'animal accomplissent leur vie sans en prendre conscience, sans s'interroger, sans la transformer par décision.

L'homme, lui, fait plus que vivre : il \emph{existe}. Parce qu'il \cle{se sait vivant}. Il a conscience de son être, réfléchit sur sa vie, se demande d'où il vient, où il va, ce qu'il veut en faire. Il peut dire « je ».

\begin{aretenir}[La distinction fondamentale : Vivre vs Exister]
\emph{Vivre} est un fait biologique : naître, se nourrir, se reproduire, mourir. La plante vit, l'animal vit. \emph{Exister} est un fait proprement humain : c'est vivre \textbf{en sachant qu'on vit}, en ayant conscience de soi, en pouvant s'interroger sur le sens de sa vie et la transformer par des choix. Tout homme vit, mais tout être vivant n'existe pas au sens philosophique.
\end{aretenir}

\begin{popfigure}
\begin{center}
\begin{tikzpicture}[scale=1]
% Cercle externe : Exister
\draw[fill=popblueL,draw=popblue,thick] (0,0) circle (3.1);
% Cercle interne : Vivre (concentrique)
\draw[fill=poporangeL,draw=poporange,thick] (0,0) circle (1.7);
% Textes
\node[popblue,font=\bfseries] at (0,2.4) {EXISTER};
\node[popdark,font=\small,align=center] at (0,1.7) {Conscience de soi};
\node[popdark,font=\small,align=center] at (0,1.2) {Projet, temporalité};
\node[popdark,font=\small,align=center] at (0,0.7) {Liberté, choix};
\node[popdark,font=\small,align=center] at (0,0.2) {Rapport à autrui};
\node[poporange,font=\bfseries] at (0,-1.0) {VIVRE};
\node[popmuted,font=\small,align=center] at (0,-1.6) {(Biologique : naître, se nourrir,\\se reproduire, mourir)};
% Flèches explicatives
\draw[->,thick,poporange] (1.9,0) -- (3.2,0);
\node[popmuted,font=\scriptsize,align=left] at (4.5,0) {Socle commun à\\tous les vivants};
\draw[->,thick,popblue] (-1.9,2.0) -- (-3.2,2.0);
\node[popmuted,font=\scriptsize,align=right] at (-4.2,2.0) {Propre à\\l'homme};
\end{tikzpicture}
\end{center}
\legende{Schéma 1 : la vie biologique (cercle central) est le socle commun à la plante, à l'animal et à l'homme ; l'existence (cercle extérieur) ajoute la conscience, le projet, la liberté et le rapport aux autres, qui sont le propre de l'homme.}
\end{popfigure}

\begin{attention}[Erreur fréquente à l'examen]
Ne jamais employer « exister » comme synonyme de « vivre » ou d'« être en vie ». Écrire « l'homme existe parce qu'il respire » est une erreur de niveau : c'est confondre existence et vie biologique. Le correcteur attend la distinction : l'homme \emph{vit} comme tout animal, mais il \emph{existe} parce qu'il a conscience de vivre et construit sa vie par des choix.
\end{attention}

\courssub{L'existence humaine comme existence consciente : le \emph{Dasein} de Heidegger}

L'homme est le seul être pour qui son être est une question. Martin Heidegger (1889-1976), dans \emph{Être et Temps} (1927), forge le concept de \cle{Dasein} (litt. « être-là »).

\begin{definition}[Le Dasein (Heidegger)]
Le \cle{Dasein} désigne l'être humain en tant qu'il est le \textbf{seul être pour qui l'être est une question}. La pierre est là sans se le demander ; l'animal vit sans s'interroger ; l'homme est « là » en se demandant ce que signifie être là. Le Dasein est l'être qui \emph{a à être}, c'est-à-dire qui doit répondre de sa propre existence et décider de ce qu'il en fera.
\end{definition}

\begin{exemplebox}[Expérience ordinaire du Dasein]
Un élève de Terminale technique à Bafoussam, seul un soir, se demande : « Qu'est-ce que je fais de ma vie ? Ce métier de technicien me correspond-il ? » À cet instant, il fait l'expérience du Dasein : il ne subit pas sa vie, il l'interroge. Aucun animal ne fait cela.
\end{exemplebox}

\courssub{Existence et temporalité : l'être projeté vers l'avenir}

La plante vit dans un présent biologique ; l'animal obéit à ses besoins immédiats. L'homme habite le temps tout entier : il se souvient, vit le présent, et surtout \cle{se projette vers l'avenir}. Exister, c'est être « en avant de soi » : l'homme n'est jamais seulement ce qu'il est, il est ce qu'il veut devenir. Sa vie est orientée par des \emph{projets} : réussir le Bac, apprendre un métier, fonder une famille.

\begin{exemplebox}[L'élève technicien et la temporalité]
Un élève de Terminale Froid/Climatisation à Douala révise chaque soir. Son effort présent n'a de sens que par rapport à un avenir anticipé : il se voit diplômé, employé ou à son compte. Il organise son présent en fonction d'un projet d'avenir. Cette structure — passé assumé, présent travaillé, avenir visé — est constitutive de l'existence.
\end{exemplebox}

\courssub{Existence et liberté : « l'existence précède l'essence » (Sartre)}

Jean-Paul Sartre (1905-1980) résume la liberté humaine par : \emph{« l'existence précède l'essence »}.

Un couteau est fabriqué selon un plan : son essence (couper) précède son existence. Pour l'homme, c'est l'inverse : il \emph{existe d'abord}, se trouve « jeté » dans le monde, et se définit \emph{ensuite} par ses choix. Il n'y a pas de « nature humaine » fixée d'avance.

\begin{propriete}[Thèse de Sartre : l'existence précède l'essence]
Selon Sartre (\emph{L'existentialisme est un humanisme}, 1946), l'homme n'a pas d'essence fixée d'avance : il \textbf{existe d'abord}, puis se définit par ses choix. « L'homme n'est rien d'autre que ce qu'il se fait. » Conséquence : l'homme est entièrement responsable de ce qu'il devient. (Thèse à connaître et mobiliser ; démonstration non exigible au Bac.)
\end{propriete}

\begin{attention}[Pièges sur la formule de Sartre]
1. Ne pas inverser : c'est l'\emph{existence} qui précède l'\emph{essence} (l'inverse vaut pour l'objet technique).\\
2. La liberté sartrienne n'est pas « faire n'importe quoi » : elle est inséparable de la \emph{responsabilité}. Choisir, c'est assumer les conséquences devant soi et devant les autres.
\end{attention}

\courssub{Synthèse : les dimensions de l'existence}

L'existence humaine repose sur quatre dimensions solidaires.

\begin{popfigure}
\begin{center}
\begin{tikzpicture}[
  root/.style={rectangle,rounded corners,draw=popblue,fill=popblueL,thick,align=center,font=\bfseries},
  branch/.style={rectangle,rounded corners,draw=poporange,fill=poporangeL,align=center,font=\small},
  link/.style={->,thick,popdark}]
\node[root, align=center] (ex) at (0,0){EXISTENCE\\HUMAINE};
\node[branch, align=center] (c) at (-5.4,-2.6){Conscience de soi\\se savoir être\\(Dasein, Heidegger)};
\node[branch, align=center] (t) at (-1.8,-2.6){Temporalité\\mémoire du passé,\\projet d'avenir};
\node[branch, align=center] (l) at (1.8,-2.6){Liberté\\se définir par\\ses choix (Sartre)};
\node[branch, align=center] (a) at (5.4,-2.6){Rapport à autrui\\exister, c'est se\\manifester devant les autres};
\draw[link] (ex) -- (c);
\draw[link] (ex) -- (t);
\draw[link] (ex) -- (l);
\draw[link] (ex) -- (a);
\end{tikzpicture}
\end{center}
\legende{Schéma 2 : arbre conceptuel de l'existence. Exister, c'est se savoir être, s'inscrire dans le temps, se construire par ses choix et se manifester devant autrui.}
\end{popfigure}

\begin{definition}[L'existence (définition philosophique)]
L'\cle{existence} est le mode d'être propre à l'homme : être conscient qui \textbf{se sait être}, qui \textbf{se projette dans l'avenir} par des projets, et qui \textbf{se définit par ses choix} dans le rapport aux autres. Exister, ce n'est pas seulement vivre : c'est assumer sa vie en conscience et en liberté.
\end{definition}

\courssub{Application : « survivre » ou « exister » ?}

À Douala, de jeunes « vendeurs à la sauvette » courent après les voitures pour gagner de quoi manger le jour même. Leur activité assure la \emph{survie} : besoins biologiques immédiats, sans horizon, sans projet, sans maîtrise de l'avenir. À l'opposé, le jeune qui entre en lycée technique, apprend un métier (électricité, informatique), épargne pour ouvrir son atelier : celui-là \emph{existe} au sens fort, car il transforme sa vie par un projet librement choisi.

\begin{exoresolu}[Analyser une situation concrète]
\textbf{Énoncé.} Deux jeunes de 18 ans à Douala. Le premier vend des sachets d'eau au carrefour, sans perspective autre que manger le soir. Le second suit une formation d'électronicien tout en faisant de menus travaux pour payer ses études ; il veut ouvrir un atelier dans cinq ans. Analysez leur situation.
\tcblower
\textbf{Solution détaillée.} Le premier \emph{survit} : son activité satisfait des besoins immédiats, sans projet. Son temps est un présent répété ; il subit sa vie. Le second \emph{existe} : (1) conscience de sa situation qu'il juge ; (2) organisation du présent en fonction d'un avenir anticipé — l'atelier dans 5 ans (temporalité, projet) ; (3) choix libre de cette voie et acceptation des sacrifices (liberté, responsabilité sartrienne). Différence existentielle, non sociale : le second fait de sa vie une œuvre. Attention : la survie est souvent contrainte, non choisie ; exister pleinement suppose des conditions matérielles minimales.
\end{exoresolu}

\coursec{Définition du concept de mort}

\courssub{La mort comme fait biologique et comme événement humain}

Biologiquement, la mort est la cessation définitive des fonctions vitales (arrêt cardiaque, mort cérébrale). C'est un processus naturel, commun à tous les vivants. Mais pour l'homme, la mort n'est pas qu'un fait : c'est un \emph{événement} qu'il anticipe, redoute, ritualise. L'animal meurt sans le savoir ; l'homme \emph{est mortel}, il sait qu'il mourra.

\begin{definition}[Mort (sens biologique vs sens existentiel)]
Au sens \textbf{biologique}, la mort est l'arrêt irréversible des fonctions d'autorégulation de l'organisme. Au sens \textbf{existentiel}, la mort est la \cle{finitude} qui structure l'existence humaine : la conscience de sa propre finitude oblige l'homme à choisir, à hiérarchiser ses valeurs, à donner du sens à son temps limité.
\end{definition}

\courssub{La mort chez Épicure : « La mort n'est rien pour nous »}

Épicure (341-270 av. J.-C.) propose une thèse rationnelle pour vaincre la peur de la mort.

\begin{propriete}[Thèse d'Épicure : l'innocuité de la mort]
« La mort n'est rien pour nous : tant que nous existons, la mort n'est pas là ; et quand la mort est là, nous n'existons plus. » (\emph{Lettre à Ménécée}). La mort est une privation de sensation ; or le bien et le mal résident dans la sensation. Donc la mort n'est ni un bien ni un mal. (Thèse à connaître ; argumentation non exigible au Bac.)
\end{propriete}

\begin{attention}[Limite de la thèse d'Épicure]
Épicure apaise l'angoisse de \emph{l'après-mort} (le néant), mais ne résout pas l'angoisse de \emph{mourir} (la souffrance, la séparation, l'inachèvement des projets). Heidegger montrera que la mort structure la vie \emph{avant} d'advenir.
\end{attention}

\courssub{La mort chez Heidegger : la « possibilité la plus propre »}

Pour Heidegger, la mort n'est pas un événement futur parmi d'autres (comme un examen), mais la \cle{possibilité la plus propre, la plus certaine et l'inrelatable} du Dasein.

\begin{itemize}
\item \textbf{Propre :} Personne ne peut mourir à ma place.
\item \textbf{Certaine :} Elle arrivera forcément (contrairement à un projet qui peut échouer).
\item \textbf{Inrelatable :} Elle me coupe de tout rapport au monde et aux autres.
\end{itemize}

L'\emph{anticipation} de la mort (la \emph{Sein-zum-Tode}, « être-pour-la-mort ») permet l'\emph{authenticité} : face à ma finitude, je cesse de fuir dans le « On » (les opinions toutes faites, la distraction) et je reprends mes choix à mon compte. La mort donne de l'urgence et de la gravité à l'existence.

\begin{savaistu}[Rites funéraires au Cameroun]
Au Cameroun, les funérailles (\"enterrements\") sont des événements sociaux majeurs (surtout dans les régions de l'Ouest et du Nord-Ouest). Elles durent plusieurs jours, mobilisent la communauté, honorent l'ancêtre et réinsèrent le vivant dans la lignée. Ce n'est pas seulement « enterrer un corps » : c'est gérer symboliquement la finitude, affirmer la continuité du clan et donner sens à la mort par le collectif. Cela illustre que la mort humaine est toujours \emph{culturelle} et \emph{relationnelle}.
\end{savaistu}

\coursec{Le sens de l'existence humaine : les grandes conceptions}

Puisque l'existence est finie et libre, quel sens lui donner ? Trois grandes familles de réponses s'affrontent.

\courssub{Le sens donné : la thèse religieuse (hétérodétermination)}

Pour les grandes religions monothéistes (judaïsme, christianisme, islam) et de nombreuses traditions africaines, le sens de la vie ne s'invente pas, il se \emph{reçoit}. L'homme est créé par Dieu pour une fin : le connaître, l'aimer, le servir, et accéder à une vie éternelle. La mort n'est pas une fin absolue mais un passage.

\begin{exemplebox}[Vision chrétienne / tradition africaine]
« Je suis venu afin que les hommes aient la vie, et qu'ils l'aient en abondance » (Jn 10,10). Dans la tradition Bamiléké ou Bamoun, la vie est un don des ancêtres et de Dieu ; la mort est le retour vers le monde des ancêtres. Le sens : fidélité aux coutumes, transmission de la vie, préparation de l'au-delà.
\end{exemplebox}

\courssub{Le sens construit : l'humanisme existentiel (Sartre, Frankl)}

Si l'existence précède l'essence (Sartre), il n'y a pas de sens préécrit. L'homme \emph{crée} son sens par ses engagements. Viktor Frankl (psychiatre, rescapé des camps), dans \emph{Le Sens de la vie}, montre que même dans l'horreur, l'homme conserve la « dernière des libertés humaines : choisir son attitude ». Le sens naît de l'œuvre, de l'amour, ou de la souffrance assumée.

\begin{propriete}[Thèse de Frankl : la volonté de sens]
L'homme ne cherche pas le plaisir (Freud) ni le pouvoir (Adler), mais du \textbf{sens}. Le sens est unique pour chacun et change d'instant en instant. Il se découvre par : (1) une œuvre/création ; (2) une rencontre/amour ; (3) l'attitude face à une souffrance inévitable.
\end{propriete}

\courssub{L'absurde et la révolte : Camus}

Albert Camus (\emph{Le Mythe de Sisyphe}, 1942) constate un divorce : l'homme réclame du sens, le monde est silencieux. C'est l'\cle{absurde}. Trois issues : le suicide physique (négation de la vie), le suicide philosophique (fuite dans l'illusion religieuse), la \emph{révolte} : vivre \emph{sans appel}, avec lucidité, en multipliant les expériences, en disant « non » à l'injustice et « oui » à la vie.

\begin{aretenir}[Synthèse des trois conceptions]
\begin{center}
\begin{tabularx}{\linewidth}{|l|X|X|X|}\hline
& \textbf{Sens donné (Religieux)} & \textbf{Sens construit (Existentialiste)} & \textbf{Absurde / Révolte (Camus)} \\ \hline
\textbf{Origine du sens} & Dieu / Ancêtres / Cosmos & L'homme lui-même (liberté) & Aucun (le monde est muet) \\ \hline
\textbf{Rôle de la mort} & Passage, accomplissement & Échéance qui donne urgence & Limite absurde, mais révolte \\ \hline
\textbf{Figure clé} & Croyant, Sage traditionnel & Sartre, Frankl, technicien engagé & Sisyphe, l'homme révolté \\ \hline
\end{tabularx}
\end{center}
\end{aretenir}

\coursec{Application aux situations concrètes de la vie courante}

\courssub{Cas 1 : Bioéthique et fin de vie (Euthanasie, soins palliatifs)}
Un patient en phase terminale demande à « partir dignement ».
\begin{itemize}
\item \textbf{Approche existentielle :} Respecter l'autonomie (liberté sartrienne) vs protéger la vie vulnérable (devoir de solidarité). La « mort digne » n'est pas la même pour le croyant (rendre l'âme à Dieu) que pour l'existentialiste (maîtriser sa fin).
\item \textbf{Contexte camerounais :} La loi pénalise l'euthanasie. Les soins palliatifs se développent (ex: Centre Pasteur de Yaoundé, unités de soins à domicile). Le débat oppose souvent modernité médicale et conception communautaire de la vie (la vie n'appartient pas à l'individu seul).
\end{itemize}

\courssub{Cas 2 : Choix d'orientation technique et sens du travail}
Un élève hésite entre une filière « porteuse » (pétrole, informatique) imposée par sa famille et une passion (menuiserie d'art, agroalimentaire).
\begin{itemize}
\item \textbf{Analyse :} Subir un choix (mauvaise foi sartrienne : « je n'ai pas le choix ») vs Assumer un projet (existence authentique). Le technicien qui \emph{choisit} son métier donne sens à son effort ; celui qui \emph{subit} vit sa formation comme une contrainte biologique (survie).
\end{itemize}

\courssub{Cas 3 : Deuil et résilience}
Face à la mort d'un parent, un jeune peut s'effondrer (nihilisme) ou transformer l'épreuve en moteur (hommage, reprise du flambeau).
\begin{itemize}
\item \textbf{Exemple :} Une étudiante en Génie Civil perd son père, maître d'œuvre. Elle achève le chantier inachevé. Elle \emph{existe} par la fidélité à un projet partagé (Frankl : sens par l'œuvre et l'amour).
\end{itemize}

\begin{experience}[Atelier de classe : « Mon projet de vie »]
\textbf{Objectif :} Passer de la survie à l'existence projetée.\\
\textbf{Matériel :} Feuille A4, stylos.\\
\textbf{Protocole :}
1. Lister 3 choses que je \emph{subis} (contraintes : pauvreté, notes moyennes, pression familiale).\\
2. Lister 3 choses que je \emph{choisis} (efforts, lecture, entraide, prière/meditation).\\
3. Formuler un \textbf{projet professionnel} à 5 ans (métier, lieu, niveau de vie).\\
4. Identifier \textbf{un obstacle majeur} et \textbf{une action concrète cette semaine} pour l'affaiblir.\\
\textbf{Interprétation :} L'élève qui remplit la colonne « choix » et définit un projet avec une action concrète fait l'expérience de l'existence (liberté + temporalité). Celui qui ne voit que des contraintes vit en mode « survie ».
\end{experience}

\coursec{Méthode : rédiger et justifier une démarche argumentative}

Au Probatoire / Baccalauréat technique, les sujets portent souvent sur : « L'existence précède-t-elle l'essence ? », « La mort donne-t-elle un sens à la vie ? », « Exister, est-ce se construire ? ».

\begin{methode}[Dissertation philosophique : 4 étapes]
\textbf{1. Analyser le sujet (15 min).}
\begin{enumerate}
\item Définir \textbf{chaque terme} du sujet (ex: « existence », « essence », « précéder »).
\item Identifier la \textbf{tension} : opposition (ex: liberté vs déterminisme) ou articulation.
\item Formuler la \textbf{problématique} : une question ouverte qui guide le devoir. Ex: « Si l'homme se définit par ses choix, est-il pour autant totalement libre de devenir n'importe qui ? »
\end{enumerate}

\textbf{2. Construire le plan dialectique (10 min).}
\begin{itemize}
\item \textbf{Thèse} : L'argument principal (ex: Sartre, l'homme se fait lui-même).
\item \textbf{Antithèse} : Les limites/objections (ex: conditionnement social, biologique, inconscient ; Heidegger : nous sommes « jetés » dans un monde déjà là).
\item \textbf{Synthèse} : Dépassement (ex: La liberté s'exerce \emph{dans} une situation ; le technicien transforme les contraintes matériaux en ouvrage).
\end{itemize}

\textbf{3. Rédiger l'introduction et la conclusion (15 min).}
\begin{itemize}
\item \textbf{Intro :} Accroche (citation, fait concret) $\rightarrow$ Définitions $\rightarrow$ Problématique $\rightarrow$ Annonce du plan.
\item \textbf{Conclusion :} Réponse synthétique à la problématique $\rightarrow$ Ouverture (autre auteur, actualité, autre discipline).
\end{itemize}

\textbf{4. Développer les parties (1h10).}
\begin{itemize}
\item Une idée directrice par paragraphe.
\item Structure : \textbf{Argument} (thèse d'auteur) $\rightarrow$ \textbf{Explication} $\rightarrow$ \textbf{Exemple concret} (Cameroun, technique, vie courante) $\rightarrow$ \textbf{Transition}.
\item Utiliser les connecteurs logiques : \emph{En effet, toutefois, par conséquent, de surcroît, en définitive}.
\end{itemize}
\end{methode}

\begin{methode}[Commentaire de texte : 3 étapes]
\textbf{1. Présentation (5 min) :} Auteur, œuvre, date, contexte, thèse générale.
\textbf{2. Analyse structurée (1h30) :} Suivre le fil du texte paragraphes par paragraphes. Pour chaque étape : \emph{Citer} $\rightarrow$ \emph{Expliquer} (vocabulaire, arguments) $\rightarrow$ \emph{Évaluer} (portée, limites).
\textbf{3. Synthèse critique (15 min) :} Reprendre la thèse, confronter à d'autres auteurs vus en cours (Sartre vs Heidegger vs Épicure), donner son avis argumenté.
\end{methode}

\begin{exercice}[Entraînement Bac - Sujets types]
1. \textbf{Dissertation :} « Exister, est-ce se construire ? »\\
2. \textbf{Dissertation :} « La conscience de la mort est-elle le principe de toute sagesse ? »\\
3. \textbf{Commentaire :} Extrait de \emph{L'existentialisme est un humanisme} (Sartre) : « L'homme n'est rien d'autre que son projet... »\\
4. \textbf{Analyse de situation :} Un technicien supérieur, après 10 ans d'entreprise, quitte tout pour devenir apiculteur dans son village natal. Son entourage le traite de « fou ». En vous appuyant sur les notions d'existence, de liberté et de projet, analysez ce choix de vie.
\end{exercice}

\begin{corrige}[Exercice 1 - Dissertation : « Exister, est-ce se construire ? »]
\textbf{Analyse :} « Exister » (mode d'être humain : conscience, projet) ; « Se construire » (devenir l'auteur de soi par des choix). Tension : liberté radicale (Sartre) vs déterminismes (social, biologique, inconscient).\\
\textbf{Problématique :} L'homme est-il le seul artisan de son existence ou n'en est-il que le produit ?\\
\textbf{Plan :}\\
I. \textbf{Thèse : L'existence comme construction de soi (Sartre).} L'existence précède l'essence. L'homme est un projet. Ex: l'élève qui choisit sa filière technique contre l'avis familial.\\
II. \textbf{Antithèse : Les limites de l'auto-construction (Heidegger, Marx, Freud).} Nous sommes « jetés » dans un monde, une langue, une classe sociale, un corps. Le technicien ne construit pas l'acier, il le travaille. L'inconscient guide des choix qu'on croit libres.\\
III. \textbf{Synthèse : Une liberté située (Merleau-Ponty, Beauvoir).} Se construire, ce n'est pas partir de zéro, c'est \emph{reprendre} sa facticité (son corps, son histoire, son milieu) pour la dépasser. Le technicien qui innove avec des matériaux locaux « se construit » en transformant sa situation.\\
\textbf{Conclusion :} Exister, c'est assumer la tension entre le donné et le choisi. La mort borne ce chantier : urgence de bâtir.
\end{corrige}

\begin{corrige}[Exercice 4 - Situation : Le technicien apiculteur]
\textbf{Analyse :} Ce technicien \emph{existe} au sens fort. (1) \textbf{Conscience :} Il juge sa vie d'entreprise aliénante (perte de sens). (2) \textbf{Projet/Temporalité :} Il anticipe un avenir aligné sur ses valeurs (nature, autonomie, transmission). (3) \textbf{Liberté/Responsabilité :} Il assume le risque (baisse de revenus, jugement social = « fou ») : c'est l'authenticité heideggérienne (sortir du « On »). (4) \textbf{Rapport à autrui :} Il ne fuit pas le monde ; il change de communauté (village vs entreprise).\\
\textbf{Nuance :} Ce choix suppose un capital (épargne, compétences techniques transférables : gestion, observation, rigueur). La liberté s'exerce \emph{dans} une situation. Ce n'est pas une fuite, mais une \emph{reconversion existentielle}.
\end{corrige}

\coursec{Définition du concept de mort}

Nous venons de définir l'existence comme le fait d'être là, de se déployer dans le temps, de se rapporter à soi-même et au monde. Mais toute définition de l'existence appelle immédiatement sa limite : ce qui existe peut cesser d'exister. Si l'existence humaine est un \emph{être-là}, la mort est ce qui en marque le terme. On ne peut donc pas penser l'existence sans penser la mort, pas plus qu'on ne peut comprendre un trajet Douala--Yaoundé sans savoir qu'il a une fin. C'est pourquoi, après avoir défini l'existence, nous devons maintenant définir rigoureusement la \cle{mort}, en distinguant ses dimensions biologique, philosophique et culturelle.

\courssub{Étymologie et premières approches du mot « mort »}

Le mot « mort » vient du latin \emph{mors}, \emph{mortis}, lui-même issu d'une racine indo-européenne \emph{mer-}, qui signifie « disparaître », « s'éteindre ». Dès l'étymologie, la mort est donc pensée comme une \emph{extinction}, un passage de la présence à l'absence. Le verbe latin \emph{mori} (mourir) a donné en français « mort », « mourir », « mortel », « immortalité ».

Dans les langues camerounaises, le vocabulaire de la mort est remarquablement riche et varié, et il n'est jamais neutre. En langue ewondo, on dit \emph{ku} ; en duala, \emph{diba} ; dans les langues du Grassfield, on emploie souvent des périphrases : « il est retourné chez les ancêtres », « il s'est assis dans la brousse », « il a rejoint les pères ». Ces expressions montrent que, dans beaucoup de cultures camerounaises, on évite de nommer directement la mort : on la désigne comme un \emph{voyage} ou un \emph{retour}. Ce fait linguistique est déjà philosophique : il révèle que la mort n'est pas pensée partout comme une simple disparition, mais parfois comme un changement de statut de la personne, qui passe du monde des vivants au monde des ancêtres.

\begin{definition}[La mort]
La \cle{mort} est la cessation définitive de la vie d'un être vivant. Biologiquement, elle est l'arrêt irréversible des fonctions vitales. Philosophiquement, elle est la \emph{limite absolue de l'existence humaine} : elle met fin à toute expérience possible pour le sujet. Selon les conceptions, elle est vécue comme une néantisation (retour au néant) ou comme un passage (vers une autre forme d'être). Elle est la source de l'angoisse humaine, mais aussi, pour certaines philosophies, de la sagesse.
\end{definition}

\courssub{La définition biologique de la mort}

Commençons par l'approche la plus immédiate : celle du médecin et du biologiste. Pour la biologie, un être vivant se définit par un ensemble de fonctions : la respiration, la circulation sanguine, l'activité cérébrale, le métabolisme. La mort est alors définie comme la \emph{cessation irréversible de ces fonctions vitales}.

\begin{propriete}[Critères biologiques de la mort]
On distingue classiquement :
\begin{itemize}
\item la \cle{mort cardiaque} : arrêt définitif du cœur et de la circulation sanguine ;
\item la \cle{mort cérébrale} : destruction irréversible de l'ensemble du cerveau, constatée par l'absence d'activité électrique et de réflexes, même si le cœur bat encore artificiellement ;
\item la \cle{mort clinique} : arrêt apparent des fonctions (cœur, respiration), parfois réversible si la réanimation intervient à temps — d'où l'importance du critère d'\emph{irréversibilité}.
\end{itemize}
\end{propriete}

L'observation quotidienne confirme cette définition : lorsqu'un vieillard meurt à l'hôpital de district, le médecin constate l'absence de pouls, l'absence de souffle, la dilatation fixe des pupilles. Mais la médecine moderne a compliqué le tableau : un patient en mort cérébrale peut être maintenu « en vie » par des machines. Cela montre que la frontière biologique de la mort n'est pas toujours simple à tracer — et c'est précisément là que la philosophie prend le relais de la biologie.

\courssub{La définition philosophique de la mort : néantisation ou passage ?}

La biologie décrit \emph{comment} meurt un organisme ; la philosophie interroge le \emph{sens} de la mort pour l'existence. Deux grandes conceptions s'affrontent.

\textbf{Première conception : la mort comme néantisation.} Pour les philosophes matérialistes (Épicure, Lucrèce, puis les existentialistes comme Sartre), la mort est l'anéantissement pur et simple de l'individu : le corps se décompose, la conscience s'éteint, il n'y a « rien » après. La mort est le \emph{néant}. Cette thèse est cohérente avec la définition biologique : quand les fonctions vitales cessent, le sujet cesse d'être.

\textbf{Seconde conception : la mort comme passage.} Pour Platon, l'âme est immortelle : la mort du corps libère l'âme, qui rejoint le monde des Idées. Le christianisme enseigne la résurrection et la vie éternelle ; l'islam enseigne le jugement et l'au-delà ; beaucoup de traditions africaines enseignent que le défunt devient \emph{ancêtre}, c'est-à-dire une puissance invisible qui continue d'agir sur les vivants. Dans ces perspectives, la mort n'est pas une fin mais un \emph{changement d'état}.

\begin{attention}[Erreur fréquente]
Ne réduis jamais la mort à un phénomène médical ou biologique dans une dissertation. Le médecin constate la mort ; le philosophe interroge son \emph{sens} pour l'existence : que signifie pour un homme de savoir qu'il va mourir ? Faut-il craindre la mort ? La mort donne-t-elle un sens à la vie ou le lui retire-t-elle ? Une copie qui se contente de définir la mort comme « l'arrêt du cœur » reste au niveau du constat et manque le problème philosophique.
\end{attention}

\courssub{La mort comme certitude : le syllogisme de la mortalité}

Une chose est sûre, et aucun homme n'en doute : \emph{tous les hommes meurent}. La tradition philosophique exprime cette certitude par un raisonnement célèbre, le \cle{syllogisme}, hérité d'Aristote et illustré par l'exemple de Socrate :

\begin{exemplebox}[Le syllogisme de la mortalité]
\begin{itemize}
\item \textbf{Prémisse majeure} : tous les hommes sont mortels ;
\item \textbf{Prémisse mineure} : or Socrate est un homme ;
\item \textbf{Conclusion} : donc Socrate est mortel.
\end{itemize}
\end{exemplebox}

Ce raisonnement est \emph{déductif} : si les deux prémisses sont vraies, la conclusion est nécessairement vraie. La majeure énonce une vérité d'expérience universelle : jamais on n'a vu un homme échapper à la mort. Les rois, les riches, les savants, les puissants meurent comme les pauvres. La mort apparaît ainsi comme la seule certitude absolue de la condition humaine : on peut douter de tout, sauf du fait que nous mourrons.

\begin{popfigure}
\begin{center}
\begin{tikzpicture}[font=\small]
\node[draw=popblue, fill=popblueL, rounded corners, text width=4.6cm, align=center] (maj) at (0,3) {\textbf{Majeure}\\ Tous les hommes sont mortels};
\node[draw=poporange, fill=poporangeL, rounded corners, text width=4.6cm, align=center] (min) at (0,1.5) {\textbf{Mineure}\\ Or Socrate est un homme};
\node[draw=popgreen, fill=popgreenL, rounded corners, text width=4.6cm, align=center] (concl) at (0,0) {\textbf{Conclusion}\\ Donc Socrate est mortel};
\draw[-{Stealth}, thick, popdark] (maj.south) -- (min.north);
\draw[-{Stealth}, thick, popdark] (min.south) -- (concl.north);
\node[text width=4.2cm, align=center, popmuted] at (6.8,1.5) {Raisonnement déductif :\\ la vérité des prémisses\\ garantit la vérité\\ de la conclusion};
\draw[-{Stealth}, popmuted, dashed] (4.6,1.5) -- (2.6,1.5);
\end{tikzpicture}
\end{center}
\legende{Schéma du syllogisme de la mortalité (Aristote). De la vérité générale « tous les hommes sont mortels » et du constat « Socrate est un homme », on déduit nécessairement que Socrate est mortel.}
\end{popfigure}

\courssub{La mort, un événement impossible à vivre : Épicure et Wittgenstein}

Voici maintenant un paradoxe troublant : la mort est certaine, mais personne ne peut jamais \emph{faire l'expérience} de sa propre mort. Je peux voir mourir les autres ; je ne me verrai jamais mort, car pour voir, il faut être vivant.

\begin{propriete}[Thèse d'Épicure : la mort n'est rien pour nous]
Épicure (341--270 av. J.-C.) soutient dans sa \emph{Lettre à Ménécée} : « Quand nous sommes, la mort n'est pas ; et quand la mort est, nous ne sommes plus. » Autrement dit, la mort ne coïncide jamais avec notre existence : tant que je vis, elle n'est pas là ; quand elle est là, je ne suis plus là pour la subir. Il est donc \emph{irrationnel de craindre la mort}, car on ne peut craindre que ce qui peut nous affecter, et la mort ne nous affecte jamais. (Cette thèse est exigible au Bac : sache la restituer et la discuter.)
\end{propriete}

Le philosophe Wittgenstein (1889--1951) prolonge cette intuition dans le \emph{Tractatus} : « La mort n'est pas un événement de la vie. On ne vit pas la mort. » Tous les événements de ma vie (naissance, mariage, réussite au baccalauréat) sont vécus de l'intérieur ; ma mort, elle, est la limite de ma vie, non un événement à l'intérieur de ma vie. On ne peut pas se représenter son propre néant, car toute représentation suppose un sujet vivant qui représente.

\begin{exoresolu}[Analyser l'argument d'Épicure]
Énoncé : Un élève affirme : « Puisque la mort est certaine, il est naturel et raisonnable d'en avoir peur. » En t'appuyant sur la thèse d'Épicure, montre que cette affirmation peut être contestée.
\tcblower
Solution détaillée. L'élève raisonne ainsi : la mort est un mal certain, donc la craindre est justifié. Épicure conteste la prémisse cachée : la mort est-elle vraiment un mal \emph{pour nous} ? Or un mal ne peut exister que pour quelqu'un qui le subit. Deux cas se présentent : soit je suis vivant, et alors la mort n'existe pas encore, donc il n'y a rien à craindre ; soit je suis mort, et alors je n'existe plus, donc il n'y a plus personne pour subir un mal. Dans les deux cas, la mort n'est jamais un mal présent pour un sujet. Conclusion : la peur de la mort porte sur un objet qui ne nous atteindra jamais ; elle est donc irrationnelle. La sagesse consiste à se libérer de cette crainte pour jouir de la vie présente. (On peut objecter : ce qui fait peur, ce n'est pas la mort elle-même, mais la souffrance qui précède, ou la perte des êtres aimés — objections qui montrent les limites de l'argument.)
\end{exoresolu}

\courssub{La conscience de la mort : une spécificité humaine}

Tous les êtres vivants meurent : l'animal meurt, la plante meurt, l'homme meurt. Mais une différence essentielle sépare l'homme de l'animal : \emph{l'animal meurt sans le savoir ; l'homme sait qu'il mourra}. La chèvre ne prévoit pas ses funérailles ; le vieillard de Bafoussam prépare sa succession, choisit parfois son linceul, réconcilie ses enfants. Cette conscience anticipée de la mort est proprement humaine.

\begin{propriete}[Thèse de Heidegger : l'être-pour-la-mort]
Heidegger (1889--1976), dans \emph{Être et Temps}, définit l'existence humaine comme un \cle{être-pour-la-mort} (\emph{Sein-zum-Tode}). La mort n'est pas un accident qui surviendra un jour : elle est la possibilité la plus propre, certaine et indépassable de l'homme, présente à chaque instant de sa vie. Deux attitudes sont possibles : fuir cette vérité dans les divertissements et les bavardages (existence \emph{inauthentique}), ou anticiper lucidement sa mort (existence \emph{authentique}). Celui qui accepte sa mortalité cesse de gaspiller son temps : chaque projet devient urgent et précieux. La conscience de la mort ne paralyse donc pas : elle \emph{donne du sérieux et de la valeur à la vie}.
\end{propriete}

Cette analyse rejoint une sagesse populaire camerounaise : « celui qui sait qu'il partira prépare bien son voyage ». Le proverbe ne dit pas d'avoir peur du départ, mais de le laisser éclairer la conduite de la vie.

\courssub{La mort comme fait social : la diversité culturelle au Cameroun}

La mort n'est pas seulement un événement biologique individuel : c'est un \emph{fait social total}. Chaque société organise la mort par des rites, qui montrent ce qu'elle pense de l'existence.

\begin{itemize}
\item \textbf{Dans les traditions chrétiennes} (très présentes au Centre et au Littoral) : messe de requiem, prières pour le repos de l'âme, croyance en la résurrection. Le deuil est marqué par l'espérance de la vie éternelle.
\item \textbf{Dans les traditions musulmanes} (Nord, Adamaoua) : toilette rituelle du corps, prière mortuaire, enterrement rapide, souvent le jour même ; la mort est soumission à la volonté d'Allah et passage vers le jugement.
\item \textbf{Dans les traditions coutumières} : le \emph{cry-die} dans l'Ouest et le Nord-Ouest (veillées funèbres où l'on pleure, chante et danse pour honorer le défunt), la \emph{levée de deuil} qui clôt officiellement la période de tristesse et réintègre la famille dans la vie sociale, les funérailles traditionnelles qui peuvent durer plusieurs jours et mobiliser tout le village.
\end{itemize}

Ces rites montrent que la mort concerne toute la communauté : elle menace le groupe, qui doit réaffirmer son unité. Le deuil n'est pas seulement une affaire privée : c'est une cérémonie collective, avec des obligations, des interdits, des dépenses, des retrouvailles familiales.

\begin{savaistu}[Les crânes des ancêtres chez les Bamiléké]
Dans certaines traditions bamiléké (Ouest-Cameroun), les crânes des ancêtres importants étaient conservés, honorés et consultés lors de cérémonies. Loin d'être macabre, cette pratique exprime une conviction profonde : la mort ne coupe pas le lien entre les vivants et les défunts. Les ancêtres demeurent membres de la famille, protecteurs et intercesseurs. La mort est ainsi pensée comme un \emph{changement de statut}, non comme une disparition : le défunt passe du rang de vivant au rang d'ancêtre.
\end{savaistu}

\courssub{Synthèse : trois dimensions de la mort}

\begin{popfigure}
\begin{center}
\begin{tikzpicture}[font=\small]
\node[draw=popblue, fill=popblueL, rounded corners, text width=3.6cm, align=center, minimum height=3.4cm] (bio) at (0,0) {\textbf{Mort biologique}\\[2pt] Cessation irréversible des fonctions vitales\\[2pt] \emph{Critères} : arrêt du cœur, du cerveau, de la respiration\\[2pt] \emph{Ex.} : constat médical du décès};
\node[draw=poporange, fill=poporangeL, rounded corners, text width=3.6cm, align=center, minimum height=3.4cm] (phi) at (5.2,0) {\textbf{Mort philosophique}\\[2pt] Limite absolue de l'existence\\[2pt] \emph{Enjeu} : néantisation ou passage ? Certitude, angoisse, sagesse\\[2pt] \emph{Ex.} : Épicure, Heidegger};
\node[draw=popgreen, fill=popgreenL, rounded corners, text width=3.6cm, align=center, minimum height=3.4cm] (cult) at (10.4,0) {\textbf{Mort culturelle}\\[2pt] Fait social organisé par des rites\\[2pt] \emph{Critères} : deuil, funérailles, statut d'ancêtre\\[2pt] \emph{Ex.} : cry-die, levée de deuil, prières};
\end{tikzpicture}
\end{center}
\legende{Les trois dimensions de la mort : biologique (le corps cesse de fonctionner), philosophique (l'existence atteint sa limite), culturelle (la société organise le passage du défunt).}
\end{popfigure}

\begin{aretenir}[L'essentiel]
\begin{itemize}
\item La mort (latin \emph{mors}) est la cessation définitive de la vie ; les langues camerounaises la nomment souvent comme un « retour aux ancêtres ».
\item Biologiquement : arrêt irréversible des fonctions vitales (cœur, cerveau, respiration).
\item Philosophiquement : limite absolue de l'existence, pensée comme néantisation (Épicure, Sartre) ou comme passage (Platon, religions, traditions africaines).
\item La mort est une certitude (syllogisme : tous les hommes sont mortels) mais un événement impossible à vivre (Épicure, Wittgenstein).
\item Seul l'homme \emph{sait} qu'il mourra : Heidegger parle d'« être-pour-la-mort » ; anticiper sa mort rend l'existence authentique.
\item Au Cameroun, la mort est un fait social : deuils chrétiens, musulmans et traditionnels (cry-die, levée de deuil) montrent que le défunt reste lié à la communauté.
\end{itemize}
\end{aretenir}

Ainsi définie, la mort apparaît non comme un simple fait médical, mais comme le problème central de l'existence humaine : c'est parce qu'elle est certaine, insaisissable et chargée de sens que se pose la question suivante — celle du \emph{sens de l'existence} elle-même, que nous examinerons dans la prochaine section.

\coursec{Le sens de l'existence humaine : les grandes conceptions}

Nous venons de définir la mort comme la cessation irréversible de la vie et, pour l'homme, comme un événement pensé d'avance : l'être humain est le seul vivant qui \emph{sait} qu'il va mourir. Mais cette certitude soulève immédiatement une question plus profonde encore : si je dois mourir, \emph{pour quoi vivre} ? Ma vie a-t-elle un sens, une direction, une valeur ? C'est à cette interrogation, la plus personnelle et la plus universelle de toutes, que répondent les grandes conceptions philosophiques et religieuses que nous allons maintenant étudier.

\courssub{La question du sens : pourquoi y a-t-il quelque chose plutôt que rien ?}

Commençons par l'intuition la plus simple. Un enfant demande : « Pourquoi le monde existe-t-il ? Pourquoi suis-je là ? » Ces questions enfantines sont en réalité les questions les plus vertigineuses de la philosophie. Le philosophe allemand \cle{Leibniz} (1646-1716) les a formulées avec une précision restée célèbre : « Pourquoi y a-t-il \emph{quelque chose} plutôt que \emph{rien} ? » Rien n'obligeait en effet l'univers à exister ; il aurait pu n'y avoir que le néant. Que quelque chose existe, et que des êtres conscients existent pour s'en étonner, voilà qui appelle une explication.

Cette question cosmique se dédouble en une question personnelle : non plus seulement « pourquoi le monde existe-t-il ? », mais « \cle{pourquoi vivre} ? », c'est-à-dire « pour \emph{quoi} vivre ? ». Remarquons la nuance entre les deux formulations : « pourquoi vivre ? » interroge la \emph{cause} (qu'est-ce qui me fait vivre ?), tandis que « pour quoi vivre ? » interroge la \emph{fin}, le but, la valeur qui justifie l'existence. C'est à cette seconde question que répond la notion de sens de l'existence.

\begin{definition}[Le sens de l'existence]
Le \cle{sens de l'existence} désigne la réponse à la question « pour quoi vivre ? » : c'est la \emph{finalité} (le but à atteindre) ou la \emph{valeur} (ce qui mérite qu'on vive pour elle) qui justifie la vie humaine et lui donne une orientation. Dire que l'existence « a un sens », c'est affirmer qu'elle n'est pas un simple fait biologique, mais qu'elle vaut d'être vécue en vue de quelque chose : Dieu, le bonheur, la liberté, l'amour, la réalisation de soi.
\end{definition}

Toutes les conceptions que nous allons examiner partagent le même point de départ — la question du sens — mais divergent radicalement sur la réponse. C'est précisément ce qui fait l'intérêt philosophique de leur comparaison.

\courssub{La conception religieuse : un sens donné par Dieu}

Pour les grandes traditions religieuses, la question du sens reçoit une réponse claire et rassurante : l'existence humaine a un sens parce qu'elle est \emph{voulue par Dieu}. L'homme n'est pas le produit du hasard ; il est une créature, dotée d'une âme, placée sur terre pour accomplir une mission.

Dans le \cle{christianisme}, la vie terrestre est un chemin vers Dieu : « Aimez-vous les uns les autres » résume la loi qui donne sens à l'existence, et la mort n'est pas une fin mais un \emph{passage} vers la vie éternelle, où chacun sera jugé selon ses actes. Dans l'\cle{islam}, la vie est une épreuve (\emph{ibtila'}) : le croyant est sur terre pour adorer Dieu et accomplir le bien, et la mort ouvre sur l'au-delà, le paradis ou l'enfer selon les œuvres. Le Coran affirme : « Toute âme goûtera la mort » (sourate III, 185), mais cette mort n'est qu'un retour vers Dieu.

Les \cle{religions traditionnelles africaines}, très présentes au Cameroun, proposent une conception originale et profonde : la mort ne rompt pas le lien entre les vivants et les défunts. Les morts deviennent des \emph{ancêtres}, qui continuent de veiller sur la communauté, d'intercéder pour elle et d'exiger le respect des traditions. Le sens de l'existence est alors \emph{communautaire} : vivre, c'est assurer la continuité du lignage, honorer les ancêtres et transmettre à son tour. Les rites funéraires (deuils, cérémonies, cultes des crânes chez certains peuples) manifestent cette conviction : le défunt n'a pas disparu, il a changé de statut.

\begin{exemplebox}[Face au deuil : trois attitudes]
Un jeune homme de Yaoundé perd son père. Le \textbf{croyant} trouve consolation dans la foi : son père est « auprès de Dieu » ou « parmi les ancêtres », la séparation n'est pas définitive. L'\textbf{existentialiste} se dit : mon père a accompli sa vie par ses choix et ses actes ; honorer sa mémoire, c'est poursuivre ce qu'il a construit. L'\textbf{absurdiste} accepte lucidement que la mort est sans retour, sans consolation promise, et décide de vivre plus intensément le temps qui reste. Même épreuve, trois réponses différentes à la question du sens.
\end{exemplebox}

\courssub{La conception existentialiste : l'homme crée lui-même le sens}

Contre la conception religieuse, le philosophe français \cle{Jean-Paul Sartre} (1905-1980) soutient une thèse radicale, résumée par sa formule célèbre : « \emph{l'existence précède l'essence} ». Que signifie-t-elle ? Pour un objet fabriqué, comme un coupe-papier, l'essence (le plan, la définition) précède l'existence : l'artisan sait ce qu'il veut faire avant de le faire. Pour l'homme, soutient Sartre, c'est l'inverse : l'homme \emph{existe d'abord}, il naît sans définition préalable, sans nature fixée par un Dieu créateur, et c'est ensuite, par ses choix et ses actes, qu'il se définit. « L'homme n'est rien d'autre que ce qu'il se fait. »

Conséquence décisive pour notre problème : l'existence n'a \emph{pas de sens donné d'avance}. Il n'existe aucun plan divin, aucune destinée écrite, aucune nature humaine qui nous dirait ce que nous devons être. C'est donc à chacun de \emph{créer} le sens de sa vie, par sa \cle{liberté} et par son \cle{engagement} — c'est-à-dire par des choix assumés qui engagent non seulement soi-même, mais une certaine image de l'homme. Cette liberté est à la fois une grandeur et un vertige : nous sommes, dit Sartre, « condamnés à être libres », responsables de tout ce que nous sommes sans pouvoir nous excuser sur Dieu, la société ou le destin.

Face à la mort, l'existentialiste n'attend aucun au-delà : la mort est simplement la limite qui clôt ma vie et la transforme en totalité achevée. C'est donc \emph{ici et maintenant}, dans le temps qui m'est donné, que je dois faire de ma vie une œuvre.

\courssub{La conception de l'absurde : Camus et Sisyphe}

\cle{Albert Camus} (1913-1960), dans \emph{Le Mythe de Sisyphe} (1942), part d'un constat encore plus radical : l'homme réclame du sens, de la clarté, de la justice — mais le monde reste silencieux, indifférent, sans réponse. Ce décalage, Camus le nomme l'\cle{absurde}.

\begin{definition}[L'absurde selon Camus]
L'\cle{absurde} n'est ni dans l'homme seul, ni dans le monde seul : il est le \emph{divorce}, la confrontation entre l'\emph{exigence humaine de sens} (notre soif de comprendre, de durée, de justice) et le \emph{silence déraisonnable du monde} qui ne répond pas. La mort, qui rend toutes nos entreprises provisoires, est la manifestation la plus évidente de l'absurde.
\end{definition}

Que faire alors ? Camus écarte deux solutions qu'il juge malhonnêtes. D'abord le \emph{suicide} : détruire sa vie parce qu'elle est absurde, c'est se laisser vaincre par l'absurde au lieu de l'affronter. Ensuite ce qu'il appelle le « suicide philosophique », c'est-à-dire le \emph{saut dans la foi} : croire en un au-delà pour échapper au vertige, c'est fuir la lucidité. Reste une troisième voie, la seule digne selon lui : la \cle{révolte}.

\begin{propriete}[Thèse de Camus : la révolte lucide]
Face à l'absurde, la réponse n'est ni le suicide (fuite par la mort) ni l'espoir religieux (fuite par la foi), mais la \cle{révolte lucide} : vivre \emph{sans appel}, c'est-à-dire sans recours à un au-delà, en maintenant les yeux ouverts sur la conscience de l'absurde, et en donnant à la vie le maximum d'intensité, d'expériences et de passion. Ce qui compte n'est pas de vivre le mieux, mais de vivre \emph{le plus}. (Cette thèse est exigible au Baccalauréat technique : il faut savoir la formuler et l'illustrer par Sisyphe.)
\end{propriete}

L'image de cette thèse est le héros mythique \cle{Sisyphe}, condamné par les dieux à pousser éternellement un rocher jusqu'au sommet d'une montagne, d'où il retombe sans cesse. Sisyphe est l'image de la condition humaine : nos efforts semblent vains, tout est à recommencer. Pourtant, au moment où Sisyphe redescend chercher son rocher, il est \emph{conscient} de sa condition — et cette lucidité est sa victoire sur les dieux. « La lutte elle-même vers les sommets suffit à remplir un cœur d'homme. \emph{Il faut imaginer Sisyphe heureux}. »

\begin{popfigure}
\begin{center}
\begin{tikzpicture}[scale=1.0]
  % Montagne (pente)
  \draw[thick,popdark] (0,0) -- (5.5,3.2) -- (7.5,3.2);
  \draw[thick,popdark,dashed] (5.5,3.2) -- (5.5,0);
  \draw[thick,popdark] (0,0) -- (7.5,0);
  % Rocher en haut de la pente
  \draw[fill=poporangeL,draw=poporange,thick] (4.6,2.55) circle (0.45);
  % Sisyphe stylisé poussant
  \draw[thick,popblue] (3.55,1.95) circle (0.14); % tête
  \draw[thick,popblue] (3.62,1.82) -- (3.95,1.35); % corps
  \draw[thick,popblue] (3.95,1.35) -- (3.7,0.95); % jambe arrière
  \draw[thick,popblue] (3.95,1.35) -- (4.25,0.98); % jambe avant
  \draw[thick,popblue] (3.75,1.62) -- (4.25,2.15); % bras vers rocher
  % Flèche montée
  \draw[-{Stealth},very thick,popgreen] (1.0,0.35) -- (4.0,2.1);
  \node[popgreen,font=\small] at (2.1,1.7) {effort vers le sommet};
  % Flèche descente du rocher
  \draw[-{Stealth},very thick,popink] (5.3,3.5) .. controls (4.2,4.3) and (2.2,4.3) .. (1.2,3.4);
  \node[popink,font=\small] at (3.2,4.45) {le rocher retombe};
  % Sisyphe heureux
  \node[popdark,font=\small\itshape,align=center] at (6.3,1.3) {« Il faut imaginer\\Sisyphe heureux »\\(Camus)};
\end{tikzpicture}
\end{center}
\legende{Le mythe de Sisyphe : l'homme pousse son rocher (effort, projets, travail), le rocher retombe (échec, mort, recommencement). La conscience lucide de ce cycle, loin d'écraser Sisyphe, fait sa grandeur : c'est la révolte heureuse de Camus.}
\end{popfigure}

\courssub{Conceptions complémentaires : hédonisme, stoïcisme, vitaliste}

\textbf{La conception hédoniste.}\par
Pour \cle{Épicure} (341-270 av. J.-C.), le \emph{bonheur} est la fin de l'existence, et le bonheur consiste dans le plaisir — mais un plaisir sage : l'absence de douleur du corps et de trouble de l'âme (\emph{ataraxie}). Quant à la mort, Épicure la dissipe par un argument célèbre : « La mort n'est rien pour nous, car tant que nous sommes, la mort n'est pas, et quand la mort est là, nous ne sommes plus. » Inutile donc de la craindre.

\textbf{La conception stoïcienne.}\par
Pour les stoïciens (\cle{Épictète}, \cle{Marc Aurèle}), le sens de la vie est la \emph{vertu}, c'est-à-dire l'accord avec la raison et l'ordre du monde. Leur règle pratique est simple : distinguer ce qui \emph{dépend de nous} (nos jugements, nos désirs, nos actes) de ce qui \emph{n'en dépend pas} (la maladie, la richesse, la mort). Il faut vouloir le premier et \emph{accepter} sereinement le second. La mort, ne dépendant pas de nous, doit être accueillie sans révolte ni terreur.

\textbf{La conception vitaliste.}\par
\cle{Nietzsche} (1844-1900) voit dans la vie une \emph{volonté de puissance} : non pas la soif de dominer les autres, mais l'élan de la vie qui veut croître, se dépasser, créer des valeurs. Le sens de l'existence n'est pas à recevoir mais à \emph{inventer}, en disant oui à la vie tout entière, y compris à la souffrance : c'est la formule de l'\emph{amor fati}, l'amour du destin — ne rien vouloir changer de ce qui arrive, mais le vouloir comme on le vit.

\begin{popfigure}
\begin{center}
\begin{tikzpicture}
  \node[font=\small\bfseries,popdark] (q) at (0,0) {La vie a-t-elle un sens ?};
  \node[draw=popblue,fill=popblueL,rounded corners,align=center,font=\scriptsize] (rel) at (-5.2,-1.8) {\textbf{Religieuse}\\Bible, Coran, traditions\\africaines};
  \node[draw=poppurple,fill=poppurpleL,rounded corners,align=center,font=\scriptsize] (exi) at (-1.75,-1.8) {\textbf{Existentialiste}\\Sartre};
  \node[draw=poporange,fill=poporangeL,rounded corners,align=center,font=\scriptsize] (abs) at (1.75,-1.8) {\textbf{Absurde}\\Camus};
  \node[draw=popgreen,fill=popgreenL,rounded corners,align=center,font=\scriptsize] (vit) at (5.2,-1.8) {\textbf{Hédoniste, stoïcienne,}\\ \textbf{vitaliste}\\Épicure, Épictète, Nietzsche};
  \draw[-{Stealth},thick,popline] (q) -- (rel);
  \draw[-{Stealth},thick,popline] (q) -- (exi);
  \draw[-{Stealth},thick,popline] (q) -- (abs);
  \draw[-{Stealth},thick,popline] (q) -- (vit);
  \node[align=center,font=\scriptsize,popdark] at (-5.2,-3.6) {Sens \emph{donné} par Dieu ;\\la mort est un passage\\vers l'au-delà (ancêtres)};
  \node[align=center,font=\scriptsize,popdark] at (-1.75,-3.6) {Pas de sens donné :\\l'homme le \emph{crée} par\\sa liberté, son engagement};
  \node[align=center,font=\scriptsize,popdark] at (1.75,-3.6) {Le monde est muet :\\révolte lucide, vivre\\sans appel ni espoir};
  \node[align=center,font=\scriptsize,popdark] at (5.2,-3.6) {Bonheur, vertu,\\volonté de puissance ;\\accepter la mort (\emph{amor fati})};
\end{tikzpicture}
\end{center}
\legende{Carte des grandes conceptions du sens de l'existence : même question, quatre familles de réponses.}
\end{popfigure}

\begin{center}
\begin{tabularx}{\linewidth}{|l|l|X|X|}
\hline
\rowcolor{popblueL} \textbf{Conception} & \textbf{Représentant} & \textbf{Sens de l'existence} & \textbf{Rapport à la mort} \\
\hline
Religieuse & Christianisme, islam, traditions africaines & Donné par Dieu ; vivre selon sa volonté, honorer les ancêtres & Passage vers l'au-delà ; survie de l'âme ou des ancêtres \\
\hline
Existentialiste & Sartre & Aucun sens donné : à créer par la liberté et l'engagement & Fin sans au-delà ; limite qui oblige à agir ici et maintenant \\
\hline
Absurde & Camus & Aucun sens ultime, mais révolte lucide et intensité de vie & Constatée sans espoir ; elle rend la vie précieuse \\
\hline
Hédoniste / stoïcienne / vitaliste & Épicure, Épictète, Nietzsche & Bonheur, vertu, ou création de valeurs (\emph{amor fati}) & À ne pas craindre, à accepter, ou à aimer comme destin \\
\hline
\end{tabularx}
\end{center}

\courssub{Discussion critique : la conscience de la mort donne-t-elle du sens ?}

Nous pouvons maintenant confronter ces thèses à une question décisive : la conscience que nous avons de notre mort donne-t-elle du sens à l'existence, ou au contraire le lui retire-t-elle ?

\textbf{Première thèse : la mort donne du sens.} Parce que le temps est compté, chaque instant devient précieux. Si nous étions immortels, nous pourrions tout remettre à demain : rien n'aurait d'urgence ni de prix. La mort fonctionne comme une limite qui concentre la vie : elle crée l'\emph{urgence de vivre}, la valeur du temps, le sérieux de nos choix. Un examen n'a de sens que parce qu'il a une date ; une vie n'a de prix que parce qu'elle est unique et brève. C'est l'intuition de Sartre (la finitude oblige à choisir) et de Camus (l'absurde invite à vivre le plus).

\textbf{Seconde thèse : la mort retire tout sens.} À quoi bon bâtir, aimer, apprendre, si tout doit être emporté ? « Vanité des vanités, tout est vanité », dit l'Ecclésiaste. La mort semble niveler toutes les différences : le riche et le pauvre, le savant et l'ignorant finissent pareils. Nos œuvres elles-mêmes périront. De ce point de vue, la conscience de la mort est un poison qui gâte toute joie : c'est le vertige de la vanité de toutes choses.

\textbf{Dépassement.} Les deux thèses saisissent chacune une part de vérité, car elles décrivent deux \emph{attitudes} possibles devant le même fait. La mort est en elle-même indifférente : c'est nous qui en faisons soit un néant qui avale tout, soit une limite qui donne son prix à chaque instant. Camus montre précisément qu'on peut passer de la seconde attitude à la première sans tricher : reconnaître que tout est provisoire, et conclure non pas « donc rien ne vaut la peine », mais « donc tout est à vivre maintenant ». La brièveté de la vie n'enlève sa valeur qu'à celui qui confond valeur et éternité ; or une rose ne vaut pas moins parce qu'elle fane : elle vaut peut-être davantage.

\begin{methode}[Comparer deux thèses philosophiques]
Pour comparer des conceptions (exercice fréquent au Baccalauréat), procéder en trois étapes :
\begin{enumerate}
  \item \textbf{Identifier le point d'accord} : toutes les thèses posent la même question (ici : la vie a-t-elle un sens, et la mort la vide-t-elle de valeur ?).
  \item \textbf{Identifier le point de désaccord} : la réponse diffère (sens donné par Dieu / sens créé par la liberté / absence de sens mais révolte / bonheur et acceptation).
  \item \textbf{Évaluer} : confronter chaque réponse à une objection ou à une situation concrète (le deuil, l'échec, la maladie), puis conclure par une position personnelle justifiée.
\end{enumerate}
\end{methode}

\begin{attention}[Pièges fréquents]
\begin{itemize}
  \item Ne pas confondre \emph{absurde} et \emph{nihilisme} : Camus ne dit pas « rien ne vaut la peine », il dit que le monde ne donne pas de sens, mais que l'homme peut vivre pleinement malgré tout. Le nihilisme détruit toute valeur ; la révolte de Camus en affirme (la vie, la solidarité).
  \item Ne pas résumer Sartre à « fais ce que tu veux » : la liberté existentialiste est inséparable de la \emph{responsabilité} et de l'engagement.
  \item Ne pas opposer caricaturalement religion et philosophie : la conception religieuse est elle aussi une réponse argumentée à la question du sens, qu'il faut exposer avec exactitude avant de la discuter.
\end{itemize}
\end{attention}

\begin{savaistu}[Albert Camus, une vie à l'image de sa philosophie]
Albert Camus, né en Algérie en 1913 dans une famille pauvre, reçoit le prix Nobel de littérature en 1957, à 44 ans. Trois ans plus tard, le 4 janvier 1960, il meurt brutalement dans un accident de voiture, à 46 ans — lui qui détestait l'automobile et avait prévu de rentrer en train. Dans sa poche, on trouva son billet de train non utilisé. Sa mort absurde, imprévisible, sans raison, illustre tragiquement ce qu'il avait philosophé toute sa vie : la fragilité de l'existence et l'urgence de vivre. « Il faut imaginer Sisyphe heureux » reste l'une des phrases les plus citées de la philosophie du XX\textsuperscript{e} siècle.
\end{savaistu}

\begin{exoresolu}[Exercice d'application : comparer et discuter]
\textbf{Énoncé.} Un élève déclare : « Puisque nous allons tous mourir, étudier, travailler et se marier ne servent à rien. » (a) À quelle conception du sens cette parole ressemble-t-elle ? (b) Montrez, en mobilisant deux auteurs étudiés, qu'on peut répondre à cet élève. (c) Formulez votre position personnelle en trois arguments.
\tcblower
\textbf{Solution détaillée.} (a) Cette parole exprime la thèse de la \emph{vanité de tout} : la mort retirerait tout sens à l'existence. Elle se rapproche du nihilisme, non de l'absurde de Camus. (b) Première réponse, \emph{Camus} : le raisonnement de l'élève est exactement celui que Camus réfute dans \emph{Le Mythe de Sisyphe} ; de l'absence de sens donné, il ne faut pas conclure au renoncement (suicide de l'action), mais à la révolte : c'est parce que la vie est brève qu'il faut la vivre avec le maximum d'intensité — « il faut imaginer Sisyphe heureux ». Seconde réponse, \emph{Sartre} : l'existence n'a pas de sens donné d'avance, mais c'est une chance : par ses études, son travail, ses engagements, l'élève \emph{crée} lui-même le sens de sa vie ; renoncer, c'est encore un choix, mais un choix dont il est responsable. (c) Position personnelle possible : la mort est une limite qui donne son prix au temps (argument 1 : valeur de la rareté) ; le sens ne se trouve pas, il se construit par des projets (argument 2 : Sartre) ; l'expérience montre que ceux qui agissent et s'engagent trouvent leur vie digne d'être vécue, témoin Sisyphe heureux (argument 3 : Camus).
\end{exoresolu}

\begin{exercice}[Synthèse personnelle guidée]
Rédigez un paragraphe argumenté (15 à 20 lignes) répondant à la question : « La conscience de la mort donne-t-elle un sens à la vie ? » Consignes : (1) formulez une thèse claire dès la première phrase ; (2) appuyez-vous sur au moins deux conceptions étudiées (religieuse, existentialiste, absurde, stoïcienne ou vitaliste) ; (3) illustrez par un exemple de la vie courante ou camerounaise ; (4) répondez à une objection avant de conclure.
\end{exercice}

\begin{corrige}[Exercice : synthèse personnelle guidée]
Éléments de correction (toute position cohérente est admissible si elle est argumentée). \emph{Introduction de la thèse} : « La conscience de la mort, loin de vider la vie de son sens, la lui donne, car elle rend le temps précieux et les choix sérieux. » \emph{Argument 1} : la rareté crée la valeur — ce qui ne finit jamais n'a pas de prix ; Sartre montre que la finitude oblige à choisir et donc à se définir. \emph{Argument 2} : Camus — face au silence du monde, la révolte lucide transforme la brièveté de la vie en appel à l'intensité ; « il faut imaginer Sisyphe heureux ». \emph{Exemple} : un candidat au Baccalauréat travaille avec sérieux précisément parce que l'examen a une date ; de même, un fils profite de son père vivant parce qu'il sait le temps compté. \emph{Objection} : « mais la mort détruit tout ce que nous construisons » — réponse : elle détruit les œuvres, non la valeur de les avoir accomplies ; la valeur d'une vie tient à sa qualité, non à sa durée (on peut citer aussi la tradition africaine : le mort accompli devient ancêtre et continue d'agir dans la communauté). \emph{Conclusion} : la mort est une limite, et la limite est la condition du sens, comme le cadre est la condition du tableau.
\end{corrige}

\begin{aretenir}[L'essentiel de la section]
\begin{itemize}
  \item Le \cle{sens de l'existence} répond à la question « pour quoi vivre ? » ; Leibniz pose la question radicale : « pourquoi y a-t-il quelque chose plutôt que rien ? ».
  \item Conception \emph{religieuse} : sens donné par Dieu, la mort est un passage (au-delà, ancêtres). Conception \emph{existentialiste} (Sartre) : l'existence précède l'essence ; le sens se crée par la liberté et l'engagement.
  \item Conception de l'\emph{absurde} (Camus) : divorce entre notre soif de sens et le silence du monde ; réponse : la révolte lucide, « il faut imaginer Sisyphe heureux ».
  \item Compléments : Épicure (bonheur, ne pas craindre la mort), les stoïciens (accepter ce qui ne dépend pas de nous), Nietzsche (volonté de puissance, \emph{amor fati}).
  \item Discussion : la conscience de la mort peut sembler tout rendre vain, mais elle donne en réalité son prix au temps et son urgence à la vie.
\end{itemize}
\end{aretenir}

\coursec{Application aux situations concrètes de la vie courante}

Dans la section précédente, nous avons parcouru les grandes conceptions du sens de l'existence humaine : la vie comme don de Dieu dans les traditions religieuses, l'existence comme projet libre chez les existentialistes, la vie comme absurdité à assumer chez Camus, la vie comme accomplissement du bonheur chez Aristote. Mais une philosophie qui resterait dans les livres serait une philosophie morte. L'examen — et surtout la vie — exige que nous sachions \emph{appliquer} ces notions à des situations réelles. C'est précisément le savoir-faire visé par le programme : appliquer les notions de l'unité à des situations concrètes de la vie courante, puis rédiger et justifier une démarche, une réponse ou une conclusion. Nous allons donc traverser cinq situations typiques, toutes puisées dans notre quotidien camerounais ou dans le débat mondial, avant de dégager une méthode générale d'application.

\courssub{Situation 1 : les funérailles au Cameroun — un rite qui donne sens à la mort}

\textbf{Description du fait.} Au Cameroun, les funérailles ne sont pas un simple enterrement : ce sont des cérémonies qui mobilisent la famille élargie, le village, les associations (tontines, réunions de quartier), parfois pendant plusieurs semaines. On organise des veillées, on collecte des contributions, on prépare le deuil, puis souvent, des mois plus tard, la « levée de deuil ». Un décès à Yaoundé peut déplacer des centaines de personnes vers le village d'origine du défunt.

\textbf{Identification du concept.} La notion en jeu est celle de \cle{mort}, mais comprise non comme simple événement biologique (arrêt du cœur et du cerveau), mais comme \emph{événement social et symbolique}. La mort biologique est un fait ; la mort humaine est un sens.

\textbf{Mobilisation d'une thèse.} Les anthropologues et les philosophes (par exemple Philippe Ariès, historien de la mort en Occident) montrent que les rites funéraires ont une double fonction. D'une part, ils \emph{donnent sens à la mort} : en inscrivant le défunt dans une continuité (les ancêtres, la communauté, Dieu), ils empêchent que la mort apparaisse comme une absurdité pure. D'autre part, ils \emph{consolident les vivants} : la solidarité matérielle et affective déployée lors des funérailles ressoude le groupe fragilisé par la perte d'un des siens.

\textbf{Conclusion justifiée.} Les funérailles camerounaises ne sont donc pas une « dépense inutile », comme le prétendent parfois certains discours économiques simplistes : elles accomplissent un travail philosophique essentiel — transformer la mort, événement qui détruit le sens, en événement qui en produit. On peut toutefois \emph{nuancer} : lorsque le faste des cérémonies endette les familles pendant des années, la fonction symbolique risque d'être détournée en compétition sociale. L'analyse philosophique doit savoir tenir les deux regards.

\begin{savaistu}[Le « cry-die » : danser la mort pour l'accepter]
Dans les Grassfields de l'Ouest Cameroun, certaines cérémonies funéraires prennent la forme du \emph{cry-die} (littéralement « pleurer-mourir ») : un deuil dansé, où les pleurs alternent avec la musique, les masques et la danse. Loin d'être un manque de respect, cette pratique traduit une sagesse pratique profonde : la mort d'un vieillard qui a accompli sa vie n'est pas seulement une perte, c'est aussi un passage, une naissance au monde des ancêtres. Transformer la tristesse en célébration, c'est refuser que la mort ait le dernier mot sur le sens de l'existence. Les philosophes stoïciens, à l'autre bout du monde, tenaient un discours comparable : accepter ce qui ne dépend pas de nous est le commencement de la sérénité.
\end{savaistu}

\courssub{Situation 2 : le jeune diplômé technique face au chômage — l'existence comme projet}

\textbf{Description du fait.} Un jeune titulaire d'un Baccalauréat technique ou d'un BTS — en génie électrique, en comptabilité, en maintenance — ne trouve pas d'emploi salarié après deux ans de recherche. Deux attitudes s'offrent à lui. La première : « C'est le pays qui est comme ça, sans piston on n'a rien, je n'y peux rien. » La seconde : « Je ne contrôle pas le marché de l'emploi, mais je contrôle ce que je fais de mes compétences : je peux me former davantage, créer une petite activité, apprendre un métier complémentaire. »

\textbf{Identification du concept.} C'est la notion d'\cle{existence} qui est en jeu, et plus précisément la question : l'existence humaine est-elle un destin subi ou un projet construit ?

\textbf{Mobilisation des thèses.} Le \emph{fatalisme} soutient que tout est écrit d'avance : la volonté humaine est impuissante face aux circonstances. Cette thèse a une part de vérité — nous ne choisissons ni notre naissance, ni notre pays, ni le contexte économique. Mais l'\emph{existentialisme} de Sartre répond que « l'existence précède l'essence » : l'homme n'est rien d'autre que ce qu'il se fait. Même dans des conditions difficiles, il reste responsable de l'usage qu'il fait de sa liberté. Entre les deux, le stoïcisme propose une sagesse réaliste : distinguer ce qui dépend de nous (nos efforts, nos décisions) de ce qui n'en dépend pas (la conjoncture), et concentrer son énergie sur le premier.

\textbf{Conclusion justifiée.} Le fatalisme pur est philosophiquement intenable, car il se réfute lui-même : celui qui dit « je n'y peux rien » et cesse d'agir \emph{produit} par sa résignation l'échec qu'il prétendait subir. Mais la volonté toute-puissante est un mythe symétrique : nier les obstacles réels (crise économique, inégalités) serait injuste envers ceux qui les subissent. La position justifiée est donc une \emph{liberté située} : l'existence est un projet, mais un projet qui se construit avec des matériaux qu'on n'a pas choisis. Le jeune diplômé qui crée un atelier de réparation de téléphones à Douala n'a pas choisi le chômage ; il choisit sa réponse au chômage — et c'est là que son existence prend sens.

\courssub{Situation 3 : les accidents de la route — la mort comme risque banalisé}

\textbf{Description du fait.} Les accidents de la circulation tuent chaque année au Cameroun plus de \num{1000} personnes selon les estimations officielles du Ministère des Transports — certaines sources internationales avancent des chiffres encore plus élevés. Pourtant, qui n'a jamais vu un motocycliste sans casque, un conducteur téléphoner au volant, un car de \num{70} places bondé de \num{100} passagers foncer sur l'axe Yaoundé–Douala ?

\begin{popfigure}
\begin{center}
\begin{tikzpicture}
\begin{axis}[
    ybar, bar width=18pt,
    width=0.85\linewidth, height=6cm,
    xlabel={Année}, ylabel={Nombre de décès (indicatif)},
    symbolic x coords={2018,2019,2020,2021,2022,2023},
    xtick=data,
    ymin=0, ymax=1600,
    nodes near coords, nodes near coords style={font=\small},
    every axis plot/.append style={fill=popblue!70, draw=popdark},
    grid=major, grid style={popline!50},
]
\addplot[popcurve,fill=popblue!70] coordinates {(2018,1050) (2019,1120) (2020,980) (2021,1150) (2022,1240) (2023,1300)};
\end{axis}
\end{tikzpicture}
\end{center}
\legende{Évolution indicative du nombre de décès liés aux accidents de la route au Cameroun (données indicatives, d'après les estimations du Ministère des Transports et des rapports de sécurité routière). La baisse de 2020 s'explique par les restrictions de circulation ; la tendance de fond reste à la hausse.}
\end{popfigure}

\textbf{Identification du concept.} La notion en jeu est la \cle{mort}, mais sous un angle paradoxal : la mort \emph{banalisée}. Nous savons que la route tue, et pourtant nous nous comportons comme si elle ne tuait pas.

\textbf{Mobilisation d'une thèse.} Heidegger a montré que l'homme ordinaire fuit la pensée de sa propre mort : « on meurt », dit-on, comme si cela ne concernait que les autres. Cette fuite explique notre insouciance : le casque « gêne », la ceinture « froisse la chemise », la prudence « fait perdre du temps ». Notre légèreté révèle une vérité philosophique profonde : nous vivons en niant notre finitude. Or, dit encore Heidegger, c'est en assumant que nous sommes des « êtres-pour-la-mort » que nous donnons du prix à chaque instant de notre existence.

\begin{exemplebox}[Un chiffre qui doit faire réfléchir]
Avec plus de \num{1000} morts par an, les accidents de la route tuent au Cameroun environ \num{20} personnes par semaine, soit l'équivalent de plusieurs classes entières de Terminale qui disparaîtraient chaque semaine de l'année. Si une épidémie faisait autant de victimes, elle serait déclarée urgence nationale. La mort routière, elle, est devenue un bruit de fond : c'est précisément cette banalisation que la philosophie nous invite à interroger.
\end{exemplebox}

\textbf{Conclusion justifiée.} Notre insouciance sur la route n'est pas un simple manque de discipline : elle est le symptôme d'un rapport refoulé à la mort. Penser la mort — loin d'être morbide — est un exercice de sagesse qui rend la vie plus précieuse et la prudence plus intelligible. Respecter le code de la route, c'est philosophiquement reconnaître la valeur de l'existence, la sienne et celle des autres.

\courssub{Situation 4 : l'euthanasie et les soins palliatifs — a-t-on le droit de choisir sa mort ?}

\textbf{Description du fait.} Un malade atteint d'un cancer en phase terminale souffre atrocement ; la médecine ne peut plus le guérir. Dans certains pays (Belgique, Pays-Bas, Canada, Espagne), la loi autorise, sous conditions strictes, l'euthanasie : un médecin peut, à la demande répétée et éclairée du patient, provoquer sa mort pour abréger ses souffrances. Dans d'autres pays, dont le Cameroun, seuls les soins palliatifs — traitement de la douleur et accompagnement — sont licites, toute provocation active de la mort restant interdite.

\textbf{Identification du concept.} Deux notions se croisent : la \cle{mort} (sa signification, sa maîtrise) et l'\cle{existence} (la liberté de disposer de soi, la dignité).

\textbf{Mobilisation des thèses.} Le débat oppose deux arguments forts. \emph{Pour} le droit de choisir sa mort : l'argument de l'autonomie (ma vie m'appartient, nul ne peut m'obliger à souffrir) et l'argument de la compassion (laisser souffrir sans espoir est une cruauté). \emph{Contre} : l'argument de la sacralité de la vie (la vie humaine a une valeur absolue qui ne dépend pas de sa qualité) et l'argument de la pente glissante (autoriser la mort choisie risque de faire pression sur les personnes âgées, pauvres ou handicapées, qui se sentiraient « de trop »). Les soins palliatifs proposent une voie médiane : soulager la douleur sans donner la mort, accompagner sans abandonner.

\textbf{Conclusion justifiée.} Il n'existe pas de réponse simple, et une bonne copie de Bac doit le reconnaître. La position la plus solide tient ensemble deux exigences : le respect de la liberté du malade (on ne peut pas condamner quelqu'un à souffrir au nom d'un principe qu'il ne partage pas) et la protection des plus vulnérables (aucune société ne doit envoyer à ses faibles le message que leur vie ne vaut plus la peine). C'est pourquoi le développement des soins palliatifs, encore insuffisants en Afrique, est une priorité éthique : beaucoup de demandes d'euthanasie naissent moins du désir de mourir que de la peur de souffrir seul.

\courssub{Situation 5 : le suicide chez les jeunes — la question de Camus}

\textbf{Description du fait.} Partout dans le monde, y compris au Cameroun, des jeunes — parfois des élèves après un échec scolaire, une rupture amoureuse, une humiliation ou une détresse familiale — mettent fin à leurs jours. Albert Camus ouvre son essai \emph{Le Mythe de Sisyphe} par une phrase célèbre : « Il n'y a qu'un problème philosophique vraiment sérieux : c'est le suicide. » Autrement dit : la première question que la philosophie doit trancher est de savoir si la vie vaut ou non la peine d'être vécue.

\textbf{Identification du concept.} C'est la question du \cle{sens de l'existence} dans sa forme la plus aiguë : face à l'absurdité — le divorce entre notre soif de sens et le silence du monde — faut-il renoncer ou persévérer ?

\textbf{Mobilisation d'une thèse.} La réponse de Camus est ferme : il faut \emph{imaginer Sisyphe heureux}. Le suicide est un refus du problème, non une solution : il supprime à la fois l'absurdité et la conscience qui pouvait la regarder en face. La révolte lucide — vivre sans espoir illusoire mais sans renoncement — est la seule attitude digne. Les traditions religieuses ajoutent que la vie est un don dont nous ne sommes pas propriétaires mais dépositaires. Les psychologues, enfin, rappellent un fait décisif : la grande majorité des personnes suicidaires ne veulent pas mourir, elles veulent que la souffrance cesse — et la souffrance, contrairement à la vie, peut cesser autrement.

\textbf{Conclusion justifiée et prévention.} La discussion philosophique du suicide doit donc toujours être \emph{encadrée} : philosopher sur la mort ne signifie jamais la banaliser. Si toi-même ou un camarade traverse une détresse profonde, la démarche juste n'est pas l'isolement mais la parole : en parler à un proche, à un enseignant, à un conseiller d'orientation, à un personnel de santé. Demander de l'aide n'est pas une faiblesse ; c'est exactement ce que Camus appelle la révolte : refuser de laisser le silence avoir le dernier mot.

\courssub{Méthode : appliquer une notion à une situation concrète}

Les cinq situations précédentes suivent toutes la même démarche. Formalisons-la, car elle est directement mobilisable à l'examen, dans les sujets du type « À partir d'une situation de la vie courante, montrez que... ».

\begin{methode}[Les quatre étapes de l'application philosophique]
\begin{enumerate}
\item \textbf{Décrire le fait} : rapporter la situation brièvement et précisément (qui ? quoi ? où ?), sans jugement hâtif. Deux ou trois phrases suffisent.
\item \textbf{Identifier le concept} : nommer la notion philosophique en jeu (existence, mort, liberté, sens...) et formuler le problème qu'elle soulève sous forme de question.
\item \textbf{Mobiliser une thèse} : convoquer un auteur ou une doctrine étudiés (Sartre, Camus, Heidegger, les traditions religieuses...) et montrer comment sa thèse éclaire la situation.
\item \textbf{Tirer une conclusion justifiée} : prendre position en argumentant, puis \emph{nuancer} en reconnaissant la part de vérité de la thèse adverse. Une conclusion sans nuance est une conclusion fragile.
\end{enumerate}
\end{methode}

\begin{popfigure}
\begin{center}
\begin{tikzpicture}[
    node distance=0.55cm,
    etape/.style={rectangle, rounded corners, draw=popdark, fill=popblueL, text width=3.1cm, align=center, font=\small, minimum height=1cm},
    fleche/.style={-{Stealth[length=3mm]}, thick, popdark}
]
\node[etape, fill=popgoldL, align=center] (fait){\textbf{1. Fait}\\ Situation concrète décrite};
\node[etape, right=of fait, align=center] (notion){\textbf{2. Notion}\\ Concept identifié, problème posé};
\node[etape, right=of notion, align=center] (these){\textbf{3. Thèse}\\ Auteur mobilisé, argument};
\node[etape, right=of these, align=center] (position){\textbf{4. Position}\\ Conclusion justifiée et nuancée};
\draw[fleche] (fait) -- (notion);
\draw[fleche] (notion) -- (these);
\draw[fleche] (these) -- (position);
\node[below=0.35cm of notion, font=\small\itshape, text=popmuted, text width=5.5cm, align=center] {Ex. : funérailles $\rightarrow$ mort comme événement symbolique $\rightarrow$ le rite donne du sens (Ariès) $\rightarrow$ les funérailles consolident les vivants, à condition de ne pas devenir un faste ruinant};
\end{tikzpicture}
\end{center}
\legende{Arbre de décision philosophique : la démarche d'application d'une notion à une situation concrète, en quatre étapes enchaînées, illustrée par l'exemple des funérailles.}
\end{popfigure}

\begin{attention}[L'erreur qui coûte des points : raconter au lieu d'analyser]
L'erreur la plus fréquente consiste à \emph{raconter} la situation pendant un paragraphe entier — « au village, quand quelqu'un meurt, on tue des chèvres, on danse, les gens viennent de loin... » — sans jamais l'analyser philosophiquement. Or l'exemple doit \emph{servir} l'argument, non le remplacer. Règle d'or : la description du fait ne doit jamais dépasser un quart du développement ; les trois quarts restants appartiennent au concept, à la thèse et à l'argumentation. Un correcteur ne note pas ce que tu as vécu ; il note ce que tu en penses.
\end{attention}

\begin{exoresolu}[Appliquer la méthode à une situation nouvelle]
Énoncé. À Bafoussam, un jeune menuisier de \num{24} ans, orphelin, refuse de vendre l'atelier que son père défunt lui a laissé, malgré une offre d'achat très avantageuse : « Cet atelier, c'est mon père qui continue de vivre à travers moi. » Appliquez la démarche en quatre étapes pour analyser philosophiquement cette situation.
\tcblower
\textbf{Solution détaillée.}
\emph{Étape 1 — Le fait.} Un jeune artisan renonce à un gain financier important pour conserver l'atelier hérité de son père, au nom d'un lien symbolique avec le défunt.
\emph{Étape 2 — La notion.} La mort est ici en jeu, non comme disparition pure, mais comme rapport continu entre les vivants et les morts : le problème posé est celui de la survie symbolique — en quel sens un mort peut-il encore « être là » ?
\emph{Étape 3 — La thèse.} La conception traditionnelle africaine de la mort (que l'on retrouve chez des penseurs comme John Mbiti) soutient que le défunt n'est pas un « néant » mais un « vivant-mort » : il demeure membre de la communauté tant que les vivants entretiennent son souvenir et poursuivent son œuvre. L'atelier n'est donc pas un simple bien économique ; il est le support matériel d'une présence.
\emph{Étape 4 — La conclusion justifiée.} Le choix du jeune menuisier est philosophiquement défendable : il donne à la mort de son père un sens (la transmission) et à sa propre existence un projet (faire fructifier un héritage). On nuancera cependant : la fidélité au défunt ne doit pas devenir une prison ; si l'atelier condamne le jeune à la misère, honorer son père pourrait au contraire consister à transformer l'héritage, voire à le vendre pour financer un projet plus vaste. La mémoire authentique est créatrice, non muséale.
\end{exoresolu}

\begin{exercice}[À toi de jouer]
Choisis l'une des deux situations suivantes et rédige une analyse complète en quatre étapes (environ une page) :
\begin{enumerate}
\item À Maroua, une famille refuse qu'on annonce à leur mère mourante la gravité de sa maladie « pour ne pas lui faire de peine ». A-t-on le droit de mentir à quelqu'un sur sa propre mort ?
\item Un lycée technique de Douala organise chaque année une journée de sensibilisation à la sécurité routière après le décès d'un élève dans un accident de moto. Pourquoi faut-il parfois une mort pour que nous pensions la mort ?
\end{enumerate}
\end{exercice}

\begin{corrige}[Exercice — éléments de correction]
\textbf{Situation 1.} Fait : un mensonge par compassion sur la mort imminente. Notion : la mort comme vérité à dire ou à taire ; croisement avec la dignité et l'autonomie. Thèses mobilisables : Kant (le mensonge est toujours condamnable car il détruit la confiance et traite l'autre comme un moyen) ; contre Kant, l'éthique du care (la vérité brutale peut être une violence ; dire la vérité suppose de choisir le moment, les mots, la dose). Conclusion nuancée : le malade a un droit fondamental à la vérité sur sa propre mort — nul ne peut mourir à sa place, donc nul ne peut décider à sa place de ce qu'il doit savoir — mais ce droit s'exerce dans un dialogue progressif et respectueux, non dans une annonce brutale ni dans un mensonge total.
\textbf{Situation 2.} Fait : la sensibilisation n'intervient qu'après un drame. Notion : la mort banalisée et la fuite devant la finitude. Thèse : Heidegger — « on meurt » concerne toujours les autres jusqu'au jour où la mort frappe proche ; la mort d'un camarade brise l'anonymat du « on » et rend chacun à sa propre mortalité. Conclusion : la journée de sensibilisation est philosophiquement un exercice de lucidité ; mais une sagesse véritable consisterait à penser la mort \emph{avant} le drame, car attendre le deuil pour devenir prudent, c'est laisser la mort faire elle-même notre éducation.
\end{corrige}

\begin{aretenir}[Ce qu'il faut retenir de cette section]
\begin{itemize}
\item Appliquer une notion, c'est suivre quatre étapes : \cle{fait} $\rightarrow$ \cle{notion} $\rightarrow$ \cle{thèse} $\rightarrow$ \cle{position justifiée et nuancée}.
\item Les funérailles camerounaises montrent que la mort humaine est un événement symbolique : le rite donne du sens aux morts et de la solidarité aux vivants.
\item Face au chômage, l'existence se pense comme \emph{projet situé} : ni fatalisme pur, ni volontarisme naïf.
\item La banalisation de la mort routière révèle notre fuite devant la finitude (Heidegger) ; penser la mort rend la vie plus précieuse.
\item L'euthanasie oppose l'autonomie à la sacralité de la vie ; les soins palliatifs dessinent une voie médiane.
\item Camus fait du suicide le problème philosophique par excellence, pour mieux le réfuter : la révolte lucide, et non le renoncement, est la réponse à l'absurdité — et la parole est le premier geste de prévention.
\end{itemize}
\end{aretenir}

\coursec{Introduction : problématiser l'existence et la mort}

\begin{experience}[Situation déclenchante : le deuil au village]
\begin{center}
\begin{tikzpicture}[
    node distance=1.5cm and 2cm,
    event/.style={draw=popblue, thick, rounded corners, fill=popblue!10, text width=3.5cm, align=center, minimum height=1.5cm, font=\small},
    arrow/.style={-Stealth, thick, popdark},
    quest/.style={draw=poporange, thick, rounded corners, fill=poporange!10, text width=4cm, align=center, minimum height=1.5cm, font=\small, font=\itshape},
]
\node[event, align=center] (deces){Décès d'un aîné\\dans un village Bamiléké};
\node[event, right=of deces] (rites) {Rites funèbres (\emph{kùm}) :\newline veillées, chants, partage,\newline discours sur l'héritage};
\node[event, right=of rites] (parole) {Parole du notable :\newline \emph{« La mort n'est pas la fin,\newline c'est le retour aux ancêtres »}};
\node[quest, align=center, below=1.5cm of rites] (q1) {La mort anéantit-elle\\ou révèle-t-elle l'existence ?};
\node[quest, align=center, below=1.5cm of q1] (q2) {L'homme a-t-il un sens\\donné ou à créer ?};
\draw[arrow] (deces) -- (rites);
\draw[arrow] (rites) -- (parole);
\draw[arrow, poporange] (parole.south) -- (q1.north);
\draw[arrow, poporange] (q1.south) -- (q2.north);
\end{tikzpicture}
\end{center}
\legende{Le deuil traditionnel : un laboratoire concret où s'éprouvent les concepts d'existence, de mort et de sens.}
\end{experience}

\begin{aretenir}[Problématique directrice de la leçon]
\textbf{Comment la conscience de la mort, loin de nier l'existence, peut-elle en fonder le sens ?} \\
Cette question traverse trois tensions majeures :
\begin{itemize}
    \item \textbf{Finitude vs Néant :} La mort est-elle une \cle{frontière} qui donne sa forme à la vie (Heidegger) ou un \cle{néant} qui la rend absurde (Camus, Épicure) ?
    \item \textbf{Subir vs Assumer :} L'homme est-il \emph{jeté} vers la mort (passivité) ou \emph{projet} vers sa fin (liberté, authenticité) ?
    \item \textbf{Individuel vs Communautaire :} Le sens se construit-il en solitaire (existentialisme) ou dans la \cle{filiation} et la transmission (pensée africaine, Frankl) ?
\end{itemize}
\end{aretenir}

\coursec{Définition du concept d'existence}

\begin{definition}[Existence (étymologie et sens philosophique)]
Du latin \emph{ex-sistere} : « se tenir hors de », « se démarquer », « apparaître ».
\begin{itemize}
    \item \textbf{Sens ontologique :} Manière d'être propre à l'homme : ne pas \emph{être} simplement (comme une pierre), mais \emph{exister}, c'est-à-dire se projeter, choisir, se faire. \cle{L'existence précède l'essence} (Sartre).
    \item \textbf{Sens existentiel (Heidegger) :} Le \cle{Dasein} (l'être-là) est cet être pour qui son être \emph{importe}. Caractères : \cle{jetéité} (né dans un monde sans l'avoir choisi), \cle{projet} (se projeter vers des possibles), \cle{compréhension} (s'entendre sur ses possibilités).
    \item \textbf{Distinction cruciale :} \textbf{Vie biologique} (métabolisme, reproduction, durée) $\neq$ \textbf{Existence humaine} (conscience, liberté, historicité, projet, responsabilité).
\end{itemize}
\end{definition}

\begin{exemplebox}[Existence technique : l'ingénieur face au chantier]
Un ingénieur des travaux publics à Douala ne se contente pas d'appliquer des normes (essence technique). Il \textbf{existe} en : (1) anticipant les risques d'érosion (projet), (2) choisissant un matériau local (responsabilité/liberté), (3) transmettant son savoir aux jeunes (historicité/filiation). Son existence \emph{est} cette tension vers un avenir qu'il façonne.
\end{exemplebox}

\begin{attention}[Piège : confondre existence et vie]
Dire « il a une belle existence » pour « il a une vie confortable » est une erreur. L'existence suppose \cle{conscience} et \cle{liberté}. Un végétatif a une vie biologique, pas une existence au sens philosophique.
\end{attention}

\coursec{Définition du concept de mort}

\begin{definition}[Mort : définitions croisées]
\begin{center}
\begin{tabularx}{\linewidth}{|l|X|X|}\hline
\textbf{Approche} & \textbf{Définition} & \textbf{Enjeu philosophique} \\ \hline
\textbf{Biologique} & Cessation irréversible des fonctions vitales (cérébrales, cardiaques, respiratoires). & Fait naturel, universel, objectivable. \\ \hline
\textbf{Phénoménologique (Heidegger)} & \cle{Fin} de la possibilité d'être-là. \emph{« La mort, comme fin du Dasein, est la possibilité propre, non-relative, certaine, indépassable, insurmontable. »} (\emph{Être et Temps}, §46-53). & La mort n'est pas un événement extérieur, elle \emph{structure} l'existence (être-pour-la-mort). \\ \hline
\textbf{Absurde (Camus)} & Le \cle{mur} infranchissable qui rend tout projet vain. \emph{« La mort est le seul problème philosophique vraiment sérieux. »} (\emph{Le Mythe de Sisyphe}). & Révèle le divorce entre l'appel au sens et le silence du monde. \\ \hline
\textbf{Épicurienne} & \emph{« La mort n'est rien pour nous : tant que nous existons, la mort n'est pas là ; quand la mort est là, nous n'existons plus. »} (\emph{Lettre à Ménécée}). & Libère de la peur : le néant ne se vit pas. \\ \hline
\textbf{Anthropologie africaine (Mbiti, sagesse orale)} & \cle{Passage} : séparation du corps, continuation de l'esprit (ancêtre). \emph{« La mort n'est pas la fin, c'est un changement de demeure. »} & La mort \emph{intègre} l'individu à la communauté des vivants et des morts (communion des ancêtres). \\ \hline
\end{tabularx}
\end{center}
\end{definition}

\begin{propriete}[Distinctions opérantes pour le Bac]
\begin{itemize}
    \item \textbf{Mourir (verbe) vs La mort (nom) :} Mourir est un processus (agonie) ; la mort est l'événement limite (aboutissement).
    \item \textbf{Ma mort vs La mort d'autrui :} \cle{Ma mort} (propre, non déléguable, Heidegger) $\neq$ \cle{La mort d'autrui} (deuil, témoignage, responsabilité éthique — Lévinas).
    \item \textbf{Finitude vs Néant :} Finitude = limite constitutive (je suis fini) ; Néant = absence de sens (le néant après la mort).
\end{itemize}
\end{propriete}

\coursec{Le sens de l'existence humaine : les grandes conceptions}

\begin{popfigure}
\centering
\begin{tikzpicture}[
    mindmap,
    concept color=popblue!50,
    level 1/.append style={level distance=4cm, sibling angle=90},
    level 2/.append style={level distance=2.8cm, sibling angle=40},
    grow cyclic,
]
\node[concept] {Sens de l'existence}
    child[concept color=popgreen!60] { node[concept, align=center]{Sens \textbf{donné}\\(Essentialisme, Religion)} 
        child { node[concept, scale=0.7] {Dieu / Ordre cosmique} }
        child { node[concept, scale=0.7] {Nature humaine fixe} }
        child { node[concept, scale=0.7] {Ex: Thomisme, Islam, Confucianisme} }
    }
    child[concept color=poporange!70] { node[concept, align=center]{Sens \textbf{absent}\\(Nihilisme, Absurde)} 
        child { node[concept, scale=0.7] {Néant final (Camus)} }
        child { node[concept, scale=0.7] {Vanité des projets} }
        child { node[concept, scale=0.7] {Ex: Cioran, Nietzsche (passif)} }
    }
    child[concept color=poppurple!70] { node[concept, align=center]{Sens \textbf{créé}\\(Existentialisme, Logothérapie)} 
        child { node[concept, scale=0.7] {Liberté radicale (Sartre)} }
        child { node[concept, scale=0.7] {Révolte / Témoignage (Camus, Frankl)} }
        child { node[concept, scale=0.7] {Ex: Engagement, Œuvre, Amour} }
    }
    child[concept color=poppink!70] { node[concept, align=center]{Sens \textbf{communautaire}\\(Pensée africaine)} 
        child { node[concept, scale=0.7] {Filiation / Ancêtres (Mbiti)} }
        child { node[concept, scale=0.7] {Ubuntu : \emph{Je suis parce que nous sommes}} }
        child { node[concept, scale=0.7] {Rites de transmission} }
    };
\end{tikzpicture}
\legende{Cartographie des quatre grandes réponses philosophiques au sens de l'existence face à la mort.}
\end{popfigure}

\courssub{Le sens donné : Essentialisme et Religions}
\begin{itemize}
    \item \textbf{Thèse :} L'homme a une \cle{nature} et une \cle{fin} (telos) prédéterminées. Le sens est \emph{objectif}, inscrit dans l'ordre du monde ou la volonté divine.
    \item \textbf{Auteurs :} Aristote (\emph{Éthique à Nicomaque} : bonheur = activité vertueuse selon la raison), Saint Thomas d'Aquin (fin surnaturelle : vision béatifique), Coran (Sourate 51:56 « Je n'ai créé les djinns et les hommes que pour qu'ils M'adorent »).
    \item \textbf{Critique (Sartre) :} « L'existence précède l'essence » : pas de nature humaine fixe, pas de Dieu pour la garantir $\rightarrow$ l'homme est \emph{condamné à être libre}.
\end{itemize}

\courssub{Le sens absent : L'absurde et le nihilisme}
\begin{itemize}
    \item \textbf{Thèse :} La mort anéantit tout. L'univers est indifférent. La quête de sens heurte le « silence irraisonné du monde » (Camus).
    \item \textbf{Auteurs :} \textbf{Camus} (\emph{Le Mythe de Sisyphe}) : l'absurde = divorce homme/monde. \textbf{Épicure} : la mort = néant, pas de malheur possible. \textbf{Nietzsche} (nihilisme passif) : « Dieu est mort », les valeurs suprêmes se dévalorisent.
    \item \textbf{Issue camusienne :} Ne pas fuir (suicide physique ou philosophique/religion), mais \cle{vivre l'absurde} : révolte, liberté, passion (quantité d'expériences).
\end{itemize}

\courssub{Le sens créé : Existentialisme et Logothérapie}
\begin{itemize}
    \item \textbf{Sartre} (\emph{L'Être et le Néant}, \emph{L'Existentialisme est un humanisme}) : \cle{Liberté radicale}. Pas de sens a priori. Le sens naît de l'\cle{engagement} dans un projet libre. « L'homme n'est rien d'autre que ce qu'il se fait. »
    \item \textbf{Heidegger} (\emph{Être et Temps}) : \cle{Authenticité} = assumer son \emph{être-pour-la-mort} (anticipation résolue de sa fin propre). Cela arrache au « On » (inauthenticité) et donne du poids à chaque choix.
    \item \textbf{Viktor Frankl} (\emph{Un sens à la vie} / Logothérapie) : \cle{Volonté de sens} (primary motivation). Trois voies : (1) Créer une œuvre / accomplir un travail, (2) Vivre une rencontre (amour), (3) Transformer une souffrance en témoignage (ex: camps de concentration). \emph{« Celui qui a un pourquoi vivre peut supporter presque n'importe quel comment. »}
\end{itemize}

\courssub{Le sens communautaire : Anthropologie philosophique africaine}
\begin{itemize}
    \item \textbf{John Mbiti} (\emph{African Religions and Philosophy}) : La mort est un \cle{rite de passage} vers l'état d'\emph{ancêtre} (\emph{living-dead}). L'individu ne meurt vraiment que quand il est oublié.
    \item \textbf{Placide Tempels} (\emph{La Philosophie bantoue}) : \cle{Force vitale}. L'existence = participation à la force vitale du clan. La mort affaiblit le clan, les rites le renforcent.
    \item \textbf{Achille Mbembe} (\emph{Critique de la raison nègre}, \emph{Nécropolitique}) : Analyse moderne du pouvoir sur la mort (colonialisme, biopolitique), mais rappelle la centralité de la \cle{filiation} dans la constitution du sujet africain.
    \item \textbf{Synthèse :} Le sens n'est pas \emph{subjectif} (mon choix seul) ni \emph{objectif} (donné par Dieu), mais \textbf{intersubjectif} : je reçois un héritage (nom, langue, terre) et je le transmets. \emph{« Je suis parce que nous sommes. »} (Ubuntu).
\end{itemize}

\begin{aretenir}[Synthèse comparative pour la dissertation]
\begin{center}
\begin{tabularx}{\linewidth}{|l|X|X|X|}\hline
\textbf{Conception} & \textbf{Rapport à la mort} & \textbf{Source du sens} & \textbf{Figure clé} \\ \hline
Essentialisme / Religion & Passage / Accomplissement & Donné (Dieu, Nature) & Aristote, St Thomas \\ \hline
Absurde (Camus) & Mur / Néant final & \textbf{Néant} (à assumer) & Camus, Épicure \\ \hline
Existentialisme (Sartre/Heidegger) & Échéance / Structure & \textbf{Créé} (Liberté, Projet) & Sartre, Heidegger \\ \hline
Logothérapie (Frankl) & Défi / Témoignage & \textbf{Découvert} (Valeurs) & Frankl \\ \hline
Pensée africaine (Mbiti) & Passage / Ancestralité & \textbf{Transmis/Partagé} (Clan) & Mbiti, Sagesse orale \\ \hline
\end{tabularx}
\end{center}
\end{aretenir}

\coursec{Application aux situations concrètes de la vie courante}

\begin{methode}[Analyser une situation concrète : Grille « Existence - Mort - Sens »]
Pour toute situation (deuil, choix de carrière, engagement, fin de vie), appliquez :
\begin{enumerate}
    \item \textbf{Identifier l'expérience vécue :} Quel événement ? (Décès, diagnostic, choix d'orientation, manifestation).
    \item \textbf{Repérer les concepts en jeu :} Existence (projet, liberté, angoisse) ? Mort (finitude, néant, passage, angoisse) ? Sens (valeur, finalité, absurdité, transmission) ?
    \item \textbf{Mobiliser une thèse d'auteur :} Quel éclairage ? (Heidegger : authenticité ; Camus : révolte ; Frankl : témoignage ; Mbiti : filiation).
    \item \textbf{Conclure sur la portée :} La situation révèle-t-elle une fuite (inauthenticité), une révolte, un engagement, une transmission ?
\end{enumerate}
\end{methode}

\begin{exemplebox}[Cas 1 : Le deuil au village (Bamiléké / Ouest Cameroun)]
\begin{itemize}
    \item \textbf{Situation :} Décès du patriarche. « Kùm » (deuil d'un an) : port du noir, interdits, veillées, \emph{tsu} (cérémonie de levée de deuil), partage de l'héritage, intronisation de l'héritier.
    \item \textbf{Analyse :} La mort n'efface pas : elle \textbf{révèle} l'existence du défunt (bilan) et \textbf{restructure} l'existence des vivants (nouveaux rôles). Le sens = \cle{continuité lignagère} (Mbiti). L'angoisse est \cle{ritualisée}, partagée, non solitaire.
    \item \textbf{Thèse mobilisée :} Heidegger (mort comme possibilité propre qui libère pour la transmission) + Mbiti (mort comme passage ancestral).
\end{itemize}
\end{exemplebox}

\begin{exemplebox}[Cas 2 : Choix d'un métier à Douala (Jeune diplômé)]
\begin{itemize}
    \item \textbf{Situation :} Ingénieur informatique, offre bien payée dans une banque vs startup à impact social (agri-tech) mal payée.
    \item \textbf{Analyse :} \textbf{Sartre} : Choix = engagement de tout son être. Pas de « nature » d'ingénieur banquier. \textbf{Frankl} : Sens par l'œuvre (créer) ou le témoignage (servir). \textbf{Heidegger} : Fuite dans le « On » (salaire, sécurité, regard social) vs Authenticité (assumer le risque pour un projet qui me définit).
    \item \textbf{Conclusion :} Le salaire assure la \emph{vie biologique} ; le choix du sens assure l'\emph{existence}.
\end{itemize}
\end{exemplebox}

\begin{exemplebox}[Cas 3 : Engagement associatif (Rotaract / Lions Club / ONG à Yaoundé)]
\begin{itemize}
    \item \textbf{Situation :} Jeunes professionnels organisent une campagne de don de sang à l'Hôpital Central.
    \item \textbf{Analyse :} Face à la \cle{finitude} (maladie, accidents), l'action crée du \cle{sens} par la \cle{solidarité}. \textbf{Camus} : Révolte contre l'injustice/la souffrance = « Je me révolte, donc nous sommes ». \textbf{Ricoeur} (\emph{Soi-même comme un autre}) : L'éthique = visée de la vie bonne \emph{avec et pour autrui} en institutions justes.
    \item \textbf{Portée :} L'engagement transforme l'angoisse de mort en \cle{puissance d'agir} collective.
\end{itemize}
\end{exemplebox}

\begin{experience}[Cas 4 : Fin de vie médicale — Débat bioéthique au Cameroun]
\begin{itemize}
    \item \textbf{Contexte :} Service de réanimation, Hôpital Général de Yaoundé. Patient en état végétatif persistant / phase terminale cancer. Famille divisée : « Tout faire » vs « Laisser partir dignement ».
    \item \textbf{Concepts :} \cle{Acharnement thérapeutique} (prolonger la vie biologique sans espoir d'existence) vs \cle{Soins palliatifs} (soulager, accompagner l'existence jusqu'au bout). \cle{Euthanasie} (interdite au Cameroun, Code pénal Art. 275-276) vs \cle{Sédation profonde} (légale, fin de vie).
    \item \textbf{Éclairages :} \textbf{Heidegger} : La mort doit rester \emph{propre} (non déléguée à la technique). \textbf{Frankl} : Dernier sens possible = \emph{attitude} face à la souffrance inévitable. \textbf{Pensée africaine} : La mort « à l'heure » (ni trop tôt, ni trop tard par machine) permet le \emph{dernier acte de sens} (bénédictions, pardon, transmission orale).
    \item \textbf{Question éthique :} La technique médicale (respirateur, réanimation) \emph{vole-t-elle} la mort propre au patient ?
\end{itemize}
\end{experience}

\coursec{Méthode : rédiger et justifier une démarche argumentative}

Après avoir exploré comment les notions d'existence et de mort éclairent des situations concrètes — le deuil au village, le choix d'un métier, l'engagement associatif à Douala ou Yaoundé, la fin de vie à l'hôpital — nous devons maintenant transformer cette compréhension en une compétence examinable : \textbf{savoir construire une réponse philosophique rigoureuse, structurée et argumentée}. Au Baccalauréat technique, la dissertation n'est pas un récit d'opinions, mais une \cle{démonstration} où chaque thèse est étayée par une référence précise, un exemple analysé et un raisonnement logique. Cette section vous donne la méthode pas à pas, les critères d'évaluation officiels et un entraînement corrigé pour éviter les pièges classiques.

\courssub{Structure type de la dissertation au Baccalauréat technique}

\begin{popfigure}
\centering
\begin{tikzpicture}[
    node distance=1.2cm and 1.5cm,
    box/.style={draw=popblue, thick, rounded corners, fill=popblue!5, text width=3.2cm, align=center, minimum height=1.8cm, font=\small},
    arrow/.style={-Stealth, thick, popdark},
    funnel/.style={draw=poporange, thick, fill=poporange!10, rounded corners, minimum height=1cm, text width=12cm, align=center, font=\small\bfseries},
]
% Introduction (Funnel)
\node[funnel] (intro) {INTRODUCTION (entonnoir)};
\node[box, below=0.8cm of intro.south west, anchor=north west, xshift=0.5cm, align=center] (accroche){\textbf{Accroche concrète}\\\emph{Exemple camerounais, citation, fait d'actualité, paradoxe}};
\node[box, right=0.5cm of accroche, align=center] (analyse){\textbf{Analyse des termes}\\\emph{Définir « conscience », « mort », « sens », « vie »}};
\node[box, right=0.5cm of analyse, align=center] (problematique){\textbf{Problématique}\\\emph{Question philosophique unique, ouverte, issue de la tension}};
\node[box, right=0.5cm of problematique, align=center] (plan){\textbf{Annonce du plan}\\\emph{« Nous examinerons d'abord... puis... enfin... »}};
\draw[arrow] (accroche.east) -- (analyse.west);
\draw[arrow] (analyse.east) -- (problematique.west);
\draw[arrow] (problematique.east) -- (plan.west);
% Développement
\node[funnel, below=2.5cm of intro.south, yshift=-0.5cm] (dev) {DÉVELOPPEMENT (3 parties équilibrées)};
\node[box, below=0.8cm of dev.south west, anchor=north west, xshift=0.5cm, fill=popgreen!5, draw=popgreen, align=center] (these){\textbf{Thèse 1}\\\emph{La conscience de la mort donne urgence et valeur à la vie (Heidegger, Sartre)}};
\node[box, right=0.5cm of these, fill=poppink!5, draw=poppink, align=center] (antithese){\textbf{Antithèse}\\\emph{La mort rend la vie absurde, le sens est illusoire (Camus, Épicure, Bouddhisme)}};
\node[box, right=0.5cm of antithese, fill=poppurple!5, draw=poppurple, align=center] (synthese){\textbf{Synthèse / Dépassement}\\\emph{Le sens ne vient pas de la mort mais de la liberté de créer des valeurs (Frankl, Ricoeur, traditions africaines)}};
\draw[arrow] (these.east) -- (antithese.west);
\draw[arrow] (antithese.east) -- (synthese.west);
% Conclusion
\node[funnel, below=2.5cm of dev.south, yshift=-0.5cm] (conc) {CONCLUSION (entonnoir inverse)};
\node[box, below=0.8cm of conc.south west, anchor=north west, xshift=0.5cm, align=center] (reponse){\textbf{Réponse explicite}\\\emph{« Oui, mais à condition de... »}};
\node[box, right=0.5cm of reponse, align=center] (resume){\textbf{Résumé de la démarche}\\\emph{Reprise synthétique des 3 temps}};
\node[box, right=0.5cm of resume, align=center] (ouverture){\textbf{Ouverture}\\\emph{Question nouvelle, domaine voisin (bioéthique, IA, rituels)}};
\draw[arrow] (reponse.east) -- (resume.west);
\draw[arrow] (resume.east) -- (ouverture.west);
% Vertical links
\draw[arrow, popblue] (intro.south) -- (dev.north);
\draw[arrow, popblue] (dev.south) -- (conc.north);
\end{tikzpicture}
\legende{Schéma de la structure dissertation : entonnoir introductionnel, développement dialectique (thèse/antithèse/synthèse), conclusion en entonnoir inverse.}
\end{popfigure}

\begin{methode}[Rédiger une introduction « entonnoir » en 4 étapes]
\textbf{Étape 1 — Accroche concrète (2--3 lignes).} Partez d'un fait observable au Cameroun : un rite funèbre \emph{Bamiléké} à Bafoussam (le « \emph{kùm} » : veillées d'un an, partage, intronisation), l'épitaphe sur une tombe au cimetière de la Madelaine à Yaoundé, le témoignage d'un malade en phase terminale à l'Hôpital Central, ou une parole de sage ewondo : « \emph{La mort n'est pas la fin, c'est le retour aux ancêtres} ».
\textbf{Étape 2 — Analyse des termes du sujet (3--4 lignes).} Définissez philosophiquement les concepts clés : \cle{conscience} (représentation réflexive, lucidité), \cle{mort} (fin biologique, néant ou passage), \cle{sens} (finalité, valeur, intelligibilité), \cle{vie} (existence biologique, existence humaine projetée). Montrez les tensions : la mort comme fin \emph{vs} la mort comme sens.
\textbf{Étape 3 — Problématique (1 phrase interrogative).} Elle doit naître de l'opposition entre les définitions. \textbf{Modèle :} « Si la mort est la fin absolue, comment la conscience que nous en avons peut-elle fonder un sens à la vie, au lieu de la détruire ? »
\textbf{Étape 4 — Annonce du plan (1 phrase).} « Nous analyserons d'abord comment la conscience de la finitude structure l'existence (Heidegger), puis pourquoi elle peut mener à l'absurde (Camus), pour enfin montrer que le sens émerge de la réponse libre que l'homme donne à sa mort (Frankl, pensée africaine). »
\end{methode}

\begin{attention}[Piège fréquent : l'accroche « hors-sujet » ou « catalogue »]
Ne commencez pas par : « Depuis l'Antiquité, les philosophes parlent de la mort... » (trop vague). Ne citez pas trois auteurs sans les lier au sujet. L'accroche doit \textbf{ancrer le sujet dans le réel camerounais ou dans une expérience universelle précise} pour montrer que la philosophie répond à une interrogation vécue.
\end{attention}

\courssub{Rédaction guidée d'une introduction : « La conscience de la mort donne-t-elle un sens à la vie ? »}

Voici une introduction modèle, annotée pour montrer la mécanique interne.

\begin{exemplebox}[Introduction type — Sujet : « La conscience de la mort donne-t-elle un sens à la vie ? »]
\textbf{[Accroche]} À Bafoussam, chez les Bamiléké, lorsqu'un aîné s'éteint, la communauté organise le « \emph{kùm} » : veillées d'un an, chants, partage des biens, discours sur « ce qu'il a laissé », intronisation de l'héritier. La mort n'efface pas l'homme ; elle révèle ce qu'il \emph{a été} et \emph{transmet}. \\
\textbf{[Analyse]} La \cle{conscience de la mort} est cette lucidité propre à l'homme : savoir qu'on mourra, non comme un animal qui meurt sans le savoir. Le \cle{sens} n'est pas une définition du dictionnaire, c'est une \cle{valeur} que l'existence acquiert quand on la regarde à l'horizon de sa fin. La \cle{vie} ici, c'est l'existence humaine : projet, relations, œuvre. \\
\textbf{[Problématique}] Dès lors, \textbf{la conscience de la mort est-elle la condition même du sens (urgence, authenticité), ou bien ne fait-elle qu'éclairer l'absurde d'une vie vouée au néant ?} \\
\textbf{[Plan]} Nous verrons d'abord que l'angoisse de la mort fonde l'authenticité (Heidegger, \emph{Être et Temps}), puis que cette conscience peut plonger dans l'absurde (Camus, \emph{Le Mythe de Sisyphe}), pour montrer enfin que le sens naît de la \cle{liberté} de répondre à la mort par un engagement (Frankl, \emph{Un sens à la vie} ; pensée africaine de la filiation).
\end{exemplebox}

\begin{aretenir}[Check-list introduction réussie]
\begin{itemize}
    \item Accroche \textbf{ancrée} (Cameroun, expérience vécue, œuvre d'art, fait précis).
    \item Termes \textbf{définis philosophiquement} (pas sens courant).
    \item Problématique \textbf{unique}, \textbf{ouverte}, issue d'une \textbf{tension} (pas une question de cours : « Qu'est-ce que la mort ? »).
    \item Plan \textbf{annoncé} avec les \textbf{grands axes} (pas « I, II, III » mais le fil logique).
\end{itemize}
\end{aretenir}

\courssub{Construire un paragraphe argumenté : la « brique » de base du développement}

Un paragraphe de dissertation n'est pas un résumé de cours. C'est une \cle{unité argumentative} : \textbf{Thèse + Référence exploitée + Exemple concret analysé + Lien logique vers le suivant}.

\begin{methode}[La méthode « T.R.E.A. » pour un paragraphe solide]
\textbf{T — Thèse (1 phrase) :} Affirmez l'idée directrice du paragraphe. \emph{Ex : « La conscience de la mort confère à chaque instant une valeur unique en nous arraignant à la dispersion. »}
\textbf{R — Référence (Auteur + Œuvre + Concept précis) :} Citez \textbf{exactement}. \emph{Ex : « Heidegger, dans \emph{Être et Temps} (§ 50-53), analyse l'\cle{être-pour-la-mort} : assumer sa finitude propre, non déléguable, libère le \emph{Dasein} de l'inauthenticité du « On » (la foule). »}
\textbf{E — Exemple camerounais analysé (3--4 lignes) :} Ne donnez pas l'exemple brut, \textbf{analysez-le} avec le concept. \emph{Ex : « Chez les Bassa, l'annonce d'une maladie grave pousse le chef de famille à réunir ses enfants pour « régler les comptes », transmettre les secrets de la lignée, pardonner. Ce n'est pas la mort qui donne du sens, c'est la \textbf{lucidité} face à elle qui transforme les derniers mois en \emph{testament vivant}, en acte de sens. »}
\textbf{A — Analyse / Transition (1 phrase) :} Montrez ce que cet exemple prouve et ouvrez sur la limite. \emph{Ex : « Ainsi, la finitude assumée structure le sens ; mais que se passe-t-il quand la mort apparaît comme \emph{absurde}, sans transmission possible ? »}
\end{methode}

\begin{exoresolu}[Paragraphe d'élève anonymisé — Correction collective]
\textbf{Sujet :} « La conscience de la mort donne-t-elle un sens à la vie ? »
\textbf{Partie I (Thèse) :} La conscience de la mort donne du sens.
\textbf{Paragraphe élève :}
\begin{quote}
\emph{« Heidegger dit que la mort donne du sens à la vie. Quand on sait qu'on va mourir, on profite mieux. Par exemple, mon oncle à Douala a eu un cancer. Il a voyagé, il a mangé ce qu'il voulait avant de mourir. Donc la mort donne du sens. Sartre aussi dit qu'on est libre. »}
\end{quote}
\textbf{Correction guidée :}
\begin{enumerate}
    \item \textbf{Référence floue :} « Heidegger dit que... » $\rightarrow$ \textbf{Précisez :} « Heidegger (\emph{Être et Temps}) distingue la \cle{mort comme fin} (biologique) de l'\cle{être-pour-la-mort} (existentiel). C'est l'\textbf{anticipation} de sa mort propre qui permet l'authenticité, pas le simple fait de savoir qu'on meurt. »
    \item \textbf{Exemple non analysé :} L'oncle « profite » (hédonisme) $\neq$ authenticité heideggérienne (responsabilité, transmission). $\rightarrow$ \textbf{Recadrez :} « L'oncle qui « profite » reste peut-être dans l'inauthenticité (fuite). L'authenticité serait : il assume sa finitude pour \textbf{choisir} ce qui compte vraiment (réconciliation, legs). »
    \item \textbf{Catalogage d'auteurs :} « Sartre aussi dit... » sans lien. $\rightarrow$ \textbf{Supprimez ou intégrez :} Si Sartre sert la thèse, montrez \textbf{comment} (liberté radicale face à la mort).
    \item \textbf{Transition absente :} Le paragraphe s'arrête net. $\rightarrow$ \textbf{Ajoutez :} « Mais cette authenticité suppose qu'on \textbf{puisse} choisir. Si la mort frappe l'enfant, le sens s'effondre-t-il ? »
\end{enumerate}
\textbf{Paragraphe réécrit (niveau Bac 16+/20) :}
\begin{quote}
\emph{« La conscience de la mort peut fonder le sens par l'authenticité qu'elle exige. Heidegger (\emph{Être et Temps}, §50) montre que l'\cle{être-pour-la-mort} — l'anticipation résolue de sa fin propre, non déléguable — arrache le \emph{Dasein} à la fuite dans le « On » (la dispersion quotidienne). Au Cameroun, un chef de famille Bassa atteint d'un cancer incurable ne « profite » pas passivement ; il convoque sa lignée pour transmettre les noms des ancêtres, désigner l'héritier, régler les dettes symboliques. C'est cette \textbf{lucidité assumée} qui transforme ses derniers mois en \emph{accomplissement}, non en simple survie. Toutefois, cette thèse suppose une mort « à temps », prévisible : que devient le sens face à la mort absurde, précoce, qui ne laisse pas le temps de l'anticipation ? »}
\end{quote}
\end{exoresolu}

\begin{attention}[Erreur fatale : « Citer sans exploiter »]
Écrire « Selon Sartre, l'existence précède l'essence » sans expliquer \textbf{ce que cela change pour le sujet} (ici : pas de sens préétabli, le sens se \emph{crée} dans l'action face à la mort) vaut \textbf{zéro point} en argumentation. Toute référence = \textbf{Thèse de l'auteur + Concept opératoire + Lien explicite avec votre argument}.
\end{attention}

\courssub{Justifier une conclusion : prendre position sans trahir le développement}

La conclusion n'est pas un résumé. C'est l'\textbf{acte de jugement} du candidat. Elle doit trois choses : \textbf{Répondre}, \textbf{Résumer}, \textbf{Ouvrir}.

\begin{methode}[Justifier sa conclusion en 3 temps]
\textbf{1. Réponse explicite à la problématique (1 phrase affirmative, nuancée).} Reprenez les termes de la problématique. \emph{Modèle :} « La conscience de la mort donne un sens à la vie, \textbf{non pas parce que la mort en contient un en soi, mais parce qu'elle oblige l'homme à en créer un par sa liberté}. »
\textbf{2. Résumé synthétique de la démarche (2--3 phrases).} Reprenez le fil \textbf{logique}, pas la liste des auteurs. \emph{Modèle :} « Nous avons vu que l'angoisse heideggérienne structure l'existence (I), que l'absurde camusien menace ce sens quand la mort paraît sans raison (II), et que la réponse victorieuse est l'engagement libre — qu'il soit existentiel (Frankl) ou communautaire (rites africains de transmission) (III). »
\textbf{3. Ouverture maîtrisée (1 phrase).} Une question nouvelle, un domaine voisin, sans lancer un nouveau débat. \emph{Modèles :} « Cette analyse invite à repenser les soins palliatifs au Cameroun : comment accompagner la \cle{fin de vie} pour qu'elle reste un acte de sens ? » \emph{ou} « Elle prolonge la réflexion sur le \cle{transhumanisme} : si la technique repousse la mort, l'existence perd-elle son urgence vitale ? »
\end{methode}

\begin{exemplebox}[Conclusion modèle pour le sujet]
\textbf{En définitive,} la conscience de la mort donne un sens à la vie, \textbf{à condition que l'homme ne la subisse pas passivement mais y réponde par un choix lucide}. Le parcours a montré que l'anticipation de la fin (Heidegger) structure l'authenticité, que le risque d'absurde (Camus) est réel si la mort apparaît comme néant pur, et que le sens surgit in fine de la \cle{liberté} de transformer la finitude en témoignage — comme le font les rites de transmission au Cameroun ou l'engagement dans une cause qui survit à soi. \textbf{Dès lors, le sens n'est pas \emph{donné} par la mort, il est \emph{arraché} à la mort par la vie.}
\textbf{Ouverture :} Cette tension éclaire le débat actuel sur la \cle{fin de vie médicale} au Cameroun : la sédation profonde ou l'euthanasie (interdite) modifient-elles la possibilité de ce « dernier acte de sens » ?
\end{exemplebox}

\courssub{Critères d'évaluation au Baccalauréat technique : la grille officielle}

\begin{popfigure}
\centering
\begin{tikzpicture}
\begin{axis}[
    width=\linewidth,
    height=10cm,
    xbar stacked,
    xmin=0, xmax=20,
    ytick={1,2,3,4},
    yticklabels={Expression \& français (4 pts), Argumentation \& logique (6 pts), Références \& notions (6 pts), Structure \& plan (4 pts)},
    legend style={at={(0.5,-0.25)}, anchor=north, legend columns=4, font=\small},
    bar width=12pt,
    nodes near coords={\pgfmathprintnumber\pgfplotspointmeta},
    nodes near coords align={horizontal},
    xlabel={Points sur 20},
    tick label style={font=\small},
    axis lines=left,
    enlarge y limits=0.3,
]
\addplot[popblue!70] coordinates {(4,1) (6,2) (6,3) (4,4)};
\addplot[popgreen!70] coordinates {(4,1) (6,2) (6,3) (4,4)};
\addplot[poporange!70] coordinates {(4,1) (6,2) (6,3) (4,4)};
\addplot[poppurple!70] coordinates {(4,1) (6,2) (6,3) (4,4)};
\legend{Niveau Insuffisant (0-8), Niveau Passable (9-11), Niveau Assez Bien (12-14), Niveau Bien/TB (15-20)}
\end{axis}
\end{tikzpicture}
\legende{Répartition indicative des 20 points au Bac Technique Philosophie (basé sur grilles MINESEC). La pondération varie légèrement selon les académies mais la hiérarchie reste : Argumentation > Notions/Références > Structure > Expression.}
\end{popfigure}

\begin{aretenir}[Ce que le correcteur attend concrètement]
\begin{itemize}
    \item \textbf{Notions (6 pts) :} Définitions exactes (existence $\neq$ vie biologique ; mort $\neq$ disparition), distinctions opérantes (finitude/néant, sens/signification, authenticité/inauthenticité).
    \item \textbf{Références (6 pts) :} \textbf{2 à 3 auteurs} bien maîtrisés (Heidegger, Camus, Frankl, Épicure, Sartre, Ricoeur, auteurs africains : Mbiti, sagesse orale) \textbf{valent mieux que 10 noms jetés}. Exigence : \textbf{Œuvre + Concept + Citation ou paraphrase précise}.
    \item \textbf{Argumentation (6 pts) :} Enchaînement logique (donc / mais / toutefois), exemples \textbf{analysés} (pas illustrés), thèse/antithèse/synthèse réelle (pas thèse/thèse/thèse).
    \item \textbf{Structure (4 pts) :} Introduction entonnoir, paragraphes T.R.E.A., transitions, conclusion justifiée.
    \item \textbf{Expression (4 pts) :} Français correct, vocabulaire philosophique (finitude, néant, absurdité, transcendance, immanence, projet, engagement), connecteurs logiques.
\end{itemize}
\end{aretenir}

\begin{savaistu}[Le savais-tu ? — Stratégie Bac Technique]
Au Baccalauréat technique, \textbf{une copie structurée, avec 2 ou 3 références exactes exploitées, des exemples camerounais analysés, et une problématique tenue du début à la fin, obtient 14--16/20}. Un « catalogue » de 10 auteurs mal compris, sans structure, sans exemples concrets, tombe à 8--10/20. \textbf{La qualité de l'exploitation prime sur la quantité des citations.} Travaillez vos 3 « auteurs-refuges » (ex : Heidegger, Camus, Frankl + un auteur africain) jusqu'à connaître par cœur : œuvre, concept-clé, citation courte, contre-argument possible.
\end{savaistu}

\courssub{Entraînement final : auto-évaluation avec la grille}

\begin{exercice}[Auto-évaluation — Sujet d'entraînement : « L'existence humaine est-elle absurde face à la mort ? »]
Rédigez \textbf{uniquement} :
\begin{enumerate}
    \item L'\textbf{introduction complète} (accroche camerounaise, analyse termes, problématique, plan).
    \item Le \textbf{premier paragraphe de la thèse} (T.R.E.A. : Thèse + Heidegger/Être-pour-la-mort + Exemple camerounais analysé + Transition).
    \item La \textbf{conclusion justifiée} (Réponse + Résumé démarche + Ouverture).
\end{enumerate}
\textit{Temps conseillé : 45 minutes. Puis auto-correz avec la grille ci-dessus.}
\end{exercice}

\begin{corrige}[Exercice n°1 — Éléments attendus pour l'auto-correction]
\textbf{Introduction :} Accroche sur les funérailles \emph{Ewondo} à Yaoundé (le « \emph{mekuk} » : veillées, chants, transmission du nom) $\rightarrow$ Analyse : absurdité (Camus) vs sens (Heidegger/Frankl) $\rightarrow$ Problématique : « Si la mort anéantit tout projet, l'existence n'est-elle que dérisoire, ou l'homme peut-il fonder un sens \emph{malgré} la mort ? » $\rightarrow$ Plan : I. L'absurde comme expérience limite (Camus) ; II. La révolte et la création de sens (Camus, Frankl) ; III. L'anthropologie africaine : la mort comme passage, non fin (Mbiti, rites).
\textbf{Paragraphe I (Thèse de l'absurde) :} Thèse : « La conscience de la mort révèle l'absurde de la condition humaine. » Référence : Camus, \emph{Le Mythe de Sisyphe} : le divorce entre l'appel de l'homme au sens et le « silence irraisonné du monde » ; la mort est le « point final » qui rend tout projet vain. Exemple : Un jeune ingénieur doualais, bâtisseur de ponts, meurt à 35 ans dans un accident. Ses projets, ses calculs, ses rêves s'arrêtent net. La conscience de cette possible brutalité rend son travail \emph{a posteriori} dérisoire \emph{si} on cherche un sens \emph{donné} par l'univers. Transition : « Mais Camus ne s'arrête pas à l'absurde : il propose la révolte. »
\textbf{Conclusion :} « L'existence n'est absurde face à la mort \textbf{que si l'on attend un sens extérieur}. Camus montre que la \cle{révolte} (vivre sans appel, multiplier les expériences), Frankl que le \cle{sens du sens} (témoignage, amour, humour), et les rites africains que la \cle{filiation} transforment l'absurde en \textbf{charge assumée}. Le sens n'est pas \emph{trouvé}, il est \emph{fait}. » Ouverture : « Cette lucidité change-t-elle notre rapport aux risques industriels ou routiers au Cameroun ? »
\end{corrige}

\begin{methode}[Rituel de révision « 3 x 20 min » avant l'épreuve]
\begin{enumerate}
    \item \textbf{20 min :} Sur un sujet zéro, écrire \textbf{uniquement l'introduction} (chronomètre). Vérifier : accroche ancrée ? termes définis ? problématique tension ? plan logique ?
    \item \textbf{20 min :} Écrire \textbf{un paragraphe T.R.E.A.} complet avec référence exacte et exemple camerounais analysé. Se relire : « Ai-je \textbf{exploité} l'auteur ? L'exemple \textbf{prouve}-t-il la thèse ? »
    \item \textbf{20 min :} Écrire \textbf{la conclusion} (réponse + résumé + ouverture). Vérifier : « Ma réponse contredit-elle mon développement ? L'ouverture est-elle pertinente ? »
\end{enumerate}
Faites ce rituel \textbf{3 fois par semaine} les 4 dernières semaines. C'est le secret des 16+.
\end{methode}

\newpage

\coursec{Activité d'intégration}

\coursec{Leçon 16 : L'existence et la mort}

\courssub{Objectifs de la leçon}
À l'issue de cette leçon, l'élève sera capable de :
\begin{itemize}
    \item définir rigoureusement les concepts d'\cle{existence} et de \cle{mort} ;
    \item expliquer et comparer trois grandes conceptions du sens de l'existence : la conception \cle{religieuse et traditionnelle}, la conception \cle{existentialiste} (Sartre) et la conception de l'\cle{absurde} (Camus) ;
    \item mobiliser ces notions pour analyser des situations concrètes camerounaises (funérailles traditionnelles, mortalité routière, engagement jeunesse) ;
    \item rédiger une \cle{introduction problématisée} et un \cle{paragraphe argumenté} conformes aux attendus de la dissertation au Baccalauréat technique ;
    \item s'autoévaluer à l'aide de la grille de critères officiels.
\end{itemize}

\courssub{Problématisation : l'existence face à la mort}
\textbf{Situation déclenchante : « Funérailles et sens de la vie au Cameroun ».}\par
Dans le cadre de la semaine philosophique de ton lycée technique de Bafoussam, ton professeur te remet un dossier de trois documents. Tu devras, à partir de ces documents, répondre à la question que se posent de nombreux jeunes de ta classe : \emph{à quoi sert de vivre, de travailler, de préparer le Probatoire, si l'on doit mourir un jour ?}

\begin{exemplebox}[Dossier documentaire]
\textbf{Document 1. Funérailles traditionnelles dans un village de l'Ouest (Bangangté).}\par
À la mort de Paï Nkouanka, 82 ans, cultivateur et notable, toute la chefferie s'est mobilisée pendant huit jours. Les funérailles ne sont pas vécues comme une disparition pure et simple : le défunt devient un \emph{ancêtre}, qui continue d'exister auprès des vivants et veille sur eux. Les danses rituelles, les sacrifices offerts, les pleurs codifiés des premières nuits puis les chants joyeux des derniers jours expriment une conviction : la mort n'est pas une fin absolue, mais un \emph{passage}. La vie du défunt prend sens à travers sa descendance, ses œuvres et son nom, que la famille s'engage à honorer. Un jeune du village déclare : « On pleure Paï, mais on célèbre aussi sa vie : il a planté des arbres dont nous mangeons les fruits. »

\textbf{Document 2. Albert Camus, \emph{Le Mythe de Sisyphe} (extrait adapté).}\par
« Il n'y a qu'un problème philosophique vraiment sérieux : c'est le suicide. Juger que la vie vaut ou ne vaut pas la peine d'être vécue, c'est répondre à la question fondamentale de la philosophie. [\ldots] Je dis que le monde est absurde, mais je jugeais trop vite. Ce monde lui-même n'est pas raisonnable, c'est tout ce qu'on peut en dire. Mais ce qui est absurde, c'est la confrontation de cet irrationnel et du désir éperdu de clarté dont l'appel résonne au plus profond de l'homme. [\ldots] Il faut imaginer Sisyphe heureux. »

\textbf{Document 3. Données de l'INS Cameroun (Annuaire statistique, données simplifiées).}\par
\begin{center}
\begin{tabularx}{\linewidth}{|l|X|X|}\hline
\rowcolor{popblueL} \textbf{Indicateur} & \textbf{Valeur} & \textbf{Observation} \\ \hline
Espérance de vie à la naissance & environ 60 ans & inférieure à la moyenne mondiale (environ 73 ans) \\ \hline
Décès annuels liés aux accidents de la route & environ 6 000 par an & l'une des premières causes de mortalité des 15--29 ans \\ \hline
Part des jeunes de moins de 25 ans dans la population & environ 63 \% & une population très jeune, confrontée tôt à la mort \\ \hline
\end{tabularx}
\end{center}
\legende{Source : INS Cameroun, données arrondies pour l'exercice pédagogique.}
\end{exemplebox}

\courssub{Définition des concepts fondamentaux}

\begin{definition}[Existence]
L'\cle{existence} (du latin \emph{ex-sistere}, « se tenir hors de », « apparaître ») désigne le mode d'être propre à l'homme. Contrairement à la simple \emph{vie biologique} (nutrition, reproduction, croissance), l'existence implique :
\begin{itemize}
    \item la \cle{conscience de soi} (je sais que je suis) ;
    \item la \cle{conscience de la finitude} (je sais que je vais mourir) ;
    \item la \cle{liberté} : l'homme doit \emph{choisir} sa vie, il n'a pas de nature fixée d'avance.
\end{itemize}
Exister, c'est être un \emph{projet} jeté dans le monde, responsable de ce qu'il devient.
\end{definition}

\begin{definition}[Mort]
La \cle{mort} est la \emph{cessation irréversible des fonctions vitales} (définition biologique). Pour la philosophie, elle est surtout un \cle{horizon ontologique} : elle structure l'existence humaine par sa certitude et son indétermination (on ne sait \emph{ni quand, ni comment}). Elle est la \emph{limite} qui donne du poids au temps et aux choix.
\end{definition}

\begin{attention}[Pièges à éviter]
\begin{itemize}
    \item Confondre \emph{vie} (biologique, partagée avec les animaux/plantes) et \emph{existence} (humaine, consciente, libre).
    \item Réduire la mort à un événement biologique sans en saisir la portée existentielle (angoisse, urgence, sens).
    \item Croire que « sens de la vie » = « bonheur » ; le sens peut inclure la souffrance, le sacrifice, la transmission.
\end{itemize}
\end{attention}

\courssub{Le sens de l'existence humaine : trois grandes conceptions}

\begin{aretenir}[Tableau synthèse des trois conceptions]
\begin{center}
\begin{tabularx}{\linewidth}{|l|X|X|X|}\hline
\rowcolor{popblueL} & \textbf{Religieuse / Traditionnelle} & \textbf{Existentialiste (Sartre)} & \textbf{Absurde (Camus)} \\ \hline
\textbf{Qu'est-ce que la mort ?} & Un \emph{passage} vers l'au-delà, le monde des ancêtres, la résurrection. & La \emph{fin définitive} qui clôt la liberté et fige l'essence. & Une \emph{fin sans au-delà}, silence du monde face à notre soif de sens. \\ \hline
\textbf{Quel sens pour l'existence ?} & \emph{Donné} (hétéro-nome) : par Dieu, les ancêtres, la communauté, la transmission. & \emph{Construit} (auto-nome) : « L'existence précède l'essence ». L'homme s'invente par ses choix. & \emph{Aucun sens donné}. Il faut vivre \emph{malgré tout}, par \cle{révolte}, \cle{liberté}, \cle{passion}. \\ \hline
\textbf{Figure emblématique} & L'Ancêtre (transmission), Le Juste (salut). & L'Homme libre (engagement), le « pour-soi ». & Sisyphe (lucide et révolté), l'homme absurde. \\ \hline
\textbf{Rôle de la conscience de la mort} & Préparer le passage, honorer les ancêtres, assurer la continuité. & Prendre conscience de la \cle{finitude} qui rend les choix urgents et définitifs. & Mesurer l'absurde, refuser le suicide philosophique (espoir illusoire) et le suicide physique. \\ \hline
\end{tabularx}
\end{center}
\end{aretenir}

\begin{experience}[Analyse guidée des documents du dossier]
\textbf{Document 1 (Tradition) :} La mort est un passage (ancêtre). Sens = transmission (arbres, descendance, nom). Continuité communautaire.
\textbf{Document 2 (Camus/Absurde) :} Le suicide = question philosophique n°1. Absurde = divorce homme/monde. Réponse = révolte lucide (« Sisyphe heureux »), pas l'espoir.
\textbf{Document 3 (Réalité camerounaise) :} Calcul : $6000 / 365 \approx 16,4$ morts/jour sur les routes. Fragilité extrême, jeunesse touchée. Urgence de donner du sens \emph{maintenant}.
\end{experience}

\courssub{Méthodologie : La dissertation au Baccalauréat technique}

\begin{methode}[Rédiger une introduction problématisée (4 étapes)]
\begin{enumerate}
    \item \textbf{Accroche} : fait concret, citation, statistique, situation vécue (ici : mortalité routière + funérailles).
    \item \textbf{Définitions} : préciser les termes du sujet (\cle{existence}, \cle{mort}, \cle{conscience}, \cle{sens}).
    \item \textbf{Problématique} : question unique, ouverte, issue de la tension entre les termes (ex : « La conscience de la mort donne-t-elle un sens à l'existence ou la rend-elle vaine ? »).
    \item \textbf{Annonce du plan} : 2 ou 3 grandes parties logiques (Thèse / Antithèse / Synthèse ou Analyse / Dépassement).
\end{enumerate}
\end{methode}

\begin{methode}[Rédiger un paragraphe argumenté (structure AREC)]
\begin{enumerate}
    \item \textbf{Idée directrice (Thèse)} : une phrase affirmative qui répond à la problématique.
    \item \textbf{Argument (Explication)} : développement logique, conceptuel (pourquoi ? comment ?).
    \item \textbf{Exemple / Illustration} : fait précis du dossier, référence culturelle, historique, littéraire.
    \item \textbf{Référence auteur} : citation ou idée précise (Camus, Sartre, Bible/Coran/Tradition, etc.).
    \item \textbf{Conclusion partielle} : phrase de transition qui valide l'idée et ouvre sur la suivante.
\end{enumerate}
\end{methode}

\courssub{Application : Production écrite guidée (Corrigé commenté)}

\begin{exoresolu}[Sujet : « La conscience de la mort donne-t-elle un sens à l'existence ? »]
\textbf{Étape 1 -- Introduction problématisée (modèle).}\par
« Chaque année au Cameroun, près de 6 000 vies sont fauchées sur les routes, majoritairement des jeunes de 15 à 29 ans, tandis que chaque famille célèbre des funérailles où l'on pleure autant qu'on danse : la mort n'est pas une abstraction, elle fait partie de notre quotidien. Or l'\emph{existence} désigne le mode d'être propre à l'homme, être conscient qui se sait mortel et qui doit choisir sa vie, tandis que la \emph{mort} est la cessation irréversible de la vie, horizon inévitable de toute existence. Dès lors, un problème se pose : la conscience que nous avons de notre finitude donne-t-elle un sens à notre vie, ou la prive-t-elle au contraire de tout sens en la rendant absurde ? Nous verrons d'abord que la mort peut sembler rendre l'existence vaine (thèse de l'absurde), puis qu'elle peut au contraire la rendre précieuse et fondatrice de valeurs (thèse existentialiste et traditionnelle), avant de montrer que le sens de l'existence dépend in fine de l'usage lucide que nous faisons de notre liberté face à la finitude. »

\textbf{Étape 2 -- Paragraphe argumenté (modèle : thèse existentialiste/traditionnelle).}\par
« La conscience de la mort donne un sens à l'existence, car c'est la finitude qui rend nos choix importants et urgents. \textbf{(Idée directrice)} En effet, si nous étions immortels, tout pourrait être remis au lendemain indéfiniment et aucun choix n'aurait de poids existentiel ; c'est parce que notre temps est limité que chaque décision engage notre vie tout entière et lui confère une valeur définitive. \textbf{(Argument)} Ainsi, au village Bangangté (Document 1), le défunt Paï Nkouanka est honoré pour ce qu'il a accompli concrètement -- ses champs, ses enfants, son nom -- : c'est bien parce que sa vie était finie qu'elle a pu prendre la forme d'une œuvre transmissible et honorée par la communauté. \textbf{(Exemple du dossier)} Comme le montre Sartre, « l'existence précède l'essence » : n'ayant aucune nature fixée d'avance par un créateur ou un destin, l'homme se définit uniquement par ses actes, et la mort est la limite qui donne à ces actes leur caractère irréversible et leur vérité. \textbf{(Référence auteur)} La conscience de la mort n'enlève donc pas le sens de la vie : elle nous oblige, au contraire, à le construire par notre engagement. \textbf{(Conclusion partielle)} »
\end{exoresolu}

\courssub{Entraînement individuel et auto-évaluation}

\begin{exercice}[À faire en classe / Devoir]
\begin{enumerate}
    \item Reprends le tableau comparatif (Religieuse / Existentialiste / Absurde) et complète-le avec tes propres mots.
    \item Rédige une introduction problématisée complète sur le sujet : « \emph{Faut-il craindre la mort pour bien vivre ?} ».
    \item Rédige un paragraphe argumenté (structure AREC) défendant l'idée que « La mort donne du prix à la vie » en t'appuyant sur le Document 1 et sur Sartre.
\end{enumerate}
\end{exercice}

\begin{corrige}[Grille d'auto-évaluation (Critères Bac Tech)]
Coche les cases après ta rédaction :
\begin{center}
\begin{tabularx}{\linewidth}{|X|c|c|}\hline
\rowcolor{popblueL} \textbf{Critère} & \textbf{Oui} & \textbf{Non} \\ \hline
Les termes du sujet sont définis avec précision (existence, mort, conscience, sens). & $\square$ & $\square$ \\ \hline
La problématique est une question claire, issue de la tension du sujet. & $\square$ & $\square$ \\ \hline
Le plan annoncé (2 ou 3 parties) répond logiquement à la problématique. & $\square$ & $\square$ \\ \hline
Chaque paragraphe suit la structure AREC (Idée -- Argument -- Exemple -- Référence -- Conclusion). & $\square$ & $\square$ \\ \hline
Au moins un auteur (Camus, Sartre) est cité ou référencé exactement. & $\square$ & $\square$ \\ \hline
Les exemples du dossier (Bangangté, statistiques INS, Sisyphe) sont mobilisés à bon escient. & $\square$ & $\square$ \\ \hline
L'expression est soignée (vocabulaire philosophique, connecteurs logiques, orthographe). & $\square$ & $\square$ \\ \hline
\end{tabularx}
\end{center}
\end{corrige}

\courssub{Bilan de la leçon}
Cette leçon a permis de mettre en évidence trois points essentiels pour le Probatoire :
\begin{enumerate}
    \item Les concepts d'\cle{existence} et de \cle{mort} ne sont pas de pures abstractions : ils s'incarnent dans des pratiques concrètes (funérailles traditionnelles de l'Ouest) et des réalités statistiques (mortalité routière, jeunesse camerounaise).
    \item Une même réalité -- la mort -- reçoit des interprétations radicalement différentes selon les conceptions : \emph{passage} vers les ancêtres (tradition), \emph{limite fondatrice de la liberté} (Sartre), \emph{scandale absurde à assumer lucidement} (Camus).
    \item L'élève a appris à transformer l'analyse de documents en une \cle{démarche argumentative} rigoureuse : définir, problématiser, comparer, argumenter avec exemples et références, conclure -- exactement les compétences exigées dans l'épreuve de dissertation de Philosophie au Baccalauréat technique.
\end{enumerate}

\begin{savaistu}[Le saviez-vous ?]
Au Cameroun, la coexistence des rites funéraires traditionnels (veillées, danses, sacrifices), des cérémonies chrétiennes (messe de requiem, inhumation au cimetière) et musulmanes (salat al-janaza, inhumation rapide) illustre concrètement le pluralisme des conceptions de la mort. Le philosophe camerounais \textbf{Fabien Eboussi Boulaga} a montré que la pensée africaine traditionnelle ne sépare pas radicalement vie et mort : les ancêtres sont les « vivants-morts » qui garantissent la cohésion du clan. Cette vision communautaire du sens (ubuntu : « Je suis parce que nous sommes ») dialogue aujourd'hui avec l'existentialisme individuel (Sartre) et la révolte solitaire (Camus).
\end{savaistu}

\newpage

\coursec{Méthodes et exercices résolus}

\coursec{L'existence et la mort}

\courssub{Savoir-faire 1 : Appliquer les notions d'existence et de mort à une situation concrète}

\begin{methode}[Analyser une situation de vie courante avec les concepts d'existence et de mort]
Pour traiter une situation concrète (un deuil, un accident, un suicide, un acharnement thérapeutique, un acte héroïque), suivre une démarche rigoureuse en cinq étapes :
\begin{enumerate}
\item \textbf{Lire et reformuler} la situation en une phrase simple : qui ? quoi ? quel problème humain ?
\item \textbf{Identifier le concept mobilisé} : s'agit-il de l'existence (le fait d'être, de vivre, de se réaliser) ou de la mort (cessation de la vie, finitude) ?
\item \textbf{Dégager le problème philosophique} : transformer le fait en question (exemple : « la mort d'un proche a-t-elle un sens ? » devient « la conscience de la mort donne-t-elle un sens à l'existence ? »).
\item \textbf{Mobiliser une thèse d'auteur} : \textbf{Épicure} (la mort n'est rien pour nous), \textbf{Heidegger} (l'être-pour-la-mort), \textbf{Sartre} (l'existence précède l'essence), ou une conception religieuse africaine (la vie continue auprès des ancêtres).
\item \textbf{Conclure par un jugement personnel justifié} : donner son avis argumenté, en lien avec la situation de départ.
\end{enumerate}
\textbf{Réflexe} : toujours boucler la boucle — la conclusion doit répondre à la situation initiale, pas à une autre.
\end{methode}

\begin{exoresolu}[Le deuil dans une famille camerounaise]
Énoncé. Dans un quartier de Yaoundé, une famille organise les funérailles d'un ancien. Pendant les cérémonies, les proches affirment que « le défunt n'est pas mort, il est parti rejoindre les ancêtres ». Un jeune lycéen, lui, déclare : « La mort, c'est la fin de tout. » Analysez cette situation en appliquant les notions d'existence et de mort.
\tcblower
Solution détaillée.
\begin{enumerate}
\item \textbf{Reformulation} : deux visions de la mort s'opposent au sein d'une même famille — la vision traditionnelle africaine (survie auprès des ancêtres) et la vision matérialiste (la mort comme anéantissement).
\item \textbf{Concept mobilisé} : il s'agit principalement du concept de \cle{mort} et de son rapport à l'\cle{existence} : la mort est-elle une fin absolue ou un passage ?
\item \textbf{Problème philosophique} : la mort met-elle fin à toute forme d'existence, ou l'existence humaine dépasse-t-elle la vie biologique ?
\item \textbf{Thèses mobilisées} :
\begin{itemize}
\item Du côté du lycéen, on peut convoquer \textbf{Épicure} : « la mort n'est rien pour nous, car quand nous sommes, la mort n'est pas là, et quand la mort est là, nous ne sommes plus ». La mort est la dissolution du corps et de la sensibilité : il n'y a rien à craindre après.
\item Du côté de la famille, la conception traditionnelle africaine soutient que la mort biologique n'abolit pas l'être : le défunt devient \emph{ancêtre}, il continue d'exister dans la mémoire et la vie de la communauté. On peut rapprocher cette idée de la distinction entre mort biologique et survie symbolique.
\end{itemize}
\item \textbf{Jugement personnel justifié} : les deux positions ne s'excluent pas totalement. Biologiquement, la mort est bien la fin de l'organisme ; mais humainement, un être continue d'exister à travers son œuvre, ses enfants, sa réputation. Le deuil camerounais, avec ses cérémonies, donne précisément un sens social à la mort et aide les vivants à continuer d'exister.
\end{enumerate}
\textbf{Conclusion} : cette situation montre que la mort n'est pas seulement un événement biologique : elle est aussi un événement culturel et philosophique qui interroge le sens de l'existence humaine.
\end{exoresolu}

\courssub{Savoir-faire 2 : Définir rigoureusement un concept philosophique}

\begin{methode}[Construire une définition philosophique correcte]
Une bonne définition au Bac respecte quatre étapes :
\begin{enumerate}
\item \textbf{Rapprocher le concept d'un genre proche} : « l'existence est le fait de... », « la mort est l'état de... ».
\item \textbf{Distinguer des notions voisines} : existence $\neq$ simple vie biologique ; mort $\neq$ sommeil ou maladie.
\item \textbf{Dégager les dimensions du concept} : pour l'existence — dimension biologique (être en vie), sociale (être reconnu), spirituelle (se réaliser) ; pour la mort — dimension biologique (arrêt des fonctions vitales), culturelle (rites), philosophique (finitude).
\item \textbf{Formuler la définition en une phrase claire et complète}, sans employer le mot à définir dans la définition.
\end{enumerate}
\textbf{Piège à éviter} : la définition circulaire (« la mort, c'est le fait de mourir ») est sanctionnée.
\end{methode}

\begin{exoresolu}[Définir l'existence et la mort]
Énoncé. Définissez de manière rigoureuse les concepts d'existence et de mort, en distinguant leurs différentes dimensions, puis montrez en une phrase le rapport qui les unit.
\tcblower
Solution détaillée.
\begin{enumerate}
\item \textbf{Genre proche} : l'existence est le \emph{fait d'être} ; la mort est un \emph{événement} ou un \emph{état}.
\item \textbf{Distinction des notions voisines} : exister ne se réduit pas à vivre au sens biologique (une plante vit, mais on ne dit pas qu'elle « existe » au sens humain) ; mourir n'est ni dormir ni être malade, car la mort est irréversible.
\item \textbf{Dimensions} :
\begin{itemize}
\item \textbf{Existence} : au sens biologique, le fait d'être en vie ; au sens social, le fait d'occuper une place parmi les hommes (travail, famille, citoyenneté) ; au sens spirituel, le fait de donner un sens à sa vie, de se réaliser par des choix libres.
\item \textbf{Mort} : au sens biologique, l'arrêt définitif des fonctions vitales (cœur, cerveau) ; au sens culturel, un passage marqué par des rites (funérailles, deuil) ; au sens philosophique, la finitude, c'est-à-dire la limite inscrite dans toute existence humaine.
\end{itemize}
\item \textbf{Définitions formulées} :
\begin{itemize}
\item L'\cle{existence} est le fait d'être au monde, de vivre en relation avec les autres et de construire sa propre vie par des choix.
\item La \cle{mort} est la cessation définitive et irréversible de la vie d'un être, qui marque la limite de son existence.
\end{itemize}
\end{enumerate}
\textbf{Conclusion} : existence et mort sont inséparables : c'est parce que l'homme sait qu'il mourra qu'il interroge le sens de son existence — la mort est, comme le dit Heidegger, la « possibilité » qui donne son urgence à la vie.
\end{exoresolu}

\courssub{Savoir-faire 3 : Problématiser un sujet de dissertation}

\begin{methode}[Transformer un sujet en problématique]
\begin{enumerate}
\item \textbf{Repérer les mots-clés} du sujet et les définir mentalement.
\item \textbf{Dégager la tension} : tout bon sujet de philosophie contient un conflit entre deux idées (exemple : vivre / mourir ; liberté / destin ; sens / absurdité).
\item \textbf{Formuler la problématique sous forme d'une question ouverte} (qui n'appelle pas simplement oui ou non), souvent de la forme : « Dans quelle mesure... ? », « Faut-il... ou... ? ».
\item \textbf{Vérifier} que la problématique annonce le plan : chaque partie de la dissertation devra apporter un élément de réponse.
\end{enumerate}
\textbf{Formule à retenir} : problématique = tension entre deux thèses + question qui les dépasse.
\end{methode}

\begin{exoresolu}[Problématiser le sujet : « La mort donne-t-elle un sens à la vie ? »]
Énoncé. Dégagez la problématique du sujet suivant : « La mort donne-t-elle un sens à la vie ? » et proposez les axes d'un plan en deux ou trois parties.
\tcblower
Solution détaillée.
\begin{enumerate}
\item \textbf{Mots-clés} : \emph{mort} (fin de la vie, finitude), \emph{sens} (signification, but, valeur), \emph{vie} (existence humaine).
\item \textbf{Tension} : d'un côté, la mort semble \emph{détruire} tout sens (à quoi bon agir si tout finit ? — thèse de l'absurde) ; de l'autre, elle semble au contraire \emph{créer} le sens (c'est parce que le temps est compté que chaque instant a du prix).
\item \textbf{Problématique} : « La mort, en limitant notre existence, la prive-t-elle de toute valeur, ou est-ce au contraire la conscience de notre finitude qui nous pousse à donner un sens à notre vie ? »
\item \textbf{Plan possible} :
\begin{itemize}
\item \textbf{I.} La mort semble ôter tout sens à l'existence : l'absurdité de la condition humaine (tout s'anéantit ; la mort rend nos efforts vains — on peut citer la vanité chez les moralistes).
\item \textbf{II.} Mais la conscience de la mort donne au contraire du prix à la vie : \textbf{Épicure} invite à ne pas la craindre et à jouir du présent ; \textbf{Heidegger} montre que l'« être-pour-la-mort » rend l'homme authentique et responsable de ses choix.
\item \textbf{III.} Synthèse : le sens de l'existence ne dépend pas de la durée mais de l'usage que nous faisons de notre liberté (\textbf{Sartre} : l'homme se définit par ses actes).
\end{itemize}
\end{enumerate}
\textbf{Conclusion} : une bonne problématique oppose deux réponses possibles et prépare un plan progressif qui va du constat à la réflexion personnelle.
\end{exoresolu}

\courssub{Savoir-faire 4 : Rédiger et justifier une démarche argumentative}

\begin{methode}[Construire un paragraphe argumenté (méthode T-E-A-J)]
Chaque paragraphe de dissertation ou de réponse construite suit quatre temps :
\begin{enumerate}
\item \textbf{T — Thèse} : énoncer clairement l'idée défendue dans la partie (« La conscience de la mort rend l'existence précieuse »).
\item \textbf{E — Explication} : développer l'idée avec ses propres mots, définir les termes.
\item \textbf{A — Appui} : citer un auteur (brièvement) ou un argument rationnel, puis un \textbf{exemple concret} (de préférence camerounais ou quotidien).
\item \textbf{J — Jugement} : conclure le paragraphe par une phrase de bilan qui relie l'idée à la problématique.
\end{enumerate}
\textbf{Réflexes} : une seule idée par paragraphe ; des connecteurs logiques (d'abord, en effet, cependant, par conséquent) ; jamais de citation sans explication.
\end{methode}

\begin{exoresolu}[Rédiger un paragraphe argumenté]
Énoncé. Rédigez un paragraphe argumenté complet (8 à 12 lignes) répondant à la question : « Faut-il craindre la mort ? »
\tcblower
Solution détaillée (application de T-E-A-J).

\textbf{(T) Thèse} : Non, l'homme raisonnable ne doit pas craindre la mort, car la craindre, c'est gâcher la seule vie dont nous disposons.

\textbf{(E) Explication} : La peur de la mort repose sur une illusion : nous imaginons la mort comme une expérience douloureuse que nous vivrions. Or la mort n'est précisément pas une expérience : elle est l'absence de toute sensation. Redouter ce que nous ne ressentirons jamais est donc une contradiction.

\textbf{(A) Appui} : C'est l'argument célèbre d'\textbf{Épicure} : « Quand nous sommes, la mort n'est pas là, et quand la mort est là, nous ne sommes plus. » La mort ne nous concernant jamais directement, il est absurde d'en avoir peur. \emph{Exemple concret} : au Cameroun, les cérémonies funéraires, loin d'être seulement des moments de terreur, sont des fêtes où la communauté célèbre la vie du défunt et se réconcilie avec l'idée de la mort ; le deuil bien conduit apaise la peur au lieu de l'entretenir.

\textbf{(J) Jugement} : Ainsi, craindre la mort revient à laisser un événement absent empoisonner un présent bien réel ; la sagesse consiste à accepter notre finitude pour mieux vivre.

\textbf{Conclusion} : le paragraphe respecte les quatre temps, cite un auteur expliqué, illustre par un exemple local et se rattache à la question posée.
\end{exoresolu}

\courssub{Savoir-faire 5 : Confronter deux conceptions du sens de l'existence}

\begin{methode}[Organiser une confrontation de thèses]
\begin{enumerate}
\item \textbf{Présenter la première conception} : son principe, son auteur ou sa tradition, sa réponse à la question du sens de la vie.
\item \textbf{Présenter la conception opposée} avec les mêmes critères (symétrie de traitement : ne pas favoriser l'une).
\item \textbf{Comparer explicitement} : points de désaccord (d'où vient le sens ? de Dieu, de la liberté, du devoir, du bonheur ?) et points de convergence éventuels.
\item \textbf{Trancher avec un jugement personnel argumenté}, éventuellement en dépassant l'opposition.
\end{enumerate}
\textbf{Grille de comparaison à retenir} : origine du sens (donné ou construit ?), place de la mort, place de la liberté, conséquences pratiques.
\end{methode}

\begin{exoresolu}[Conception religieuse et conception existentialiste du sens de la vie]
Énoncé. Confrontez la conception religieuse du sens de l'existence et la conception existentialiste (Sartre : « l'existence précède l'essence »), puis défendez votre position.
\tcblower
Solution détaillée.
\begin{enumerate}
\item \textbf{Conception religieuse} : pour les croyants (christianisme, islam, religions traditionnelles africaines), le sens de l'existence est \emph{donné} par Dieu ou par l'ordre du monde : l'homme est créé pour une fin (adorer Dieu, servir la communauté, rejoindre les ancêtres). La mort n'est pas une fin absolue mais un passage ; la vie a donc un sens qui la dépasse.
\item \textbf{Conception existentialiste} : pour \textbf{Sartre}, « l'existence précède l'essence » : l'homme naît sans nature définie, sans mission préétablie ; il est « condamné à être libre » et doit \emph{construire lui-même} le sens de sa vie par ses choix et ses actes. La mort étant la fin de tout, le sens ne peut venir que de l'existence elle-même.
\item \textbf{Comparaison} :
\begin{itemize}
\item Désaccord sur l'origine du sens : donné d'en haut (religion) contre fabriqué par la liberté (existentialisme).
\item Désaccord sur la mort : passage vers une autre vie contre fin définitive.
\item Convergence : dans les deux cas, l'homme est \emph{responsable} de sa conduite — devant Dieu et la communauté, ou devant lui-même et les autres.
\end{itemize}
\item \textbf{Jugement personnel} : on peut soutenir que, même pour un croyant, le sens donné par Dieu doit être \emph{assumé librement} : une vie n'a de valeur que si l'on choisit d'y répondre par des actes. La liberté et la foi ne s'opposent donc pas nécessairement : l'essentiel est que chacun donne à son existence une direction au lieu de la subir.
\end{enumerate}
\textbf{Conclusion} : que le sens soit reçu ou construit, l'existence humaine n'a de valeur que vécue de manière responsable — c'est là le point de rencontre des deux conceptions.
\end{exoresolu}

\courssub{Savoir-faire 6 : Rédiger une introduction et une conclusion de dissertation}

\begin{methode}[Introduction en 4 temps et conclusion en 2 temps]
\textbf{Introduction} :
\begin{enumerate}
\item \textbf{Amorce} : partir d'un constat général ou d'une situation concrète liée au sujet.
\item \textbf{Analyse des termes} : définir brièvement les mots-clés du sujet.
\item \textbf{Problématique} : poser la question philosophique (la tension).
\item \textbf{Annonce du plan} : présenter les deux ou trois parties en une phrase.
\end{enumerate}
\textbf{Conclusion} :
\begin{enumerate}
\item \textbf{Bilan} : répondre clairement à la problématique en résumant le chemin parcouru.
\item \textbf{Ouverture} : poser une question nouvelle qui prolonge la réflexion (sans la traiter).
\end{enumerate}
\textbf{Interdits} : pas de définition de dictionnaire en amorce, pas de nouvelle thèse dans la conclusion, pas de plan annoncé au style télégraphique (« Dans un premier temps... dans un deuxième temps... » sans lien avec le sujet).
\end{methode}

\begin{exoresolu}[Introduction et conclusion sur « L'existence humaine a-t-elle un sens ? »]
Énoncé. Rédigez l'introduction et la conclusion d'une dissertation sur le sujet : « L'existence humaine a-t-elle un sens ? »
\tcblower
Solution détaillée.

\textbf{Introduction.} \emph{(Amorce)} Chaque jour, au Cameroun comme partout, des hommes se lèvent pour travailler, élever leurs enfants, construire leur avenir, tout en sachant qu'ils mourront un jour. \emph{(Analyse des termes)} L'existence désigne le fait d'être au monde et de vivre parmi les autres ; le sens renvoie à la fois à la signification (comprendre pourquoi) et à la direction (vers quel but). \emph{(Problématique)} Mais cette vie limitée par la mort possède-t-elle une signification donnée d'avance, ou chacun doit-il en inventer le sens par sa liberté ? \emph{(Annonce du plan)} Nous verrons d'abord que l'existence peut sembler absurde, puis que diverses conceptions lui attribuent un sens, avant de montrer que le sens se construit par l'engagement responsable.

\textbf{Conclusion.} \emph{(Bilan)} Notre réflexion a montré que si la mort menace toujours de rendre l'existence vaine, elle la rend aussi précieuse : le sens de la vie n'est ni un donné à subir ni une illusion, mais une tâche à accomplir par nos choix, notre travail et notre engagement envers les autres. \emph{(Ouverture)} Reste alors à se demander si la recherche du bonheur suffit à remplir cette tâche, ou si le sens exige aussi le devoir.

\textbf{Conclusion de l'exercice} : l'introduction respecte les quatre temps et la conclusion répond à la problématique tout en ouvrant sur la question du bonheur et du devoir.
\end{exoresolu}

\begin{attention}[Les 5 erreurs qui coûtent des points au Bac]
\begin{enumerate}
\item \textbf{La définition circulaire} : écrire « la mort, c'est mourir » ou « exister, c'est avoir une existence ». Une définition doit passer par un genre proche et des dimensions précises, sans reprendre le mot défini.
\item \textbf{La paraphrase du sujet} : recopier le sujet en introduction sans définir les termes ni dégager de tension. Sans problématique explicite (une vraie question), la copie ne peut pas dépasser la moyenne.
\item \textbf{La citation sans explication} : écrire « Épicure a dit que la mort n'est rien pour nous » puis passer à autre chose. Toute référence doit être expliquée et reliée à l'argumentation ; une citation non commentée est considérée comme du remplissage.
\item \textbf{Le hors-sujet existentiel} : transformer la dissertation en récit personnel (histoire familiale, témoignage de deuil) sans aucun concept ni raisonnement. Les exemples concrets doivent \emph{illustrer} une thèse, jamais la remplacer.
\item \textbf{La conclusion absente ou contradictoire} : terminer sans répondre à la problématique, ou défendre en conclusion le contraire de ce que le développement a montré. La conclusion doit être le bilan logique du chemin parcouru, suivi d'une ouverture courte.
\end{enumerate}
\end{attention}

\newpage

\coursec{Exercices}

\courssub{Je vérifie mes connaissances}

\begin{exercice}[QCM : les notions de base]
Pour chaque question, une seule réponse est correcte. Entoure la lettre correspondante.
\begin{enumerate}
\item L'\emph{existence}, au sens strict, désigne :
\begin{enumerate}
\item[a)] le simple fait d'être vivant biologiquement ;
\item[b)] le fait pour l'homme d'être là, de se définir lui-même par ses choix ;
\item[c)] l'ensemble des choses matérielles du monde.
\end{enumerate}
\item Pour Sartre, « l'existence précède l'essence » signifie :
\begin{enumerate}
\item[a)] l'homme naît avec une nature déjà fixée par Dieu ;
\item[b)] l'homme existe d'abord et se définit ensuite par ses actes ;
\item[c)] l'essence de l'homme est de mourir.
\end{enumerate}
\item Le \emph{Dasein} chez Heidegger désigne :
\begin{enumerate}
\item[a)] la vie animale ;
\item[b)] l'être-là humain, celui qui se questionne sur son être ;
\item[c)] l'âme immortelle.
\end{enumerate}
\item Selon Épicure, il ne faut pas craindre la mort parce que :
\begin{enumerate}
\item[a)] l'âme est immortelle ;
\item[b)] la mort n'est rien pour nous : tant que nous sommes, elle n'est pas, et quand elle est, nous ne sommes plus ;
\item[c)] la mort est un passage vers une vie meilleure.
\end{enumerate}
\end{enumerate}
\end{exercice}

\begin{exercice}[Vrai ou faux justifié]
Indique si chaque affirmation est vraie ou fausse, puis justifie ta réponse en deux ou trois phrases.
\begin{enumerate}
\item La mort biologique et la mort vécue comme angoisse sont une seule et même chose.
\item Pour Heidegger, l'homme authentique est un « être-pour-la-mort » : anticiper sa mort donne sens à sa vie.
\item L'absurde, chez Camus, naît de la confrontation entre l'exigence humaine de sens et le silence du monde.
\item Exister et vivre sont des termes parfaitement synonymes en philosophie.
\end{enumerate}
\end{exercice}

\begin{exercice}[Définitions]
Définis en trois lignes maximum chacune des notions suivantes, en donnant pour chacune un exemple tiré de la vie courante camerounaise :
\begin{enumerate}
\item l'existence ;
\item la mort ;
\item l'être-pour-la-mort ;
\item le sens de la vie.
\end{enumerate}
\end{exercice}

\begin{exercice}[Auteurs et thèses]
Complète le tableau suivant en associant à chaque auteur sa thèse sur l'existence ou la mort, puis rédige une phrase qui résume chaque thèse avec tes propres mots.
\begin{center}
\begin{tabularx}{\linewidth}{|l|X|}
\hline
\rowcolor{popblueL} \textbf{Auteur} & \textbf{Thèse} \\
\hline
Épicure & \ldots \\
\hline
Heidegger & \ldots \\
\hline
Sartre & \ldots \\
\hline
Camus & \ldots \\
\hline
\end{tabularx}
\end{center}
\end{exercice}

\courssub{Je m'entraîne}

\begin{exercice}[L'accident de la N3]
À la sortie de Douala, un jeune conducteur de moto-taxi perd la vie dans un accident de la route. Ses camarades, réunis au carrefour, s'interrogent : « À quoi bon travailler, se lever chaque matin, si tout peut s'arrêter en une seconde ? »
\begin{enumerate}
\item Identifie les deux notions philosophiques en jeu dans cette situation.
\item Rédige un paragraphe argumenté (dix à quinze lignes) dans lequel tu mobilises la pensée d'Épicure \emph{ou} celle de Heidegger pour répondre à la question des camarades.
\end{enumerate}
\end{exercice}

\begin{exercice}[Débat encadré : faut-il craindre la mort ?]
La classe se divise en deux camps. Le camp A défend la thèse d'Épicure : « La mort n'est rien pour nous, il est insensé de la craindre. » Le camp B défend la thèse de Heidegger : « Penser la mort, c'est apprendre à vivre authentiquement. »
\begin{enumerate}
\item Chaque camp rédige deux arguments et un exemple concret tiré de la réalité camerounaise (deuils, funérailles, croyances traditionnelles ou religieuses).
\item Rédige individuellement une conclusion personnelle justifiée en cinq lignes.
\end{enumerate}
\end{exercice}

\begin{exercice}[Existence et engagement]
Dans un quartier de Yaoundé, un jeune diplômé sans emploi décide de créer une coopérative de recyclage plastique avec ses amis. Il déclare : « Personne ne décidera à ma place de ce que ma vie doit être. »
\begin{enumerate}
\item Montre en quoi cette attitude illustre la formule de Sartre : « l'existence précède l'essence ».
\item Ce jeune est-il, selon toi, responsable du sens de sa propre existence ? Justifie ta réponse en un paragraphe argumenté.
\end{enumerate}
\end{exercice}

\begin{exercice}[Vivre ou seulement exister ?]
On considère les deux situations suivantes : (a) un homme passe ses journées à regarder passer les motos sans aucun projet ; (b) une mère de famille de Bafoussam se bat chaque jour pour scolariser ses enfants et développer son champ de maïs.
\begin{enumerate}
\item À partir de ces deux cas, établis une distinction claire entre « vivre » au sens biologique et « exister » au sens philosophique.
\item Laquelle des deux personnes donne, selon toi, un sens plus fort à son existence ? Argumente en mobilisant une notion du cours.
\end{enumerate}
\end{exercice}

\courssub{Je me prépare au Bac}

\begin{exercice}[Dissertation : la mort et le sens de l'existence]
\textbf{Sujet :} \emph{La mort retire-t-elle tout sens à l'existence humaine ?}

En t'appuyant sur les notions du cours (existence, mort, être-pour-la-mort, absurde) et sur au moins deux auteurs étudiés (Épicure, Heidegger, Sartre, Camus), tu rédigeras une dissertation complète comportant :
\begin{enumerate}
\item une introduction (accroche, présentation du sujet, problématique, annonce du plan) ;
\item un développement en trois parties : (I) la mort semble anéantir tout sens ; (II) mais la conscience de la mort peut au contraire fonder le sens de l'existence ; (III) position personnelle argumentée ;
\item une conclusion qui répond clairement à la problématique et propose une ouverture.
\end{enumerate}
Illustre ton propos par au moins deux exemples tirés de la vie courante au Cameroun.
\end{exercice}

\begin{exercice}[Analyse de texte : Sartre]
\textbf{Texte :} « L'existence précède l'essence. [\ldots] L'homme existe d'abord, se rencontre, surgit dans le monde, et se définit après. L'homme, tel que le conçoit l'existentialiste, s'il n'est pas définissable, c'est qu'il n'est d'abord rien. Il ne sera qu'ensuite, et il sera tel qu'il se sera fait. Ainsi, il n'y a pas de nature humaine, puisqu'il n'y a pas de Dieu pour la concevoir. L'homme est non seulement tel qu'il se conçoit, mais tel qu'il se veut. » (Jean-Paul Sartre, \emph{L'Existentialisme est un humanisme}, adapté.)

\begin{enumerate}
\item Dégage la thèse défendue par Sartre dans ce texte et reformule-la en une phrase.
\item Explique la formule « l'existence précède l'essence » à l'aide d'un exemple simple.
\item Releve dans le texte deux arguments qui soutiennent la thèse.
\item Ce texte implique-t-il que l'homme est entièrement responsable de ce qu'il devient ? Justifie ta réponse.
\item Donne ton avis personnel argumenté : un jeune Camerounais peut-il toujours « se faire » librement malgré les contraintes sociales et économiques ?
\end{enumerate}
\end{exercice}

\begin{exercice}[Rédaction d'introduction et construction de problématique]
\textbf{Sujet :} \emph{Vivre, est-ce seulement exister ?}

\begin{enumerate}
\item Définis les deux termes du sujet (« vivre » et « exister ») et montre qu'ils ne sont pas synonymes.
\item Formule la contradiction qui fait problème dans ce sujet, puis rédige la problématique sous forme d'une question.
\item Rédige l'introduction complète de la dissertation (accroche, analyse des termes, problématique, annonce d'un plan en trois parties).
\item Propose, pour chaque partie annoncée, un argument principal, un auteur de référence et un exemple concret camerounais.
\end{enumerate}
\end{exercice}

\newpage

\coursec{Corrigés des exercices}

\begin{corrige}[Exercice 1 — QCM : les notions de base]
\begin{enumerate}
\item \textbf{Réponse b).} Au sens strictement philosophique, l'\cle{existence} ne se confond pas avec la vie biologique : elle désigne le fait pour l'homme d'« être-là », de se projeter et de se définir par ses choix. La réponse a) décrit la simple \emph{vie} (végétative et animale), et la réponse c) décrit l'\emph{étant} en général, non l'existence humaine.
\item \textbf{Réponse b).} Chez Sartre, l'homme surgit d'abord dans le monde, puis se construit par ses actes : il n'existe pas de nature humaine préétablie. La réponse a) correspond à la thèse essentialiste (et théologique) que Sartre réfute ; la réponse c) est un contresens.
\item \textbf{Réponse b).} Le \emph{Dasein} (« être-là ») désigne chez Heidegger l'être humain en tant qu'il est le seul étant capable de s'interroger sur le sens de son être.
\item \textbf{Réponse b).} Pour Épicure, la crainte de la mort est irrationnelle : « la mort n'est rien pour nous », car nous ne la rencontrons jamais. Les réponses a) et c) supposent une immortalité de l'âme qu'Épicure, matérialiste, rejette.
\end{enumerate}
\textbf{Points de vigilance :} bien distinguer \emph{vivre} (fait biologique) et \emph{exister} (fait proprement humain) ; ne pas attribuer à Épicure une croyance en l'au-delà.
\end{corrige}

\begin{corrige}[Exercice 2 — Vrai ou faux justifié]
\begin{enumerate}
\item \textbf{Faux.} La mort biologique est un événement organique (arrêt des fonctions vitales) ; la mort vécue comme angoisse est un rapport de la conscience à sa propre finitude. On peut mourir sans jamais avoir pensé sa mort, et l'on peut penser sa mort toute sa vie sans être mourant.
\item \textbf{Vrai.} Pour Heidegger, anticiper sa mort (« être-pour-la-mort ») arrache l'homme à l'anonymat du « on » et lui permet d'assumer ses possibles : la conscience de la finitude donne sérieux et orientation à l'existence.
\item \textbf{Vrai.} Chez Camus, l'absurde n'est ni dans l'homme seul ni dans le monde seul : il naît de leur confrontation, entre l'appel humain au sens et le silence déraisonnable du monde.
\item \textbf{Faux.} \emph{Vivre} désigne le fait biologique commun à tous les organismes ; \emph{exister} désigne le mode d'être spécifiquement humain, fait de conscience, de liberté et de projet. Tout homme vit, mais tous n'existent pas au sens fort : certains se contentent de subsister sans projet.
\end{enumerate}
\end{corrige}

\begin{corrige}[Exercice 3 — Définitions]
\begin{enumerate}
\item \textbf{L'existence :} mode d'être propre à l'homme, qui n'est pas seulement vivant mais conscient de lui-même, libre et capable de se définir par ses choix et ses projets. \emph{Exemple :} un jeune de Douala qui refuse l'oisiveté et monte un atelier de couture pour construire son avenir.
\item \textbf{La mort :} cessation définitive de la vie biologique ; pour la philosophie, elle est aussi l'horizon de finitude qui donne son urgence et sa valeur à l'existence. \emph{Exemple :} les cérémonies de deuil et funérailles qui, dans nos villages, rappellent à chacun la fragilité de la vie.
\item \textbf{L'être-pour-la-mort :} formule de Heidegger désignant l'attitude de celui qui anticipe sa propre mort et, par là, assume authentiquement son existence au lieu de fuir dans la distraction. \emph{Exemple :} un chef de famille qui, après un deuil, organise sérieusement la succession et l'éducation de ses enfants.
\item \textbf{Le sens de la vie :} la signification, la valeur ou la finalité qu'un homme donne à son existence par ses projets, ses engagements et ses choix. \emph{Exemple :} une mère de Bafoussam qui trouve le sens de sa vie dans la scolarisation de ses enfants.
\end{enumerate}
\end{corrige}

\begin{corrige}[Exercice 4 — Auteurs et thèses]
\begin{center}
\begin{tabularx}{\linewidth}{|l|X|}
\hline
\rowcolor{popblueL} \textbf{Auteur} & \textbf{Thèse} \\
\hline
Épicure & La mort n'est rien pour nous : tant que nous vivons elle n'est pas là, et quand elle est là nous ne sommes plus ; il est donc insensé de la craindre. \\
\hline
Heidegger & L'homme est un « être-pour-la-mort » : anticiper sa finitude rend l'existence authentique et lui donne son sérieux. \\
\hline
Sartre & L'existence précède l'essence : l'homme n'a pas de nature préétablie, il se définit lui-même par ses actes et est responsable de ce qu'il devient. \\
\hline
Camus & L'absurde naît du conflit entre l'exigence humaine de sens et le silence du monde ; il faut vivre et créer malgré l'absurde, sans espoir d'au-delà. \\
\hline
\end{tabularx}
\end{center}
\textbf{Reformulations possibles :} Épicure : « Craindre la mort, c'est craindre quelque chose que l'on ne rencontrera jamais. » Heidegger : « Penser sa mort, c'est apprendre à vivre vraiment. » Sartre : « Nous ne sommes pas nés tout faits : nous nous fabriquons nous-mêmes. » Camus : « Le monde ne répond pas à notre soif de sens, mais c'est à nous de vivre pleinement quand même. »
\end{corrige}

\begin{corrige}[Exercice 5 — L'accident de la N3]
\textbf{1.} Les deux notions en jeu sont l'\cle{existence} (que vaut une vie humaine, quel sens lui donner ?) et la \cle{mort} (sa brutalité, son imprévisibilité, qui semble anéantir tous les projets).

\textbf{2. Paragraphe argumenté (piste Heidegger).} La question des camarades — « à quoi bon ? » — suppose que seule une vie sans fin mériterait d'être vécue. Or c'est précisément l'inverse que soutient Heidegger. Pour lui, l'homme est un « être-pour-la-mort » : c'est parce que notre temps est compté que nos choix ont du poids. Une existence sans terme serait une existence indifférente, où tout pourrait être remis à demain. La mort du jeune conducteur ne prouve donc pas l'inutilité de l'effort ; elle rappelle au contraire que chaque matin où l'on se lève pour travailler est un matin où l'on décide de soi. Fuir cette vérité dans la distraction, c'est vivre de façon inauthentique, selon le « on ». L'assumer, c'est donner à ses projets — le travail, la famille, le quartier — un sérieux et une urgence. La mort ne retire pas le sens de l'existence : elle en est la condition. \emph{(Variante Épicure acceptée : tant que nous sommes vivants, la mort n'est pas ; il faut donc consacrer notre énergie à la vie présente, à l'amitié et au travail, au lieu de nous paralyser par une crainte vaine.)}

\textbf{Points de vigilance :} citer explicitement l'auteur choisi, mobiliser au moins une notion du cours (\emph{être-pour-la-mort} ou \emph{la mort n'est rien pour nous}), et conclure en répondant effectivement à la question des camarades.
\end{corrige}

\begin{corrige}[Exercice 6 — Débat encadré : faut-il craindre la mort ?]
\textbf{1. Arguments du camp A (Épicure).} \emph{Argument 1 :} la crainte de la mort est irrationnelle, car on ne peut craindre que ce qu'on éprouvera ; or la mort, nous ne l'éprouverons jamais. \emph{Argument 2 :} cette crainte gâche le seul bien que nous possédions, la vie présente ; celui qui l'écarte peut goûter l'amitié et le plaisir simple. \emph{Exemple :} au village, les veillées mortuaires où l'on mange, chante et raconte le défunt montrent que la communauté choisit de célébrer la vie plutôt que de s'effondrer dans la peur.

\textbf{Arguments du camp B (Heidegger).} \emph{Argument 1 :} fuir la pensée de la mort, c'est vivre selon le « on », sans jamais décider soi-même ; l'anticipation de la mort individualise et responsabilise. \emph{Argument 2 :} la conscience de la finitude donne de la valeur au temps : c'est parce que les jours sont comptés que les projets comptent. \emph{Exemple :} dans les traditions camerounaises, les rites funéraires et le culte des ancêtres obligent chacun à se situer dans une filiation et à réfléchir à l'héritage qu'il laissera.

\textbf{2. Conclusion personnelle (modèle).} Les deux thèses ne s'excluent pas totalement : Épicure a raison de dénoncer la peur panique qui paralyse, mais Heidegger a raison de montrer que l'oubli de la mort rend la vie frivole. Il ne faut donc ni craindre la mort ni l'ignorer : l'accepter comme horizon, c'est apprendre à vivre pleinement et librement le temps qui nous est donné.
\end{corrige}

\begin{corrige}[Exercice 7 — Existence et engagement]
\textbf{1.} La formule de Sartre signifie que l'homme n'a pas de nature fixée d'avance : il existe d'abord, puis se définit par ses actes. Le jeune diplômé illustre cette thèse : rien dans sa situation (le chômage, l'attente d'un emploi « tout fait ») ne détermine ce qu'il doit être. En créant sa coopérative, il refuse de se laisser définir par les circonstances ou par le regard des autres : il se fait lui-même par son projet. Sa phrase — « personne ne décidera à ma place » — exprime exactement la liberté comme autodéfinition.

\textbf{2.} Oui, selon Sartre, il est responsable du sens de son existence : puisqu'il n'y a pas d'essence humaine préalable, chacun est « condamné à être libre » et répond de ce qu'il fait de lui-même. Ce jeune ne peut pas rejeter la faute de son oisiveté sur la société s'il choisit de ne rien tenter ; inversement, son engagement fait de lui un homme qui construit, pour lui et pour les autres. On peut nuancer : les conditions économiques pèsent réellement. Mais la responsabilité sartrienne ne signifie pas que tout est facile ; elle signifie que, même dans la contrainte, c'est moi qui décide du sens que je donne à ma situation — subir ou entreprendre.
\end{corrige}

\begin{corrige}[Exercice 8 — Vivre ou seulement exister ?]
\textbf{1.} \emph{Vivre} au sens biologique, c'est accomplir les fonctions de l'organisme : manger, dormir, se reproduire, durer. L'homme (a) vit pleinement en ce sens : son corps fonctionne. \emph{Exister} au sens philosophique, c'est en plus se rapporter à soi-même, choisir, se projeter, donner une signification à ses jours. La mère de Bafoussam (b) existe au sens fort : elle a un projet (scolariser ses enfants), elle assume des responsabilités, elle se définit par un combat. On peut donc vivre sans véritablement exister : la vie est la condition de l'existence, mais ne lui suffit pas.

\textbf{2.} C'est la mère de famille qui donne le sens le plus fort à son existence. En mobilisant la notion de \cle{sens de la vie} (et la thèse de Sartre), on peut dire qu'elle se définit par ses actes et son engagement, là où le premier personnage se laisse porter par le temps, proche de ce que Heidegger appelle la vie inauthentique du « on ». Son existence a une direction, une finalité, une valeur pour elle-même et pour autrui : c'est cela, exister et non seulement vivre.
\end{corrige}

\begin{corrige}[Exercice 9 — Dissertation : la mort retire-t-elle tout sens à l'existence ?]
\textbf{Barème indicatif (sur 20) :} introduction (accroche, analyse des termes, problématique, plan) : 4 pts ; développement en trois parties équilibrées, avec auteurs et exemples : 12 pts ; conclusion avec réponse et ouverture : 2 pts ; langue et organisation : 2 pts.

\textbf{Corrigé rédigé (tracé détaillé).}

\emph{Introduction.} Accroche : au Cameroun, les funérailles occupent une place centrale — on y pleure, mais on y célèbre aussi la vie du défunt. Présentation : tout homme sait qu'il mourra ; cette certitude semble menacer la valeur de tout ce qu'il entreprend. Problématique : la mort, en mettant fin à tous nos projets, retire-t-elle tout sens à l'existence, ou est-ce au contraire elle qui rend le sens possible ? Plan : nous verrons d'abord que la mort paraît anéantir tout sens (I), puis que sa conscience peut fonder le sens (II), enfin nous défendrons une position personnelle (III).

\emph{I. La mort semble anéantir tout sens.} Argument 1 : la mort détruit tous les projets — le jeune tué sur la N3 ne verra jamais le fruit de son travail ; « à quoi bon ? » demandent ses camarades. Argument 2 : le sentiment de l'\cle{absurde} chez Camus : l'homme crie au monde sa demande de sens, le monde répond par le silence ; la mort scelle ce silence. Argument 3 : l'oubli — les générations passent et effacent les noms.

\emph{II. Mais la conscience de la mort peut fonder le sens.} Argument 1 : Épicure — bien comprise, la mort n'est rien pour nous ; cesser de la craindre libère la joie de vivre le présent. Argument 2 : Heidegger — l'homme est « être-pour-la-mort » : anticiper sa finitude arrache à la vie inauthentique du « on » et donne son sérieux à chaque choix ; c'est parce que le temps est compté qu'il a du prix. Argument 3 : Sartre — puisque rien n'est écrit d'avance, c'est à moi de faire sens avant de mourir ; l'existence précède l'essence. Exemple : la mère de Bafoussam qui se bat pour ses enfants donne à sa vie une valeur que la mort ne pourra pas effacer de la mémoire des siens.

\emph{III. Position personnelle.} La mort ne retire le sens qu'à celui qui attend un sens éternel et tout fait. Si le sens est quelque chose que l'existence produit et non qu'elle reçoit, alors la finitude n'est pas un scandale mais une condition : elle rend nos choix urgents, rares, précieux. Le deuil camerounais, qui honore le défunt en réaffirmant les liens des vivants, montre que la mort peut rassembler autour de ce qui compte. La mort limite l'existence, mais c'est la limite qui dessine la figure.

\emph{Conclusion.} La mort ne retire donc pas tout sens à l'existence : elle le menace, mais cette menace même est ce qui nous oblige à en créer. Ouverture : reste à savoir si le sens que nous créons peut survivre à notre mort — question de la transmission et de l'héritage.

\textbf{Points de vigilance :} au moins deux auteurs cités avec leur thèse exacte ; deux exemples camerounais ; la troisième partie doit trancher, pas répéter ; la conclusion répond mot pour mot à la problématique.
\end{corrige}

\begin{corrige}[Exercice 10 — Analyse de texte : Sartre]
\textbf{Barème indicatif :} question 1 : 3 pts ; question 2 : 4 pts ; question 3 : 4 pts ; question 4 : 4 pts ; question 5 : 5 pts.

\textbf{1.} Thèse : l'homme n'a pas de nature préétablie ; il existe d'abord, puis se définit lui-même par ses choix et ses actes — il est ce qu'il se fait.

\textbf{2.} « L'existence précède l'essence » signifie que, contrairement à un objet fabriqué (le tailleur fabrique un couteau selon un modèle préalable : l'essence précède l'existence), l'homme surgit dans le monde sans modèle. \emph{Exemple :} un élève de Terminale technique n'est pas « né » menuisier ou commerçant : c'est par ses efforts, ses stages, ses décisions qu'il deviendra l'un ou l'autre.

\textbf{3.} Deux arguments du texte : (a) « il n'y a pas de nature humaine, puisqu'il n'y a pas de Dieu pour la concevoir » — sans créateur, pas de définition préalable de l'homme ; (b) « il ne sera qu'ensuite, et il sera tel qu'il se sera fait » — l'identité humaine est un résultat, non un point de départ : l'homme est « tel qu'il se veut ».

\textbf{4.} Oui, ce texte implique une responsabilité totale : si l'homme est ce qu'il se fait, il ne peut excuser ses échecs par une « nature » ou un destin. Être libre, c'est répondre de soi devant soi et devant les autres ; Sartre dira ailleurs que nous sommes « condamnés à être libres ».

\textbf{5.} Avis personnel argumenté : on peut reconnaître que les contraintes sociales et économiques (pauvreté, chômage, poids des familles) sont réelles et limitent les possibles du jeune Camerounais. Mais la thèse de Sartre ne dit pas que tout est possible ; elle dit que, quelle que soit la situation, c'est encore moi qui lui donne un sens : subir ou entreprendre, fuir ou lutter. Le jeune qui crée sa coopérative de recyclage à Yaoundé ne choisit pas sa situation de départ, mais il choisit ce qu'il en fait. La liberté sartrienne n'est pas la toute-puissance : c'est la responsabilité dans la contrainte.
\end{corrige}

\begin{corrige}[Exercice 11 — Rédaction d'introduction et construction de problématique]
\textbf{Barème indicatif :} définitions et distinction : 4 pts ; contradiction et problématique : 4 pts ; introduction complète : 6 pts ; tableau des parties (argument, auteur, exemple) : 6 pts.

\textbf{1.} \emph{Vivre} : accomplir les fonctions biologiques de l'organisme — naître, se nourrir, durer. \emph{Exister} : au sens philosophique, être-là comme conscience libre, capable de choix, de projet et de rapport à soi. Ils ne sont pas synonymes : une plante vit sans exister ; et un homme peut vivre (biologiquement) sans véritablement exister, s'il n'a ni projet ni choix.

\textbf{2.} Contradiction : le langage courant traite « vivre » et « exister » comme équivalents (« il vit encore »), pourtant la philosophie les distingue — on peut vivre sans exister pleinement. Problématique : \emph{la vie biologique suffit-elle à faire une existence humaine, ou faut-il, pour exister vraiment, se donner un projet et un sens ?}

\textbf{3. Introduction rédigée (modèle).} Chaque jour, à l'aube, les rues de Douala s'animent : des millions d'hommes vivent — ils respirent, mangent, travaillent, dorment. Mais vivre suffit-il à faire une existence ? Le langage courant confond volontiers les deux termes : « vivre » désigne le fait biologique d'être en vie, tandis qu'« exister », au sens philosophique, désigne le mode d'être propre à l'homme, qui se définit par sa conscience, sa liberté et ses projets. Dès lors, un problème se pose : si tout homme vit, tout homme existe-t-il pour autant ? La vie biologique suffit-elle à faire une existence humaine, ou exiger que l'on se choisisse soi-même ? Nous verrons d'abord que vivre est la condition naturelle de toute existence (I), puis qu'exister suppose un dépassement de la simple vie par le projet et la liberté (II), enfin que l'existence pleine est celle qui assume sa finitude en se donnant un sens (III).

\textbf{4. Plan détaillé.}
\begin{center}
\begin{tabularx}{\linewidth}{|l|X|X|X|}
\hline
\rowcolor{popblueL} \textbf{Partie} & \textbf{Argument} & \textbf{Auteur} & \textbf{Exemple} \\
\hline
I & La vie biologique est la condition première : sans elle, rien n'est possible. & Épicure (le corps, le plaisir de vivre) & Le nourrisson, le malade soigné à l'hôpital : vivre est déjà un bien. \\
\hline
II & Mais exister, c'est se définir par ses choix, non seulement subsister. & Sartre (l'existence précède l'essence) & L'homme qui regarde passer les motos vs le jeune qui crée sa coopérative. \\
\hline
III & L'existence pleine assume sa finitude et se donne un sens. & Heidegger (être-pour-la-mort) & La mère de Bafoussam qui se bat pour scolariser ses enfants. \\
\hline
\end{tabularx}
\end{center}

\textbf{Points de vigilance :} ne pas annoncer le plan avant la problématique ; la problématique doit être une \emph{question} ; chaque partie du plan doit correspondre à une étape de la réponse, pas à une simple liste d'auteurs.
\end{corrige}

\newpage

\coursec{Fiche bilan}

\begin{popfigure}
\begin{tikzpicture}[mindmap, grow cyclic, every node/.style={concept, font=\small},
  root concept/.append style={concept color=popblue, font=\bfseries},
  level 1/.append style={level distance=3.4cm, sibling angle=60},
  level 2/.append style={level distance=2.2cm, font=\scriptsize}]
\node[root concept]{L'existence et la mort}
  child[concept color=poporange]{ node {Définir l'existence}
    child { node {Exister $\neq$ vivre} }
    child { node {Être-là, Dasein (Heidegger)} }
  }
  child[concept color=popgreen]{ node {Définir la mort}
    child { node {Mort biologique} }
    child { node {Mort de la personne} }
  }
  child[concept color=poppurple]{ node {Le sens de l'existence}
    child { node {Sens donné (Dieu)} }
    child { node {Sens à construire} }
  }
  child[concept color=poppink]{ node {L'absurde}
    child { node {Camus : le divorce} }
    child { node {Révolte et lucidité} }
  }
  child[concept color=popgold]{ node {L'existentialisme}
    child { node {Sartre : existence précède essence} }
    child { node {Liberté et responsabilité} }
  }
  child[concept color=popblue!60]{ node {Vie quotidienne}
    child { node {Deuils, rites funéraires} }
    child { node {Projets, choix de vie} }
  };
\end{tikzpicture}
\legende{Carte mentale de la leçon : les six branches à mobiliser pour toute dissertation sur l'existence et la mort.}
\end{popfigure}

\begin{bilan}
\textbf{Définitions clés à connaître par cœur}
\begin{itemize}
  \item \cle{Existence} : manière d'être propre à l'homme, être-là qui se sait vivant, se projette et donne un sens à sa vie ; à distinguer de la simple vie biologique.
  \item \cle{Mort} : cessation irréversible de la vie (mort biologique) ; pour la philosophie, fin de la personne et horizon qui donne son prix à l'existence.
  \item \cle{Absurde} (Camus) : confrontation entre l'appel humain au sens et le silence du monde.
  \item \cle{Angoisse} : sentiment par lequel l'homme découvre sa finitude et sa liberté (Heidegger, Sartre).
  \item \cle{Finitude} : caractère limité de l'existence humaine, vouée à la mort.
\end{itemize}

\textbf{Les grandes thèses et leurs auteurs}
\begin{center}
\begin{tabularx}{\linewidth}{|l|X|}
\hline
\rowcolor{popblueL}\textbf{Auteur} & \textbf{Thèse à retenir} \\
\hline
Heidegger & L'homme est un « être-pour-la-mort » : assumer sa finitude rend l'existence authentique. \\
\hline
Sartre & « L'existence précède l'essence » : l'homme n'est rien d'autre que ce qu'il se fait ; il est condamné à être libre. \\
\hline
Camus & L'absurde naît du divorce entre l'homme et le monde ; il faut imaginer Sisyphe heureux : révolte, liberté, passion. \\
\hline
Épicure & « La mort n'est rien pour nous » : quand nous sommes, elle n'est pas ; quand elle est, nous ne sommes plus. \\
\hline
Traditions africaines & La mort n'abolit pas le lien social : les ancêtres continuent d'exister dans la communauté (rites funéraires). \\
\hline
\end{tabularx}
\end{center}

\textbf{Méthodes express}
\begin{itemize}
  \item Dissertation : problématiser (l'existence a-t-elle un sens ? la mort la prive-t-elle de sens ?), opposer deux thèses, trancher avec nuance.
  \item Toujours distinguer \emph{vivre} (fait biologique) et \emph{exister} (se rapporter à sa vie).
  \item Illustrer chaque argument par un exemple concret (deuil, choix professionnel, engagement).
  \item Conclure en justifiant sa position personnelle argumentée.
\end{itemize}
\end{bilan}

\begin{aretenir}[Checklist avant le Bac]
\begin{itemize}
  \item[$\square$] Je sais définir l'existence et la distinguer de la simple vie biologique.
  \item[$\square$] Je sais définir la mort (biologique et philosophique) en une phrase précise.
  \item[$\square$] Je connais la thèse de Heidegger sur l'être-pour-la-mort.
  \item[$\square$] Je sais expliquer la formule de Sartre « l'existence précède l'essence ».
  \item[$\square$] Je sais définir l'absurde chez Camus et citer le mythe de Sisyphe.
  \item[$\square$] Je connais l'argument d'Épicure pour apprivoiser la peur de la mort.
  \item[$\square$] Je sais mobiliser la conception africaine des ancêtres et des rites funéraires.
  \item[$\square$] Je distingue les trois réponses au problème du sens : sens donné, sens à créer, absence de sens.
  \item[$\square$] Je sais construire une problématique à partir d'un sujet sur l'existence ou la mort.
  \item[$\square$] Je sais rédiger une introduction (accroche, définitions, problématique, plan).
  \item[$\square$] Je sais illustrer chaque thèse par un exemple de la vie courante camerounaise.
  \item[$\square$] Je sais rédiger une conclusion qui justifie ma position personnelle.
\end{itemize}
\end{aretenir}

\findecours

\end{document}
